\documentclass[11pt]{article}

\usepackage[T1]{fontenc}
\usepackage{amsmath,amssymb,amsthm}
\usepackage[margin=1in]{geometry}
\usepackage{xcolor}
\usepackage{hyperref}
\hypersetup{colorlinks=true,linkcolor=blue,urlcolor=blue,citecolor=blue}

\theoremstyle{plain}
\newtheorem{theorem}{Theorem}[section]
\newtheorem{lemma}[theorem]{Lemma}
\newtheorem{proposition}[theorem]{Proposition}
\newtheorem{corollary}[theorem]{Corollary}
\newtheorem{conjecture}[theorem]{Conjecture}
\theoremstyle{definition}
\newtheorem{definition}[theorem]{Definition}
\theoremstyle{remark}
\newtheorem{remark}[theorem]{Remark}

\newcommand{\Z}{\mathbb Z}
\newcommand{\Q}{\mathbb Q}
\newcommand{\F}{\mathbb F}
\newcommand{\C}{\mathbb C}
\newcommand{\OK}{\mathcal O_K}
\newcommand{\OL}{\mathcal O_L}
\newcommand{\Cl}{\operatorname{Cl}}
\newcommand{\Tr}{\operatorname{Tr}}
\newcommand{\Gal}{\operatorname{Gal}}
\newcommand{\ord}{\operatorname{ord}}
\newcommand{\Recon}{\operatorname{Recon}}
\newcommand{\Sq}{\operatorname{Sq}}
\newcommand{\PG}{\mathrm{PG}}
\newcommand{\paf}{\mathrm{PAF}}

\title{Toward a converse of Turyn's theorem:\\
obstructions to Turyn pairs at composite orders}
\author{Jeffery Kline}
\date{July 2026 --- consolidated working draft}

\begin{document}
\maketitle


\begin{abstract}
Turyn (1972) constructed Hadamard matrices of order $4t$ from pairs of
symmetric circulants $R,S$ of odd order $t$ with $R^2+S^2=qI$, $q=2t-1$,
whenever $q$ is a prime power, and conjectured that such pairs exist
\emph{only} then. This paper assembles a program of exact obstructions
toward that converse. The organizing result is a selection theorem:
nonexistence at $q$ is equivalent to emptiness of a finite marginal system
at some divisor of $t$, reducing the conjecture to per-family firing
criteria. Written criteria close every self-conjugacy-blind multi-prime
composite $q\le2000$, including apparently first determinations at
$q=1469$, $1937$, $1325$ (and, under GRH, $q=5185$ and $62305$); a
closed-table caveat applies and is stated in the text. On the lane
$q=6\ell-1$ with $\ell\equiv3\pmod4$ prime, the criteria fire
deterministically: a Turyn pair at composite $q$ forces the multiplier
image modulo $\ell$ down to $\{\pm1\}$. Consequences include an
unconditional semiprime nonexistence family and a conditional density
program whose remaining analytic inputs are stated as explicit,
falsifiable hypotheses. Every principal claim carries a verification-status
tier recorded in Appendix~B.
\end{abstract}

\tableofcontents

\section{The converse problem}\label{sec:intro}

\subsection{Turyn pairs and Turyn's theorem}

Let $t$ be odd, let $G=\langle x\rangle$ be cyclic of order $t$, and put
$q=2t-1$, so that $q\equiv1\pmod4$ (as $t$ is odd). Identify circulant matrices
of order $t$ with elements of the group ring $\Z[G]$.

\begin{definition}\label{def:turyn-pair}
A \emph{Turyn pair} of order $t$ is a pair $r=\sum_g r_g g,\ s=\sum_g s_g g$ in
$\Z[G]$ such that
\begin{itemize}
\item $r_1=0$ and $r_g\in\{\pm1\}$ for $g\ne1$;
\item $s_g\in\{\pm1\}$ for all $g$;
\item $r_{g^{-1}}=r_g$ and $s_{g^{-1}}=s_g$ (both symmetric); and
\item $r^2+s^2=q\cdot1$ in $\Z[G]$.
\end{itemize}
\end{definition}

From a Turyn pair one obtains the Williamson quadruple~\cite{Williamson44}
$A=R+I$, $B=R-I$, $C=D=S$, which satisfies $A^2+B^2+C^2+D^2=4tI$ and hence
yields a Hadamard matrix of order $4t$.

\begin{theorem}[Turyn 1972~\cite{Turyn72}]\label{thm:turyn}
If $q=2t-1$ is a prime power ($q\equiv1\pmod4$), a Turyn pair of order $t$
exists.
\end{theorem}

Turyn built the pair by ordering the quadratic-character matrix on
$\PG(1,q)$ along a Singer cycle (Section~\ref{sec:mary});
Whiteman~\cite{Whiteman73} re-proved it cyclotomically, and
Yamamoto--Yamada~\cite{YY85} identified the character values as relative Gauss
sums for $\F_{q^2}/\F_q$.

\subsection{The conjecture}

The closing remark of~\cite{Turyn72} states that the methods of~\cite{Turyn68}
show, ``in certain cases,'' that no analogue of the pair exists when $q$ is not
a prime power, and that ``it seems likely this is true in general.'' We take
this as the destination.

\begin{conjecture}[Turyn's converse]\label{conj:converse}
If a Turyn pair of order $t$ exists, then $q=2t-1$ is a prime power.
\end{conjecture}

Conjecture~\ref{conj:converse} is open. It has the flavour of, but is not
implied by, the computational classifications discussed below.

\subsection{The same object in other clothing}\label{sec:dictionary-intro}

The Turyn-pair equation is a hub; the following identifications are standard and
are used to import necessary conditions and to position the results
(Section~\ref{sec:related}).
\begin{itemize}
\item \emph{Negacyclic conference matrices.} By alternating-sign conjugation a
  Turyn pair of order $t$ is monomially equivalent to a negacyclic $W(2t,2t-1)$;
  see Seberry~\cite{Seberry16}, Lemmas 1.3--1.4. The
  Delsarte--Goethals--Seidel--Eades computation~\cite{DGS71,Seberry16} found
  that the only negacyclic $W(v,v-1)$ with $v<1000$ have $v=q+1$ with $q$ an
  odd prime power---the converse, verified by machine below $v=1000$.
\item \emph{Almost-perfect autocorrelation sequences.} The pair corresponds to
  a binary sequence of period $4t$ whose out-of-phase periodic autocorrelations
  all vanish except at the half-period shift $2t$, a value of $4-4t$; these are
  Wolfmann's almost-perfect sequences~\cite{Wolfmann92}, identified with cyclic
  divisible difference sets by Bradley--Pott~\cite{PB95}.
\item \emph{Two-circulant-core Hadamard matrices.} The pair is exactly the input
  to the two-circulant-core construction and its cocyclic analysis~\cite{BAOD22}.
\item \emph{Complex/Butson conference matrices.} With $\tau=r+is$, the matrix
  $P+iI$ (where $P$ is the conference matrix built from $R,S$) is a complex
  Hadamard matrix of order $q+1$; see Section~\ref{sec:mary}.
\end{itemize}


\subsection{Main results and reading guide}

The contribution is a theorem roadmap, not a proof of the conjecture.
Principal results, with verification tiers defined in Appendix~B:
\begin{itemize}
\item[(1)] \textbf{Selection theorem} (Theorem~\ref{thm:divisor-selection},
  Tier~0): nonexistence at $q$ is equivalent to emptiness of some divisor
  node's marginal system; feasibility is monotone on the divisor lattice
  and the orbit model at every node is computed by the multiplier group
  $M(q)$ alone. This converts the converse into per-family firing theorems.
\item[(2)] \textbf{Closed-form firing criteria} (Tier~0): the transitive
  two-orbit representation-pair criterion, the $h=3$
  support/augmentation hierarchy, the full-torus row-sum bound --- sharp
  exactly at Turyn's $q=25$ --- with its exact sub-threshold form, the
  full-unit prime-power nested criterion, and lattice-reduced period
  joins.
\item[(3)] \textbf{Finite closure} (Tier~0 per member): every
  self-conjugacy-blind multi-prime composite $q\le2000$ is closed by
  written quantified criteria, including apparently first determinations
  $q=1469$, $1937$, $1325$; under GRH (Tier~1) also $q=5185$ and
  $q=62305$.
\item[(4)] \textbf{The lane $q=6\ell-1$, $\ell\equiv3\pmod4$ prime}: a
  Turyn pair at composite $q$ forces the multiplier image in
  $(\Z/\ell)^\times$ down to $\{\pm1\}$
  (Proposition~\ref{prop:lane-kill}, Tier~0); the semiprime case gives an
  unconditional nonexistence family
  (Theorem~\ref{thm:semiprime-converse}, Tier~0) whose infinitude is
  parity-blocked; a density program under the explicit hypothesis
  $\mathrm{EQ}_3$ (Tier~2) gives the small-factor obstruction count
  (Proposition~\ref{prop:small-factor-sieve}), with the passage to an
  infinite family isolated in one open structure count
  (Remark~\ref{rem:lane-density-status}) --- open also under GRH.
\item[(5)] \textbf{Construction side} (Tier~0): a uniform
  character-lifting normal form and fillability dichotomy, in
  Appendix~A.
\end{itemize}
Sections 2--3 build the character and multiplier machinery; Section~4 is
the obstruction hierarchy and contains results (1)--(4); Sections 6--8
record the closed-order ledger, related work, and what remains.

\section{The character identity}\label{sec:character}

All obstructions flow from one identity. Put
\[
  \tau=r+is\in\Z[i][G],\qquad \tau^*=r-is.
\]
Because $r,s$ are symmetric, coefficient conjugation and the group-ring adjoint
agree on $\tau$, and the pair equation is
\begin{equation}\label{eq:tau}
  \tau\tau^*=q.
\end{equation}
For a character $\chi$ of $G$ of order $d$, the values $\chi(r),\chi(s)$ are
real, so
\begin{equation}\label{eq:charnorm}
  \chi(\tau)\,\overline{\chi(\tau)}=\chi(r)^2+\chi(s)^2=q,
\end{equation}
an equation in $\Z[\zeta_d,i]=\Z[\zeta_{4d}]$. Thus each character value
$\chi(\tau)$ is a cyclotomic integer of absolute value $\sqrt q$ at every
archimedean place: the Turyn-pair equation is a family of \emph{norm-$q$
conditions}, one per character, coupled by integrality of the coefficients.
Every result below analyzes these conditions---at the trivial character
(\S\ref{sec:twosquares}), over inert towers (\S\ref{sec:selfconj}), under common
decomposition subgroups of the full Galois action (\S\ref{sec:multiplier}), by
orbit-reduced full-group and marginal decisions
(\S\S\ref{sec:t3-decision}--\ref{sec:t2-marginal}), or after descent to a prime
quotient (\S\S\ref{sec:rigidity}--\ref{sec:defect}).

\section{Multiplier theorems}\label{sec:multiplier}

\subsection{The prime-power case: exact signs}

Before the composite obstructions we record the rigidity that the prime-power
case already carries; it both frames the method and pins the symmetry the
converse must contend with. For $u\in(\Z/t)^\times$ let $r^{(u)}$ denote
decimation $x\mapsto x^u$. Say $u$ is a \emph{multiplier} of $(r,s)$ if
$r^{(u)}=\epsilon r$ and $s^{(u)}=\delta s$ for signs $\epsilon,\delta\in\{\pm1\}$.
Symmetry makes $-1$ a trivial multiplier with signs $(1,1)$.

\begin{theorem}[Frobenius multiplier, exact signs]\label{thm:multiplier}
Let $(r,s)$ be a Turyn pair of order $t$ with $q=2t-1=p^e$, $p$ prime. Then $p$
is a multiplier, and
\[
  p\equiv1\!\!\pmod4\ \Rightarrow\ r^{(p)}=r,\ s^{(p)}=s;
  \qquad
  p\equiv3\!\!\pmod4\ \Rightarrow\ r^{(p)}=-r,\ s^{(p)}=s.
\]
In the second case $e$ is even, $\sum_g r_g=0$, and $\sum_g s_g=\pm p^{e/2}$.
\end{theorem}

When $e=1$ this only restates the symmetry $p\equiv-1\pmod t$; the first
nontrivial content is at $e\ge2$. Let $\sigma_p\in\Gal(\Q(\zeta_{4t})/\Q)$ send
$\zeta_{4t}\mapsto\zeta_{4t}^p$ (well-defined since $p\mid q$ is coprime to
$4t$), and define the Frobenius twist
$\tau_p=\sum_g\sigma_p(\tau_g)g^p=r^{(p)}+i^p s^{(p)}$, so that
$\chi(\tau_p)=\sigma_p(\chi(\tau))$ for every character $\chi$.

\begin{lemma}\label{lem:per-character-unit}
For every character $\chi$ of $G$ of order $d$, the quotient
$u_\chi=\sigma_p(\chi(\tau))/\chi(\tau)$ is a root of unity in $\Q(\zeta_{4d})$.
\end{lemma}

\begin{proof}
Let $\alpha=\chi(\tau)\in\Z[\zeta_{4d}]$. Every Galois conjugate of $\alpha$ has
absolute value $\sqrt q$: a conjugate sends $\chi$ to some power $\chi^m$ and
$i$ to $\pm i$, so its squared absolute value is $\chi^m(r)^2+\chi^m(s)^2=q$
by~\eqref{eq:charnorm}; in particular $\alpha\ne0$. The identity
$\alpha\overline\alpha=q=p^e$ shows $(\alpha)$ is a product of primes above $p$.
As $p\nmid4d$, $p$ is unramified in $\Q(\zeta_{4d})$ and $\sigma_p$ is its
Frobenius, so $\sigma_p$ fixes each prime above $p$; hence
$\sigma_p((\alpha))=(\alpha)$ and $u_\chi$ is a unit. All its conjugates have
absolute value $1$, so Kronecker's theorem makes $u_\chi$ a root of unity; as
$d$ is odd, these are $\mu_{4d}$.
\end{proof}

\begin{lemma}\label{lem:coherence}
There exist $a\in\Z/4$ and $c\in\Z/t$ with $u_\chi=i^a\chi(x^c)$ for every
character $\chi$.
\end{lemma}

\begin{proof}
For $m\in(\Z/t)^\times$ let $\gamma_m$ fix $i$ and send $\zeta_t\mapsto\zeta_t^m$;
since the cyclotomic Galois group is abelian and $\tau$ has $\Z[i]$
coefficients, $\gamma_m(u_\chi)=u_{\chi^m}$. The characters of exact order $d$
form one orbit under these $\gamma_m$ (as $(\Z/t)^\times\to(\Z/d)^\times$ is
onto), so for each $d\mid t$ there are $a_d\in\Z/4$ and $c_d\in\Z/d$ with
$u_\psi=i^{a_d}\psi(x^{c_d})$ for every $\psi$ of exact order $d$. We show these
data cohere. Set $\beta=\tau_p\tau^*\in\Z[i][G]$, so
$\chi(\beta)=\sigma_p(\chi(\tau))\overline{\chi(\tau)}=q\,u_\chi$. Mapping
$\Z[i][G]\to\Z[i][C_e]$ and inverting the Fourier transform on $C_e$ gives, for
the coefficient of $y^j$,
\[
  \pi_e(\beta)_j=\frac qe\sum_{d\mid e}i^{a_d}R_d(c_d-j),\qquad
  R_d(n)=\sum_{m\in(\Z/d)^\times}\zeta_d^{mn}
\]
the Ramanujan sum. Since $\pi_e(\beta)_j\in\Z[i]$ and $(q,e)=1$,
\begin{equation}\label{eq:integrality}
  e\ \Big|\ \sum_{d\mid e}i^{a_d}R_d(c_d-j)\qquad\text{in }\Z[i]
\end{equation}
for all $j$ and all $e\mid t$. Induct on the number of prime factors of $e$
(with multiplicity); $e=1$ is empty. If $e=P$ is prime, choose $j\not\equiv
c_P\pmod P$; then $R_1\equiv1$, $R_P(c_P-j)=-1$, so~\eqref{eq:integrality}
gives $P\mid i^{a_1}-i^{a_P}$, impossible unless $a_P=a_1$ (a nonzero multiple
of $P$ has norm $\ge P^2>4$, while a difference of fourth roots of unity has
norm $\le4$). For general $e$, by induction $a_d=a_1=:a$ for proper $d\mid e$
and the $c_d$ are compatible; choose $c_*$ with $c_*\equiv c_d\pmod d$ for all
proper $d\mid e$. Using $\sum_{d\mid e}R_d(n)=e$ or $0$ according as $e\mid n$
or not, \eqref{eq:integrality} reduces to
\begin{equation}\label{eq:top-step}
  e\mid i^{a_e}R_e(c_e-j)-i^aR_e(c_*-j)\qquad\text{for all }j.
\end{equation}
The formula $R_e(n)=\mu(e/(e,n))\varphi(e)/\varphi(e/(e,n))$ shows, for odd $e$:
the value $\varphi(e)$ occurs only when $e\mid n$, the value $-\varphi(e)$ never
occurs, and all values exceed $-\varphi(e)/(P-1)$ ($P$ the least prime factor).
If $a_e-a$ is odd the two terms in~\eqref{eq:top-step} lie in $\Z$-independent
directions, forcing $e\mid R_e(c_*-j)$ for all $j$---impossible at $j=c_*$ where
the value is $\varphi(e)\in(0,e)$. If $a_e-a\equiv2\pmod4$ then $e\mid
R_e(c_e-j)+R_e(c_*-j)$; the bounds confine the sum to $(-e,2e)$, so it is $0$ or
$e$, and summing over a period gives $0$, whence it vanishes identically---but
at $j=c_*$ this reads $R_e(c_e-c_*)=-\varphi(e)$, impossible for odd $e$. Thus
$a_e=a$, and~\eqref{eq:top-step} gives $e\mid R_e(c_e-j)-R_e(c_*-j)$. If $e$ has
two distinct prime factors $P<Q$ the difference is bounded by
$\varphi(e)P/(P-1)<e$, so the two Ramanujan functions agree and $c_*\equiv
c_e\pmod e$; if $e=P^k$ the values are $\varphi(e),-P^{k-1},0$ and one checks
$c_*\equiv c_e\pmod{P^{k-1}}$, exactly the compatibility needed. The induction
closes.
\end{proof}

\begin{proof}[Proof of Theorem~\ref{thm:multiplier}]
By Lemma~\ref{lem:coherence}, $\chi(\beta)=q\,i^a\chi(x^c)$ for all $\chi$, so
$\beta=\tau_p\tau^*=q\,i^a x^c$ in $\Z[i][G]$. Multiplying by $\tau$ and
using~\eqref{eq:tau} gives $\tau_p=i^a x^c\tau$. The identity coefficient of
$\tau$ is $is_1$ (modulus $1$), while every other has modulus $\sqrt2$; the same
holds for $\tau_p$ and decimation fixes the identity, so $c=0$. Comparing
identity coefficients gives $i^p s_1=i^{a+1}s_1$, so $a\equiv p-1\pmod4$ and
$\tau_p=i^{p-1}\tau$. For $p\equiv1$ this is $r^{(p)}+is^{(p)}=r+is$; for
$p\equiv3$ it is $r^{(p)}-is^{(p)}=-(r+is)$, i.e.\ $r^{(p)}=-r$, $s^{(p)}=s$.
Finally, if $p\equiv3\pmod4$, at the trivial character
$(\sum r_g+i\sum s_g)(\sum r_g-i\sum s_g)=p^e$; $p$ is inert in $\Z[i]$, so $e$
is even and the left factor is $i^\kappa p^{e/2}$; as $\sum s_g$ is odd its
imaginary part is nonzero, forcing $\sum r_g=0$ and $\sum s_g=\pm p^{e/2}$.
\end{proof}

\begin{corollary}\label{cor:orbits}
If $q=p^e$ then every Turyn pair is determined by its values on the
$\langle p\rangle$-orbits of $G$: constant along them for $p\equiv1\bmod4$,
sign-alternating ($r_{g^p}=-r_g$) for $p\equiv3\bmod4$.
\end{corollary}

Empirically the multiplier group is \emph{exactly} $\langle p\rangle$ on all
$328$ pairs of the exhaustive sweep over odd $t\le41$---a strengthening, since
the theorem asserts only that $p$ is a multiplier---consistent with the
uniqueness up to equivalence observed by Lang--Schneider~\cite{LS12} through
order $99$. Theorem~\ref{thm:multiplier} is best viewed as a normal-form
refinement of classical multiplier theory~\cite{Hall67}: no group-development hypothesis is
made (any solution of~\eqref{eq:tau} already carries the Frobenius symmetry of
the finite-field examples), the zero diagonal removes the usual translate
ambiguity, and the signs are pinned exactly. For the group-developed objects
behind the finite-field examples (Bose's affine relative difference set and its
balanced-generalized-weighing-matrix avatars) see
Jungnickel--Kharaghani~\cite{JK}.

\subsection{Composite common-decomposition multipliers}
\label{sec:composite-multiplier}

At composite $q$ no single Frobenius $\sigma_p$ fixes the prime ideals above all
other rational primes dividing $q$. The part of the Galois group that survives
is the intersection of the relevant decomposition groups.

\begin{definition}\label{def:composite-M}
Let $q=2t-1=\prod_{p\mid q}p^{e_p}$ be composite. For each $p\mid q$ write
$\langle p\rangle$ for the subgroup generated by $p$ in $(\Z/4t)^\times$, and
put
\[
  M(q)=\bigcap_{p\mid q}\langle p\rangle\subseteq(\Z/4t)^\times .
\]
For $\sigma\in M(q)$ write $\sigma=(\epsilon,u)$, where
$\epsilon=1$ or $-1$ according as $\sigma$ is $1$ or $3$ modulo $4$, and
$u=\sigma\bmod t\in(\Z/t)^\times$.
\end{definition}

\begin{lemma}[Composite coherence]\label{lem:composite-coherence}
Let $\sigma=(\epsilon,u)\in M(q)$ and put
\[
  \tau_\sigma=\sum_g\sigma(\tau_g)g^u .
\]
Then there are $a\in\Z/4$ and $c\in\Z/t$ such that
\[
  \tau_\sigma\tau^*=q\,i^a x^c .
\]
\end{lemma}

\begin{proof}
For a character $\chi$ of order $d$, put
$\alpha=\chi(\tau)\in\Z[\zeta_{4d}]$.  Since
$\alpha\overline\alpha=q$, the ideal $(\alpha)$ is supported only at primes
above rational primes dividing $q$.  The image of $\sigma$ in
$(\Z/4d)^\times$ lies in $\langle p\rangle$ for every $p\mid q$, hence in the
decomposition group of every prime above $p$ in the abelian field
$\Q(\zeta_{4d})$.  Thus $\sigma((\alpha))=(\alpha)$ and
$u_\chi=\sigma(\alpha)/\alpha$ is a unit.  All conjugates of $u_\chi$ have
absolute value $1$, so Kronecker's theorem makes $u_\chi$ a root of unity.

The roots $u_\chi$ satisfy the same two coherence inputs used in
Lemma~\ref{lem:coherence}.  First, decimation commutes with $\sigma$, so for
every $m\in(\Z/t)^\times$,
\[
  \gamma_m(u_\chi)=u_{\chi^m}.
\]
Second, with $\beta=\tau_\sigma\tau^*$, the coefficient of $y^j$ in the image
of $\beta$ in $\Z[i][C_e]$ is
\[
  \pi_e(\beta)_j=\frac qe\sum_{d\mid e} i^{a_d}R_d(c_d-j),
\]
once the exact-order-$d$ character orbit is written as
$u_\psi=i^{a_d}\psi(x^{c_d})$.  Since $\pi_e(\beta)_j\in\Z[i]$ and $(q,e)=1$,
the integrality congruence~\eqref{eq:integrality} holds for every $e\mid t$
and every $j$.  The Ramanujan-sum induction in the proof of
Lemma~\ref{lem:coherence} uses only this congruence, the compatibility of the
character orbits, and the coefficient ring $\Z[i]$; it does not use that the
automorphism came from one prime Frobenius.  Hence the $a_d$ and $c_d$ cohere:
there are $a\in\Z/4$ and $c\in\Z/t$ with
$u_\chi=i^a\chi(x^c)$ for every $\chi$.  Fourier inversion then gives
$\beta=q\,i^a x^c$, as claimed.
\end{proof}

\begin{theorem}[Composite multiplier]\label{thm:composite-multiplier}
Let $(r,s)$ be a Turyn pair of order $t$ and let $q=2t-1$ be composite. For every
$\sigma=(\epsilon,u)\in M(q)$,
\[
  r^{(u)}=\epsilon r,\qquad s^{(u)}=s .
\]
\end{theorem}

\begin{proof}
Let $\tau=r+is$ and
$\tau_\sigma=\sum_g\sigma(\tau_g)g^u=r^{(u)}+\epsilon i\,s^{(u)}$.
By Lemma~\ref{lem:composite-coherence}, $\tau_\sigma\tau^*=q\,i^a x^c$ for some
$a\in\Z/4$ and $c\in\Z/t$.
Multiplying by $\tau$ and using $\tau\tau^*=q$ gives
$\tau_\sigma=i^a x^c\tau$. The identity coefficient of $\tau_\sigma$ has
modulus $1$, while a nonidentity coefficient of $\tau$ has modulus $\sqrt2$;
hence $c=0$. Comparing identity coefficients gives $i^a=\epsilon$, and so
$\tau_\sigma=\epsilon\tau$. Unpacking the real and imaginary parts gives the two
displayed identities.
\end{proof}

\begin{remark}
For $q=p^e$ the group $M(q)$ is just $\langle p\rangle$, so
Theorem~\ref{thm:composite-multiplier} recovers the multiplier part of
Theorem~\ref{thm:multiplier}; the prime-power theorem additionally records the
trivial-character consequences in the sign-alternating case. For composite $q$
the theorem should be read with the same caution as the prime-power multiplier:
it is an unconditional normal-form constraint, while novelty relative to the
divisible-difference-set multiplier literature is not audited here.
\end{remark}

\section{A hierarchy of obstructions at composite \texorpdfstring{$q$}{q}}
\label{sec:hierarchy}

\subsection{The trivial character}\label{sec:twosquares}

At $\chi=1$, \eqref{eq:charnorm} reads $\bigl(\sum_g r_g\bigr)^2+\bigl(\sum_g
s_g\bigr)^2=q$.

\begin{proposition}\label{prop:twosquares}
If some prime $\ell\equiv3\pmod4$ divides $q$ to an odd power, no Turyn pair of
order $t$ exists.
\end{proposition}

This elementary sum-of-two-squares test closes $211$ of the composite $q\le2000$.

\begin{lemma}[Signed branch after two squares]\label{lem:signed-branch-reduction}
Suppose the subgroup generated by $M(q)$ and the trivial symmetry contains an
element with $\epsilon=-1$.  Then every prime divisor of $q$ is congruent to
$3\bmod4$.  Consequently, after Proposition~\ref{prop:twosquares}, the only
possible sign-alternating composite branch has
\[
  q=n^2,\qquad p\mid n\Rightarrow p\equiv3\pmod4 .
\]
In particular, every composite $q$ with some prime divisor $p\equiv1\pmod4$
is automatically in the all-plus multiplier regime.
\end{lemma}

\begin{proof}
The trivial symmetry has $\epsilon=+1$, so an $\epsilon=-1$ element in the
generated group already comes from $M(q)$.  Let
$\sigma=(\epsilon,u)\in M(q)$ have $\epsilon=-1$.  For every prime $p\mid q$,
the definition of $M(q)$ gives $\sigma\in\langle p\rangle\subset(\Z/4t)^\times$,
so $\sigma\equiv p^a\pmod{4t}$ for some $a$.  Reducing modulo $4$ gives
$p^a\equiv3\pmod4$, hence $p\equiv3\pmod4$ and $a$ is odd.  Thus all prime
divisors of $q$ are $3\bmod4$.  If a Turyn pair exists, the trivial-character
equation makes $q$ a sum of two squares, so Proposition~\ref{prop:twosquares}
forces every exponent of every such prime in $q$ to be even.  This is exactly
the displayed square branch.
\end{proof}

\subsection{The self-conjugacy boundary}\label{sec:selfconj}

The classical refinement is the self-conjugacy (inert-tower) argument of
Turyn~\cite{Turyn65,Turyn68}, in the present normal form. It is the class of
argument Turyn indicated for composite $q$ in~\cite{Turyn72}.

\begin{proposition}\label{prop:selfconj}
Let $q=2t-1$ and suppose some prime $\ell\equiv3\pmod4$ has $\ell^k\,\|\,q$ with
either \emph{(i)} $k$ odd, or \emph{(ii)} $k$ even and $\ord_d(\ell)$ odd for
every divisor $d\mid t$. Then no Turyn pair of order $t$ exists.
\end{proposition}

\begin{proof}
If $k$ is odd, the trivial character gives a sum of two squares equal to $q$,
impossible (Proposition~\ref{prop:twosquares}). Assume (ii). Let $\chi$ have
exact order $d\mid t$ and put $\alpha=\chi(\tau)\in K_d=\Q(\zeta_{4d})$. Since
$\ell\mid q=2t-1$ is odd and coprime to $t$, $\ell\nmid4d$ is unramified in
$K_d$. Let $\lambda\mid\ell$ in $\Q(\zeta_d)$; its residue degree is
$\ord_d(\ell)$, odd by hypothesis, so the residue field has order
$\ell^{\ord_d(\ell)}\equiv3\pmod4$ and lacks a primitive fourth root of unity;
hence $\lambda$ is inert in $K_d=\Q(\zeta_d)(i)$ and every prime
$\Lambda\mid\ell$ of $K_d$ satisfies $\Lambda=\overline\Lambda$. From
$\alpha\overline\alpha=q$, $2\nu_\Lambda(\alpha)=\nu_\Lambda(q)=k$, so
$\ell^{k/2}\mid\chi(\tau)$ for every $\chi$. Fourier inversion
$t\,\tau_g=\sum_\chi\chi(\tau)\chi(g^{-1})$ shows $\ell^{k/2}$ divides
$t\,\tau_g$, hence (as $\Z[i]$ is integrally closed and $(\ell,t)=1$) every
$\tau_g$; but $|\tau_g|^2\in\{1,2\}$ while $\ell^k\ge9$, a contradiction.
\end{proof}

This closes $q=45$ ($t=23$, $\ord_{23}(3)=11$ odd), the first composite case
beyond two squares, and $12$ further $q\le2000$.

\begin{remark}[The blind spot]\label{rem:blindspot}
Both Propositions~\ref{prop:twosquares} and~\ref{prop:selfconj} require a prime
divisor $\ell\equiv3\pmod4$ in a suitable position---an odd power for the
two-squares obstruction, an order-parity condition for self-conjugacy. A
composite $q$ survives both exactly when every $\ell\equiv3\pmod4$ dividing it
occurs to an even power and the self-conjugacy test does not fire. The simplest
such orders have all prime divisors $\equiv1\pmod4$---e.g.\ $q=65=5\cdot13$
($t=33$), for which no inert prime appears in any $\Q(\zeta_{4d})$ layer---but
the blind spot is larger: of the $111$ surviving composite $q\le2000$, $80$ are
all-$\equiv1\pmod4$ and $31$ carry a $\equiv3\pmod4$ prime to an even power
(e.g.\ $q=245=5\cdot7^2$). These $111$ orders are the domain of
Sections~\ref{sec:t3-decision}--\ref{sec:defect}. (Exact quotient descent
nonetheless closes $q=65$ via $\Z[C_{33}]\to\Z[C_{11}]$; the T3 and quotient
mechanisms developed next attack this same blind spot by different means.)
\end{remark}

\subsection{The composite-multiplier reduced decision}\label{sec:t3-decision}

Theorem~\ref{thm:composite-multiplier} also gives a full-group, field-free
decision method when the multiplier signs are all positive. Let
\[
  V=\langle\,u:(1,u)\in M(q),\ -1\,\rangle\subseteq(\Z/t)^\times .
\]
If every element of the multiplier group generated by $M(q)$ and the trivial
symmetry $(-1)$ has $\epsilon=+1$, then any Turyn pair is constant on the
$V$-orbits of $C_t$. The pair equation is equivalently
\[
  \paf_r(\delta)+\paf_s(\delta)=0\qquad(\delta\ne0),
\]
and the periodic autocorrelation functions are themselves $V$-invariant. Thus
there is one equation for each nonzero $V$-orbit of shifts, and the $r$-side and
$s$-side can be enumerated independently.

\begin{proposition}[Orbit-reduced decision]\label{prop:t3-decision}
Assume the multiplier signs above are all $+1$. Enumerate the possible signs of
$r$ and $s$ on $V$-orbits, with $r_1=0$, store the vectors
$(\paf_r(\delta))$ on nonzero shift-orbit representatives, and meet in the
middle against $-(\paf_s(\delta))$. If this finite search has no match, no Turyn
pair of order $t$ exists. If it has a match, the resulting rows are a Turyn pair
exactly when the full integer check $r^2+s^2=q$ passes.
\end{proposition}

\begin{proof}
Theorem~\ref{thm:composite-multiplier} and the symmetry $j\mapsto-j$ force every
putative pair into the enumerated $V$-invariant search space. The equation
$r^2+s^2=q$ has zero off-identity coefficients precisely when the displayed PAF
equalities hold for all nonzero shifts; since PAF is $V$-invariant, the listed
representatives are enough. The final statement is the same identity checked in
the full group ring, with exact integer arithmetic.
\end{proof}

The direct sign-reversing kill mechanism does not fire in the current range:
the sweep through $q\le2000$ finds zero multi-prime composite cases surviving
the classical filters for which an $\epsilon=-1$ multiplier is present.  The
reduced all-plus decision is nonetheless strong. Over the $111$
self-conjugacy-blind composite $q\le2000$, it closes $85$ cases. Of these,
$52$ overlap the rigidity scan and $33$ are new relative to quotient rigidity;
the new list is recorded in Section~\ref{sec:results}. The closures are
field-free and unconditional.  The implementation keeps a hard enumeration cap
and refuses sign-alternating inputs rather than applying an unsound
constant-orbit reduction; Proposition~\ref{prop:signed-marginal} is the
theorem-level replacement for that guard when an $\epsilon=-1$ branch must be
handled.

\subsection{Sub-torus marginals and cross-prime gluing}\label{sec:t2-marginal}

The full T3 decision can be too large when $M(q)$ is small, but a putative pair
has smaller marginals at every divisor of $t$. Let $t'\mid t$, put $h=t/t'$, and
let $\pi:\Z[C_t]\to\Z[C_{t'}]$ be induced by $x\mapsto y$. Set
$A=\pi(r)$ and $B=\pi(s)$. Then
\[
  AA^*+BB^*=q\cdot1\qquad\text{in }\Z[C_{t'}],
\]
where $*$ is inversion in the group ring. The coefficients have forced boxes:
the origin coefficient of $A$ is an even sum of $h-1$ signs and has absolute
value at most $h-1$, while every other coefficient of $A$ and every coefficient
of $B$ is an odd sum of $h$ signs and has absolute value at most $h$.

Let $V$ be the multiplier group used in Section~\ref{sec:t3-decision}, and let
$V'$ be its image in $(\Z/t')^\times$. By
Theorem~\ref{thm:composite-multiplier}, $A$ and $B$ are constant on the
$V'$-orbits of $C_{t'}$.

\begin{definition}[Marginal orbit algebra]\label{def:marginal-orbit-algebra}
Let $\Omega=\Omega(t',V')$ be the set of $V'$-orbits on $\Z/t'$. For
$O\in\Omega$ put
\[
  E_O=\sum_{j\in O} y^j,\qquad
  \mathcal A(t',V')=\Z[C_{t'}]^{V'}=\bigoplus_{O\in\Omega}\Z E_O .
\]
For $O,P,Q\in\Omega$, define $N_{O,P}^Q$ to be the number of pairs
$(a,b)\in O\times P$ with $a-b=k$ for any fixed $k\in Q$. This number is
independent of the choice of $k\in Q$, and
\[
  E_OE_P^*=\sum_{Q\in\Omega}N_{O,P}^Q E_Q .
\]
\end{definition}

The marginal equation is therefore a finite quadratic system in the orbit
values. Writing
\[
  A=\sum_{O\in\Omega}a_OE_O,\qquad
  B=\sum_{O\in\Omega}b_OE_O,
\]
we get
\[
  \sum_{O,P}N_{O,P}^Q(a_Oa_P+b_Ob_P)
  =\begin{cases} q,& Q=\{0\},\\ 0,& Q\ne\{0\},\end{cases}
\]
for each orbit $Q\in\Omega$. Equivalently,
\[
  \paf_A(\delta)+\paf_B(\delta)=0
  \quad(\delta\ne0),
  \qquad
  \sum_j A_j^2+\sum_j B_j^2=q,
\]
with one PAF equation for each nonzero $V'$-orbit of shifts. The last equation
is automatic at $h=1$, but is a real counting constraint for proper marginals.

\begin{lemma}[Marginal algebra lemma]\label{lem:marginal-algebra}
Assume the composite multiplier signs are all $+1$. Every Turyn pair of order
$t$ induces, for every divisor $t'\mid t$, orbit-value vectors
$a=(a_O)_{O\in\Omega}$ and $b=(b_O)_{O\in\Omega}$ in
$\mathcal A(t',V')$ satisfying the coefficient equations above, the trivial
character equation
\[
  \left(\sum_O |O|a_O\right)^2+\left(\sum_O |O|b_O\right)^2=q,
\]
and the following boxes and parities: if $O=\{0\}$, then $a_O$ is even and
has $|a_O|\le h-1$; if $O\ne\{0\}$, then $a_O$ is odd and $|a_O|\le h$; and
every $b_O$ is odd with $|b_O|\le h$.
\end{lemma}

\begin{proof}
The map $\pi:\Z[C_t]\to\Z[C_{t'}]$ sends a Turyn pair to a marginal pair
$(A,B)$ satisfying $AA^*+BB^*=q\cdot1$. Theorem~\ref{thm:composite-multiplier}
and the definition of $V'$ put $A,B$ in $\mathcal A(t',V')$, so expanding in the
orbit-sum basis and using Definition~\ref{def:marginal-orbit-algebra} gives the
coefficient equations. Evaluating at the trivial character gives the displayed
two-square equation. The boxes and parities are exactly the fiber-sum boxes:
the origin coefficient of $A$ omits the zero diagonal and sums $h-1$ signs,
while all other coefficients sum $h$ signs.
\end{proof}

\begin{definition}[Signed marginal module]\label{def:signed-marginal-module}
Let $\widetilde M(q)$ be the subgroup generated by $M(q)$ and the trivial
symmetry $(+1,-1)$, and let
\[
  E'=\{(\epsilon,u\bmod t'):(\epsilon,u)\in\widetilde M(q)\}
  \subseteq\{\pm1\}\times(\Z/t')^\times .
\]
Let $H'$ be its decimation image.  Define the signed marginal module
\[
  \mathcal S(t',E')=
  \{X=\sum_j X_jy^j\in\Z[C_{t'}]: X_{uj}=\epsilon X_j
    \text{ for all }(\epsilon,u)\in E'\}.
\]
Thus a negative stabilizer of $j$ forces $X_j=0$.  When all signs in $E'$ are
$+1$, this module is the invariant algebra $\mathcal A(t',V')$ above.
\end{definition}

\begin{proposition}[Signed marginal obstruction]\label{prop:signed-marginal}
Let $q=2t-1$ be composite and let $t'\mid t$.  Every Turyn pair induces
marginals
\[
  A=\pi(r)\in\mathcal S(t',E'),\qquad
  B=\pi(s)\in\mathcal A(t',H')
\]
satisfying
\[
  AA^*+BB^*=q\cdot1,\qquad A(1)^2+B(1)^2=q,
\]
together with the same fiber boxes and parities as in
Lemma~\ref{lem:marginal-algebra}.  Consequently, if this finite signed
marginal system is infeasible for some $t'\mid t$, then no Turyn pair of order
$t$ exists.
\end{proposition}

\begin{proof}
For $(\epsilon,u)\in\widetilde M(q)$, Theorem~\ref{thm:composite-multiplier}
and the symmetry $j\mapsto-j$ give $r_{uj}=\epsilon r_j$ and $s_{uj}=s_j$.
Fiber summation commutes with decimation because $u$ is prime to $t$ and hence
permutes the fibers over $C_{t'}$; therefore the marginals satisfy
$A_{ua}=\epsilon A_a$ and $B_{ua}=B_a$.  The group-ring equation, the
trivial-character equation, and the boxes/parities are exactly the images of
the Turyn pair equation and the fiber sums, as in
Lemma~\ref{lem:marginal-algebra}.  Thus any Turyn pair gives a point of the
finite signed system, and emptiness of that system is an obstruction.
\end{proof}

\begin{corollary}[Signed full-torus sign criterion]
\label{cor:signed-full-torus-sign}
Assume the hypotheses of Proposition~\ref{prop:signed-marginal} and take the
full torus $t'=t$.  Then $h=1$, so the boxes force
\[
  A_0=0,\quad A_j\in\{\pm1\}\ (j\ne0),\qquad
  B_j\in\{\pm1\}\quad(0\le j<t),
\]
with $A$ in the signed module $\mathcal S(t,E)$ and $B$ invariant under the
decimation image.  If the resulting finite signed sign-join
\[
  AA^*+BB^*=q\cdot1,\qquad A(1)^2+B(1)^2=q
\]
is empty, then no Turyn pair of order $t$ exists.
\end{corollary}

\begin{proof}
This is Proposition~\ref{prop:signed-marginal} with $h=1$.  The origin
coefficient of $A$ is forced to be $0$, all other coefficients are signs, and
the signed/invariant orbit relations are exactly those defining
$\mathcal S(t,E)$ and $\mathcal A(t,H)$.
\end{proof}

\begin{corollary}[Signed full-unit prime-torus obstruction]
\label{cor:signed-full-unit-prime}
Let $t$ be prime and $t>5$.  Suppose the image of $\widetilde M(q)$ in
$(\Z/t)^\times$ is all of $(\Z/t)^\times$ and that some element of
$\widetilde M(q)$ has sign $\epsilon=-1$.  Then no Turyn pair of order $t$
exists.
\end{corollary}

\begin{proof}
If two elements of $\widetilde M(q)$ have the same decimation and opposite
signs, then the multiplier equations force $r=-r$, impossible because
$r_j=\pm1$ for $j\ne0$.  Thus, for any putative pair, the sign descends to a
nontrivial character of $(\Z/t)^\times$.  Since $(\Z/t)^\times$ is cyclic, this
character is the quadratic character $\eta$ on $\F_t^\times$; it is even
because the trivial symmetry $j\mapsto-j$ has sign $+1$.  The multiplier
equations force $s_j=s_*$ for all $j\ne0$, and $r_j=r_1\eta(j)$ for all
$j\ne0$, with $r_0=0$.
For any nonzero shift $\delta$, the standard quadratic-character correlation is
\[
  \sum_{j\in\F_t}\eta(j)\eta(j+\delta)=-1,
\]
so $\paf_r(\delta)=-1$.  Writing $s_0$ for the origin value of $s$, the
constant off-origin row gives
\[
  \paf_s(\delta)=(t-2)s_*^2+2s_0s_*=t-2+2s_0s_* .
\]
The Turyn equation would require
$0=\paf_r(\delta)+\paf_s(\delta)=t-3+2s_0s_*$.  Since
$s_0s_*\in\{\pm1\}$ and $t>5$, this is impossible.
\end{proof}

\begin{proposition}[Sub-torus marginal obstruction]\label{prop:t2-marginal}
Assume the composite multiplier signs are all $+1$. If for some divisor
$t'\mid t$ the $V'$-invariant marginal system above has no solution satisfying
the stated boxes and parities, then no Turyn pair of order $t$ exists.
\end{proposition}

\begin{proof}
Every Turyn pair maps to a marginal pair $(A,B)$ with the displayed group-ring
identity, boxes, parities, and $V'$-invariance. If the finite marginal system is
infeasible for one $t'$, no such pair can have existed upstairs.
\end{proof}

\begin{corollary}[Exact secondary marginal certificate]\label{cor:exact-secondary-marginal}
In the setting of Proposition~\ref{prop:t2-marginal}, suppose the finite
orbit-value system of Lemma~\ref{lem:marginal-algebra} is exhaustively joined
over the stated boxes and parities for a divisor $t'\mid t$.  If the exact join
is empty, then no Turyn pair exists at $q$, even if some other marginal divisor
belongs to a quantified lane that passes locally.
\end{corollary}

\begin{proof}
The exhaustive join is precisely the finite system whose infeasibility is
forbidden by Proposition~\ref{prop:t2-marginal}.  Local feasibility at another
divisor does not alter the necessity of the marginal equations at this
divisor.
\end{proof}

\begin{proposition}[$h=3$ support-augmentation obstruction]\label{prop:h3-support-augmentation}
Assume the hypotheses of Proposition~\ref{prop:t2-marginal} and let
$t'=m$ be a divisor with fiber size $h=t/t'=3$, so $q=6m-1$.  Let
$\mathcal O$ be the set of $V'$-orbits on $C_m$, and let
$\mathcal O_0=\{0\}$ be the zero orbit.  For orbit-constant marginal values
$A_O,B_O$, define
\[
  N(A,B)=
  \sum_{\substack{O\in\mathcal O\\O\ne\mathcal O_0}}
  |O|\,{\bf 1}_{|A_O|=3}
  +
  \sum_{O\in\mathcal O}|O|\,{\bf 1}_{|B_O|=3}.
\]
Every marginal solution must satisfy
\[
  A_{\mathcal O_0}\in\{\pm2\},\qquad N(A,B)=\frac{m-1}{2},
  \qquad A(1)^2+B(1)^2=q,
\]
where
\[
  A(1)=\sum_{O\in\mathcal O}|O|A_O,\qquad
  B(1)=\sum_{O\in\mathcal O}|O|B_O.
\]
Consequently, if the finite support-augmentation join over
\[
  A_{\mathcal O_0}\in\{\pm2\},\quad
  A_O\in\{\pm1,\pm3\}\ (O\ne\mathcal O_0),\quad
  B_O\in\{\pm1,\pm3\}
\]
has no assignment satisfying the two displayed equations, then no Turyn pair
exists at $q$.
\end{proposition}

\begin{proof}
The coefficient at the identity in $AA^*+BB^*=q\cdot1$ is
\[
  q=\sum_{j\in C_m}A_j^2+\sum_{j\in C_m}B_j^2 .
\]
For $h=3$, the fiber boxes give
$A_{\mathcal O_0}\in\{-2,0,2\}$, all other $A$-values are odd of absolute value
at most $3$, and all $B$-values are odd of absolute value at most $3$.  If all
odd entries had magnitude $1$, the non-origin $A$ entries and all $B$ entries
would contribute $(m-1)+m=2m-1$.  Each entry of magnitude $3$ instead of
magnitude $1$ adds $8$.  Thus
\[
  6m-1=A_{\mathcal O_0}^2+(2m-1)+8N(A,B).
\]
The case $A_{\mathcal O_0}=0$ would force $N(A,B)=m/2$, impossible because
$m$ is odd.  Hence $A_{\mathcal O_0}=\pm2$ and
$N(A,B)=(m-1)/2$.  The trivial-character equation in
Lemma~\ref{lem:marginal-algebra} gives $A(1)^2+B(1)^2=q$.  These are necessary
conditions for the marginal system, so emptiness of this weaker finite join is
an obstruction by Proposition~\ref{prop:t2-marginal}.
\end{proof}

\begin{corollary}[$h=3$ support-PAF firing criterion]\label{cor:h3-support-paf}
In the setting of Proposition~\ref{prop:h3-support-augmentation}, restrict the
finite marginal system to rows satisfying
\[
  A_{\mathcal O_0}\in\{\pm2\},\qquad N(A,B)=\frac{m-1}{2},
  \qquad A(1)^2+B(1)^2=q .
\]
If no such pair also satisfies the full orbit-algebra equation
\[
  AA^*+BB^*=q\cdot1
\]
in $\Z[C_m]^{V'}$, then no Turyn pair exists at $q$.
\end{corollary}

\begin{proof}
Proposition~\ref{prop:h3-support-augmentation} shows that every marginal
solution with $h=3$ lies in this reduced finite row set.  The displayed
orbit-algebra equation is the remaining condition of
Lemma~\ref{lem:marginal-algebra}.  Hence an empty join on the reduced row set
is an empty marginal system, and Proposition~\ref{prop:t2-marginal} applies.
\end{proof}

At a prime node the support stage collapses to a single congruence.

\begin{corollary}[Prime support-sum closed form]\label{cor:h3-support-sum-prime}
In the setting of Proposition~\ref{prop:h3-support-augmentation}, suppose
$t'=\ell$ is prime and let $w$ be the order of the image of $V'$ in
$(\Z/\ell)^\times$; since that image acts freely on the nonzero residues,
every nonzero orbit has size $w$.
The support condition $A_{\mathcal O_0}=\pm2$, $N(A,B)=(\ell-1)/2$ is
satisfiable if and only if
\[
  \frac{\ell-1}2\bmod w\in\{0,1\}.
\]
In particular, if $(\ell-1)/2\bmod w\notin\{0,1\}$, the $h=3$ marginal at
$\ell$ fires and no Turyn pair exists at $q=6\ell-1$; and if $V_\ell$ is
transitive ($w=\ell-1$, $\ell\ge5$), the condition always fails, so every
transitive $h=3$ node fires.
\end{corollary}

\begin{proof}
All nonzero orbits have size $w$ and the $B$-origin has size $1$, so
$N(A,B)=wj+\beta$ with $j$ the number of magnitude-$3$ nonzero orbits over
both sides, $0\le j\le2(\ell-1)/w$, and $\beta={\bf1}_{|B_0|=3}\in\{0,1\}$.
Thus $N(A,B)=(\ell-1)/2$ is solvable iff $(\ell-1)/2\equiv\beta\pmod w$ for
some $\beta\in\{0,1\}$; the range constraint never binds because
$((\ell-1)/2-\beta)/w<2(\ell-1)/w$.  In the transitive case one has
$(\ell-1)/2\bmod(\ell-1)=(\ell-1)/2\ge2$ for $\ell\ge5$.
\end{proof}

The recorded support-sum kills are its instances: $q=425$ at $\ell=71$
($w=10$, $35\bmod10=5$) and $q=1625$ at $\ell=271$ ($w=18$,
$135\bmod18=9$), while $q=1445$ at $\ell=241$ ($w=20$, $120\bmod20=0$)
passes the support stage and is closed by the PAF layer
(Proposition~\ref{prop:q1445-h3-paf}); the transitive instances
$q=785$ at $131$ and $q=905$ at $151$ fire with no computation.  The
closed form is machine-checked against the support DP on every prime
$h=3$ node: $19$ nodes at $q\le2000$ and $71$ at $q\le10000$, with no
disagreement.

Note that the support congruence at prime $\ell$ fires automatically
whenever $d=(\ell-1)/w$ is odd and $w\ge4$, since
$(\ell-1)/2=(w/2)d\equiv w/2\pmod w$.  At composite $t'$ the same
counting argument gives a weaker congruence.

\begin{corollary}[Composite support congruence]\label{cor:h3-support-sum-gcd}
In the setting of Proposition~\ref{prop:h3-support-augmentation}, let $g$
be the greatest common divisor of the nonzero orbit sizes of
$\Omega(t',V')$ (computable from Lemma~\ref{lem:orbit-census}).  Then
$N(A,B)\equiv{\bf1}_{|B_0|=3}\pmod g$, so if
$\frac{t'-1}2\bmod g\notin\{0,1\}$, the $h=3$ marginal at $t'$ fires.
This congruence is deliberately coarse: at mixed orbit sizes $g$ is often
$2$ and the criterion is silent, while the exact subset-sum condition of
Proposition~\ref{prop:h3-support-augmentation} remains decisive.
\end{corollary}

The augmentation stage also has a closed form at prime nodes, with the
same representation-pair skeleton as Theorem~\ref{thm:two-orbit-reps}:
row sums are rigid modulo $w$, so feasibility is a residue condition on
the Gaussian representations of $q$.

\begin{proposition}[$h=3$ augmentation representation criterion]
\label{prop:h3-augmentation-reps}
Assume the setting of Proposition~\ref{prop:h3-support-augmentation} with
$t'=\ell$ prime, $w$ the order of the image of $V'$ in
$(\Z/\ell)^\times$, $d=(\ell-1)/w$, and suppose the support congruence
passes: $\beta=\frac{\ell-1}2\bmod w\in\{0,1\}$.  Put
$J=\bigl(\frac{\ell-1}2-\beta\bigr)/w$ and $b^*=3$ if $\beta=1$, else
$b^*=1$.  For $0\le j\le d$ define the exact achievable sum set
\[
  \Sigma(d,j)=\bigl\{\,3(2k-j)+\bigl(2l-(d-j)\bigr):\
  0\le k\le j,\ 0\le l\le d-j\,\bigr\}.
\]
Then the support-augmentation join of
Proposition~\ref{prop:h3-support-augmentation} is nonempty if and only if
there are a representation $q=X^2+Y^2$ with $X$ even and $Y$ odd, signs
$\epsilon_A,\epsilon_B\in\{\pm1\}$, and integers $j_A,j_B\ge0$ with
$j_A+j_B=J$, $j_A,j_B\le d$, such that
\[
  \epsilon_A X\in 2\{\pm1\}+w\,\Sigma(d,j_A),\qquad
  \epsilon_B Y\in b^*\{\pm1\}+w\,\Sigma(d,j_B).
\]
\end{proposition}

\begin{proof}
Every nonzero orbit has size $w$, so the weighted support is
$N(A,B)=w(j_A+j_B)+{\bf1}_{|B_0|=3}$, where $j_A,j_B$ count the
magnitude-$3$ orbits per side.  The support equation
$N(A,B)=\frac{\ell-1}2$ forces ${\bf1}_{|B_0|=3}=\beta$, hence
$|B_0|=b^*$ and $j_A+j_B=J$.  A side row with $j$ magnitude-$3$ orbits,
$k$ of them positive, and $l$ positive among the $d-j$ magnitude-$1$
orbits has orbit-sum contribution exactly $3(2k-j)+(2l-(d-j))$, and every
element of $\Sigma(d,j)$ arises this way; thus the achievable row sums
are $A(1)\in\pm2+w\Sigma(d,j_A)$ and $B(1)\in\pm b^*+w\Sigma(d,j_B)$.
The augmentation equation is $A(1)^2+B(1)^2=q$ with $A(1)$ even and
$B(1)$ odd (as $w$ is even), which is the displayed representation
condition up to the signs $\epsilon_A,\epsilon_B$.
\end{proof}

\begin{corollary}[$h=3$ augmentation residue criterion]
\label{cor:h3-augmentation-residue}
In the setting of Proposition~\ref{prop:h3-augmentation-reps}: if no
representation $q=X^2+Y^2$ with $X$ even and $Y$ odd satisfies
\[
  X\equiv\pm2\pmod w
  \qquad\text{and}\qquad
  Y\equiv\pm b^*\pmod w,
\]
then the $h=3$ marginal at $\ell$ fires and no Turyn pair exists at
$q=6\ell-1$.
\end{corollary}

\begin{proof}
Drop the support coupling in Proposition~\ref{prop:h3-augmentation-reps}:
membership in $\pm2+w\Sigma(d,j_A)$ forces $X\equiv\pm2\pmod w$, and
likewise for $Y$.  An empty relaxation is an empty join.
\end{proof}

The closed forms are machine-verified against the support-augmentation DP
on every evaluable prime $h=3$ node through $q\le10000$: all $71$ nodes
agree exactly.  The criterion fires on $68$ of the $71$ ($61$ already by
the residue corollary alone); the three survivors passing to the PAF
layer are $q=1445$ at $241$ (closed by
Proposition~\ref{prop:q1445-h3-paf}), and, outside the finite ledger,
$q=3845$ at $641$ and $q=3965$ at $661$, whose reduced PAF joins ($20$
and $66$ nonzero orbits) exceed present enumeration budgets.

\begin{proposition}[Closure of $q=1445$ at $t'=241$]\label{prop:q1445-h3-paf}
Let $q=1445$, so $t=723=3\cdot241$, and take the marginal divisor
$t'=241$ with $h=3$.  The $h=3$ support-PAF join of
Corollary~\ref{cor:h3-support-paf} is empty.  Therefore no Turyn pair exists
at $q=1445$.
\end{proposition}

\begin{proof}
Here the image $V'\le(\Z/241\Z)^\times$ has order $20$, and the nonzero
residues split into $12$ orbits, each of size $20$.  The $h=3$ energy equation
forces support $120=(241-1)/2$.  The augmentation equation and row-sum
congruences force
\[
  A(1)=\pm38,\qquad B(1)=\pm1 .
\]
The exact reduced join enumerates $2{,}148{,}432$ admissible $A$-rows
($1{,}074{,}216$ distinct orbit-PAF side vectors) and $2{,}257{,}816$
admissible $B$-rows.  No $B$ side vector has a complementary $A$ vector at the
required support.  Thus the reduced finite system of
Corollary~\ref{cor:h3-support-paf} is empty.
\end{proof}

\begin{proposition}[Component join criterion]\label{prop:component-join}
Assume the composite multiplier signs are all $+1$ and let
\[
  \Q[C_{t'}]^{V'}\simeq\prod_{\kappa\in\mathcal C}K_\kappa
\]
be the rational Wedderburn decomposition of the fixed algebra, with component
maps $\lambda_\kappa$.  Since $-1\in V'$, inversion is trivial on the fixed
algebra.  A $V'$-invariant marginal solution exists if and only if there are
component pairs
\[
  \lambda_\kappa(A)^2+\lambda_\kappa(B)^2=q\qquad(\kappa\in\mathcal C)
\]
whose inverse image in the orbit-sum basis lies in
$\Z[C_{t'}]^{V'}$ and satisfies the fiber boxes and parities of
Lemma~\ref{lem:marginal-algebra}.  For every real embedding
$\rho:K_\kappa\hookrightarrow\mathbb R$, any such component value obeys
$|\rho(\lambda_\kappa(A))|,|\rho(\lambda_\kappa(B))|\le\sqrt q$; hence each
component search is finite in the corresponding order.  If the finite component
join is empty for some $t'$, then no Turyn pair exists at $q$.
\end{proposition}

\begin{proof}
The semisimple algebra $\Q[C_{t'}]^{V'}$ is the product of the character
components fixed by $V'$.  Because $-1\in V'$, each fixed element is also fixed
by inversion, so the image of $AA^*+BB^*=q\cdot1$ in every component is exactly
$\lambda_\kappa(A)^2+\lambda_\kappa(B)^2=q$.  Conversely, component values that
invert to integral orbit coefficients satisfying the boxes/parities give an
element of the marginal system.  Applying any real embedding to the component
equation gives a sum of two real squares equal to $q$, which gives the displayed
bound and finiteness.  The final obstruction is
Proposition~\ref{prop:t2-marginal}.
\end{proof}

\begin{corollary}[Full-torus sign firing criterion]\label{cor:full-torus-sign}
Assume the hypotheses of Proposition~\ref{prop:t2-marginal} and take the full
torus $t'=t$.  If $t$ is prime and the multiplier signs are all $+1$, then
the fiber parameter is $h=1$ and the marginal boxes force
\[
  A_0=0,\quad A_j\in\{\pm1\}\ (j\ne0),\qquad
  B_j\in\{\pm1\}\quad(0\le j<t),
\]
constant on the $V$-orbits.  Hence full-torus existence is decided by the
finite sign join
\[
  AA^*+BB^*=q\cdot1,\qquad A(1)^2+B(1)^2=q
\]
over these orbit-sign rows.  If this finite join is empty, then no Turyn pair
exists at $q$.
\end{corollary}

\begin{proof}
For $t'=t$ one has $h=1$ in the fiber boxes of
Lemma~\ref{lem:marginal-algebra}; the $A$-origin is therefore forced to be
$0$, every other $A$ coefficient is odd of absolute value at most $1$, and
every $B$ coefficient is odd of absolute value at most $1$.  The composite
multiplier theorem makes the rows constant on $V$-orbits.  The displayed
finite join is precisely the marginal algebra equation and the trivial
character equation from Lemma~\ref{lem:marginal-algebra}, so emptiness is an
instance of Proposition~\ref{prop:t2-marginal}.
\end{proof}

The sign join is finite but can be wide.  For prime $t$ the trivial
character alone already decides the join whenever the multiplier image is
large: the row sums of $V$-invariant sign rows are so rigid that the
two-square equation overshoots.  This gives the prime-$t$ full-torus branch
a zero-computation firing criterion, and, read contrapositively, a
square-root bound on the multiplier image of any genuine Turyn pair of
prime order.

\begin{theorem}[Full-torus row-sum bound]\label{thm:full-torus-rowsum}
Let $q=2t-1$ with $t\ge3$ prime and all composite multiplier signs $+1$.
Let $w$ be the order of the image of $V$ in $(\Z/t)^\times$ ($w$ is even
because $-1\in V$), and let $d=(t-1)/w$.  If
\[
  (w-1)^2\ge t,\qquad\text{equivalently}\qquad w\ge d+2,
\]
then the full-torus sign join of Corollary~\ref{cor:full-torus-sign} is
empty, and no Turyn pair of order $t$ exists.  Contrapositively: if a
Turyn pair of prime order $t$ exists, then $w<1+\sqrt t$.
\end{theorem}

\begin{proof}
Since $t$ is prime, $V$ acts freely on the nonzero residues, so there are
exactly $d$ nonzero orbits, all of size $w$, and $wd=t-1$.  The stated
equivalence is
$w\ge d+2\iff w(w-1)\ge wd+w=t-1+w\iff w^2-2w\ge t-1\iff(w-1)^2\ge t$.

At $h=1$ the boxes force sign rows: $A_0=0$, $A$ equal to a sign on each
nonzero orbit, $B$ a sign everywhere.  Hence
\[
  \lambda_1(A)=w\,s_A,\qquad
  \lambda_1(B)=\beta+w\,s_B,
\]
where $s_A$ and $s_B$ are sums of $d$ signs --- so
$s_A\equiv s_B\equiv d\pmod2$ and $|s_A|,|s_B|\le d$ --- and
$\beta=B_0=\pm1$.  The trivial-character equation of
Lemma~\ref{lem:marginal-algebra} reads
\[
  w^2s_A^2+(\beta+w\,s_B)^2=q=2wd+1 .
\]
Put $\sigma=\beta s_B$, of the same parity and range as $s_B$.  Expanding
and dividing by $2w$,
\[
  w\,(s_A^2+\sigma^2)=2(d-\sigma).
\]
If $d$ is even, then $s_A,\sigma$ are even; $(s_A,\sigma)=(0,0)$ would
force $d=0$, so $s_A^2+\sigma^2\ge4$ and
$4w\le2(d-\sigma)\le4d$, giving $w\le d$.  If $d$ is odd, then
$s_A,\sigma$ are odd, so $s_A^2+\sigma^2\ge1+\sigma^2$, and $w\ge d+2$
would give
\[
  2(d-\sigma)\;=\;w\,(s_A^2+\sigma^2)\;\ge\;(d+2)(1+\sigma^2),
\]
i.e.\ $0\ge d(\sigma^2-1)+2(\sigma^2+\sigma+1)$, false because
$\sigma^2\ge1$ and $\sigma^2+\sigma+1>0$.  In both parities the equation
has no solution when $w\ge d+2$, so the sign join is empty, and
Proposition~\ref{prop:t2-marginal} applies.
\end{proof}

\begin{remark}[Sharpness and the small-image regime]\label{rem:rowsum-sharp}
(i) The bound is sharp.  For odd $d$ and $w=d+1$ the displayed equation is
solvable: $s_A=\pm1$, $\sigma=-1$ gives $2w=2(d+1)$, and then
$q=2wd+1=w^2+(w-1)^2$.  Turyn's pair at $q=25$ realizes the boundary:
$t=13$, the image of $\langle M(25),-1\rangle$ in $(\Z/13)^\times$ has
$w=4$, $d=3$, and $25=4^2+3^2$.  The row-sum bound thus separates the
construction regime from the obstruction regime exactly at $w=d+1$.
(ii) The case $d=1$ recovers Corollary~\ref{cor:transitive-full-torus}:
$w=t-1$ and $(t-2)^2\ge t$ for $t\ge4$.
(iii) On the finite panel the theorem closes the entire prime-$t$
full-torus branch with zero computation: all $29$ branch members at
$q\le2000$ satisfy $w\ge d+2$, including the higher-degree joins $q=697$
($w=58$, $d=6$), $q=1521$ ($w=76$, $d=10$), and $q=1765$ ($w=126$,
$d=7$), previously closed by complete integer MITM.
(iv) What remains at prime-$t$ top nodes is the small-image regime
$w\le d+1$, i.e.\ $w<1+\sqrt t$.  This is exactly where Turyn's
construction lives ($q=p^e$ makes $\langle p,-1\rangle$ small), so no
row-sum-type refinement can close it uniformly.  The regime does contain
multi-prime composites: the members at $q\le30000$ are
$q=5185=5\cdot17\cdot61$ ($t=2593$, $w=8$, $d=324$), $q=10045$
($t=5023$, $w=54$), $q=13833$ ($t=6917$, $w=26$), and $q=15205$
($t=7603$, $w=42$).  Thus no uniform composite-$q$ lower bound on the
multiplier image can exist.  Three of the four fall to the exact form
of the trivial-character predicate
(Corollary~\ref{cor:full-torus-exact-char} below), so the regime
demanding genuinely non-multiplier input is thinner than the threshold
suggests; its one surviving member in range, $q=5185$, falls to the
fold-tower/stabilizer-amplification mechanism below (GRH-conditional
through one class-group computation).
Novelty of the row-sum mechanism relative to classical multiplier-size
arguments is not audited here; the closures stand regardless.
\end{remark}

The proof of Theorem~\ref{thm:full-torus-rowsum} in fact establishes a
predicate strictly stronger than the threshold: the displayed scalar
equation is \emph{necessary} for the sign join, so its exact emptiness
fires below the threshold as well.

\begin{corollary}[Exact trivial-character firing]
\label{cor:full-torus-exact-char}
In the setting of Theorem~\ref{thm:full-torus-rowsum}, if the equation
\[
  w\,(s^2+\sigma^2)=2(d-\sigma)
\]
has no integer solution with $s\equiv\sigma\equiv d\pmod2$ and
$|s|,|\sigma|\le d$, then the full-torus sign join is empty and no
Turyn pair of order $t$ exists.  For $w\ge d+2$ no solution exists for
any parameters (Theorem~\ref{thm:full-torus-rowsum}); below the
threshold the exact check is a finite $O(d)$ enumeration and can still
fire.
\end{corollary}

\begin{proof}
The equation was derived in the proof of
Theorem~\ref{thm:full-torus-rowsum} as a consequence of the
trivial-character equation on $V$-invariant sign rows; emptiness of its
admissible solution set therefore empties the sign join, and
Proposition~\ref{prop:t2-marginal} applies.
\end{proof}

The exact form is decisive on the small-image census: of the four
members at $q\le30000$ listed in Remark~\ref{rem:rowsum-sharp}(iv),
three fire below the threshold --- $q=10045$ ($w=54$, $d=93$),
$q=13833$ ($w=26$, $d=266$), and $q=15205$ ($w=42$, $d=181$) all have
empty admissible solution sets (complete enumeration; each is all-plus
and blind-spot, independently re-verified) --- while $q=5185$ genuinely
passes, with $(s,\sigma)=(8,4)$ realizing $5185=64^2+33^2$.  The census
extends unconditionally to $q\le100000$: of $1027$ blind-spot prime-$t$
all-plus members, $1011$ die by the threshold, $9$ more by the exact
predicate, and exactly four survive:
\[
  q=5185,\quad 37825=5^2\cdot17\cdot89,\quad
  62305=5\cdot17\cdot733,\quad 93721=17\cdot37\cdot149,
\]
all, curiously, divisible by $17$.  Two of the four are decided by the
fold-tower mechanism below
(Propositions~\ref{prop:q5185-closure},~\ref{prop:q62305-closure}); the
other two are the sharp wall witnesses of
Remark~\ref{rem:fold-gaps}.

\paragraph{Small-image top nodes: stabilizer amplification.}
In the small-image regime the multiplier group is too small for any
$M(q)$-invariance argument, but it is not static: class-group input can
force partial Galois stability of the ideal generated by the character
values, and stability converts --- by the next lemma --- into
\emph{honest new multipliers}, re-entering the regime where the row-sum
bound fires.  Throughout, $t$ is prime, the multiplier signs are all
$+1$, $K_V\subset\Q(\zeta_t)$ is the fixed field of $V$, and
$L=K_V(i)$, a CM field with $\mu_L=\langle i\rangle$ and
$\Gal(L/\Q(i))\simeq(\Z/t)^\times/V$, cyclic of order
$d=(t-1)/w$.  A Turyn pair gives the
Weil number $\gamma=\chi(\tau)\in\mathcal O_L$ (any character $\chi$ of
order $t$), with $\gamma\bar\gamma=q$; for squarefree $q$ the ideal
$(\gamma)$ selects exactly one prime from each conjugate pair above each
$p\mid q$, a \emph{pattern ideal}.

\begin{lemma}[Stabilizer rigidity]\label{lem:stabilizer-rigidity}
Let $t\ge9$ be prime with all multiplier signs $+1$, and let
$\sigma=\sigma_u\in\Gal(L/\Q(i))$ satisfy $\sigma((\gamma))=(\gamma)$.
Then $\sigma(\gamma)=\gamma$, and $u$ is an honest multiplier of the
pair: $\tau^{(u)}=\tau$.
\end{lemma}

\begin{proof}
$\sigma(\gamma)/\gamma$ is an algebraic integer all of whose archimedean
absolute values equal $1$ (numerator and denominator are Weil
$q$-numbers), hence a root of unity in $L$ (Kronecker), so
$\sigma(\gamma)=i^a\gamma$ with $a\in\Z/4$.  Applying this identity to
every character conjugate and inverting the Fourier transform on $C_t$
gives
\[
  \tau^{(u)}=i^a\tau+\frac{(1-i^a)\,\chi_0(\tau)}{t}\sum_{g\in C_t}g .
\]
Integrality of the correction requires $t\mid(1-i^a)\chi_0(\tau)$ in
$\Z[i]$.  For $a\ne0$ the norm of the left side is
$N(1-i^a)\cdot q\le4(2t-1)<t^2$ for $t\ge9$, while
$(1-i^a)\chi_0(\tau)\ne0$ because $\chi_0(\tau)\bar\chi_0(\tau)=q\ne0$;
so no such nonzero multiple of $t$ exists, forcing $a=0$ and
$\tau^{(u)}=\tau$.
\end{proof}

\begin{proposition}[The $F_4$ class-group cut at $q=5185$]
\label{prop:q5185-f4-cut}
Let $q=5185=5\cdot17\cdot61$, $t=2593$, $w=8$.  In $L$ (degree $648$)
the pattern ideal factors as: a pinned Gaussian $61$-part
($(5\pm6i)\mathcal O_L$; one conjugate pair), a $17$-pattern on
$C_4=\Gal(L/\Q(i))/D_{17}$ (four pairs), and a $5$-pattern on $C_{81}$
($81$ pairs).  Let $F_4=k_4(i)$ be the degree-$8$ fold field, $k_4$ the
quartic subfield of $\Q(\zeta_t)$.  Assume the class group computed by
\texttt{bnfinit} under GRH:
$\Cl(F_4)=C_{408}\times C_2\times C_2$ of order $1632$ --- consistent
with the \emph{unconditional} analytic minus class number
$h^-(F_4)=1632\,Q$ computed by generalized Bernoulli numbers, which pins
$Q=1$, $h(k_4)=1$.  Then: both primes above $5$ and above $61$ are
principal in $F_4$; the four $17$-side prime classes
$v_1,\dots,v_4$ satisfy $\sum v_j=0$ (exactly: their product is
$(1+4i)\mathcal O_{F_4}$, verified as an ideal identity) and all four
first coordinates are $\equiv1\pmod4$; hence the folded ideal of a
pattern with sign set $S$ has class $162\sum_{j\in S}v_j\equiv
2|S|\pmod 8$ in the $C_{408}$ coordinate, nonzero for every proper
nonempty $S$.  Consequently every non-constant $17$-pattern gives a
non-principal folded ideal, and the $17$-pattern of any Turyn pair at
$q=5185$ is constant.  The order-$8$ character of $C_{408}$ is the
minimal certifying quotient: an order-$4$ character misses $|S|=2$ and
the order-$17$ character misses two pairs.
\end{proposition}

\begin{proof}
The fold data and the fourteen non-constant folded ideals were computed
exactly in PARI/gp and verified two independent ways: the class-linear
arithmetic displayed, and direct \texttt{bnfisprincipal} on each exact
folded ideal $(\prod_{j\in S}P_j)^{162}$ (all non-principal, classes
matching).  Existence of the pattern ideal downstairs forces the folded
ideal principal, which by the above forces $S\in\{\emptyset,\text{all}\}$:
a constant pattern.
\end{proof}

\begin{proposition}[Closure of $q=5185$, conditional on $\Cl(F_4)$]
\label{prop:q5185-closure}
Assume the class group of Proposition~\ref{prop:q5185-f4-cut}.  Then no
Turyn pair of order $t=2593$ exists.
\end{proposition}

\begin{proof}
By Proposition~\ref{prop:q5185-f4-cut} the $17$-pattern is constant, so
the stabilizer of $(\gamma)$ in $\Gal(L/\Q(i))\simeq C_4\times C_{81}$
contains $C_4\times\{0\}$: the $61$-part is always stable, the constant
$17$-part is stable, and $C_4\times\{0\}$ is exactly the decomposition
group $D_5$, which acts trivially on the $81$ pairs above $5$ --- so the
$5$-pattern imposes no condition.  (This coincidence, $D_5$ equal to the
$17$-pattern space, is what makes this top node decidable at one fold
level.)  By Lemma~\ref{lem:stabilizer-rigidity}, every element of
$C_4\times\{0\}$ is an honest multiplier, so the pair's true multiplier
image contains the unique order-$32$ subgroup of $(\Z/t)^\times$, i.e.\
its order $w'$ lies in $\{32,96,288,864,2592\}$.  All die: at $w'=32$
the trivial-character equation reads $1024\,s_A^2+\lambda_1(B)^2=5185$
with $s_A$ odd, impossible ($5185-1024=4161$ is not a square, and
$s_A^2\ge9$ overshoots); every $w'\ge96$ satisfies $(w'-1)^2\ge t$ and
dies by Theorem~\ref{thm:full-torus-rowsum}.
\end{proof}

\begin{remark}[Conditionality ledger and the general mechanism]
\label{rem:small-image-conditional}
Every link in the $q=5185$ chain is unconditional exact arithmetic ---
the pattern normal form, the pinned trivial character
$5185=64^2+33^2$ (the unique representation compatible with the
$V$-invariant row sums, and the unique one divisible by $5+6i$), the
stabilizer lemma, the exact ideal identities, and the amplification
cascade --- except the single input $\Cl(F_4)$, which currently rests
on \texttt{bnfinit} (GRH/Bach).  \texttt{bnfcertify} passed for the
smaller fold field $F_3$ but extrapolates to weeks for $F_4$
(discriminant $\approx7.8\cdot10^{22}$); two upgrade paths are
documented: the $h^-$ pincer (unconditional $h^-\ge1632$ plus a flag-$1$
quotient certification pins $\Cl(F_4)$ exactly) and an explicit octic
class-field certificate (necessarily octic: $C_8$ is the minimal
certifying quotient, and $\zeta_8\notin F_4$ makes it a Kummer-descent
construction).  Until one lands, $q=5185$ carries the same
GRH-conditional label as the class-group layer of the prime-quotient
program, in contrast to the field-free ledger of
\S\ref{sec:finite-panel-accounting}.  The general lemma to extract, the
\emph{fold-tower/stabilizer-amplification dichotomy}, is: if some prime
$p\mid q$ has decomposition group equal to the pattern space of another
prime $p'$, and the fold level $L^{D_{p'}}$ forces the $p'$-pattern
constant, then the multiplier image amplifies by the full pattern space
of $p'$ and the trivial-character/row-sum cascade applies at the
amplified level.  Quantifying when the fold level fires (here: prime
classes $\equiv1\bmod4$ in the relevant cyclic quotient) is the next
theorem obligation for this family; novelty is not yet audited.
\end{remark}

The mechanism is now stated once, with its single class-group input
isolated as a hypothesis, and it has a second instance.

\begin{proposition}[Fold-amplification dichotomy]
\label{prop:fold-dichotomy}
Let $q=2t-1$ be squarefree with every prime divisor $\equiv1\bmod4$,
$t\ge9$ prime, all multiplier signs $+1$, and suppose the exact
predicate of Corollary~\ref{cor:full-torus-exact-char} is solvable at
$w=|V|$.  Write $C=\Gal(L/\Q(i))\simeq(\Z/t)^\times/V$, let
$D_p\subseteq C$ be the decomposition group of $p\mid q$ and
$X_p=C/D_p$ its pattern space.  Fix a target prime $p'\mid q$, put
$F=L^{D_{p'}}=k_e(i)$ with $e=|X_{p'}|$, and
\[
  U(p')=\bigcap_{p''\ne p'}D_{p''}.
\]
Assume the \emph{firing condition} at $F$: for every non-constant
$p'$-pattern, no choice of the other primes' fold weights makes the
pushed-down pattern ideal principal in $\Cl(F)$.  Then any Turyn pair
of order $t$ has multiplier image containing the subgroup of order
$w\cdot|U(p')|$, and if the exact predicate of
Corollary~\ref{cor:full-torus-exact-char} is empty at every order $w'$
with $w|U(p')|\mid w'\mid t-1$, no Turyn pair of order $t$ exists.
Every step except the evaluation of the firing condition --- which
requires $\Cl(F)$ and carries that computation's certification tier
--- is unconditional exact arithmetic.
\end{proposition}

\begin{proof}
A pair gives the Weil number $\gamma\in\mathcal O_L$ with
$\gamma\bar\gamma=q$; squarefreeness makes
$v_P(\gamma)+v_{\bar P}(\gamma)=1$ for every prime $P\mid q$ of $L$, so
$(\gamma)$ is a pattern ideal.  Pushing down by $N_{L/F}$ forces the
folded pattern ideal principal in $\Cl(F)$; the $p'$-pairs fold
bijectively, while each other prime contributes a free fold weight.
The firing condition therefore forces the $p'$-pattern constant.  A
constant $p'$-pattern is $\Gal$-stable, every $D_{p''}$ acts trivially
on its own pattern space, and the $61$-type pinned parts are always
stable; hence $\mathrm{Stab}_C((\gamma))\supseteq U(p')$, and
Lemma~\ref{lem:stabilizer-rigidity} converts each stabilizing element
into an honest multiplier.  The enlarged image has order divisible by
$w|U(p')|$, and Corollary~\ref{cor:full-torus-exact-char} at the
enlarged image gives the final contradiction.
\end{proof}

\begin{proposition}[Closure of $q=62305$, conditional on
$\Cl(k_3(i))$]\label{prop:q62305-closure}
Let $q=62305=5\cdot17\cdot733$, $t=31153$, $w=44$.  Take $p'=5$
($e=3$, $F=k_3(i)$ of degree $6$, discriminant
$t^4\cdot4^3\approx6\cdot10^{19}$) and assume the class group computed
by \texttt{bnfinit} under GRH: $\Cl(F)=C_{5799}$ --- matching the
\emph{unconditional} analytic minus class number $h^-(F)=5799$
exactly, which pins $Q=1$ and $h(k_3)=1$.  Then the rescue-class set
generated by the $17$- and $733$-fold weights has size one, only the
two constant $5$-patterns are rescued, and the firing condition of
Proposition~\ref{prop:fold-dichotomy} holds.  With
$U=D_{17}\cap D_{733}$ of order $3$, the multiplier image contains the
order-$132$ subgroup, and all six admissible amplified orders are
empty for the exact predicate.  Hence no Turyn pair of order
$t=31153$ exists, conditional exactly on the one \texttt{bnfinit}
input.
\end{proposition}

\begin{proof}
Fold data (exact): the three conjugate pairs above $5$ fold
bijectively with exponent $236$; $17$ contributes weights in $[0,4]$
with exponent $59$; $733$ weights in $[0,59]$ with exponent $4$.  The
rescue computation is exact linear arithmetic in $C_{5799}$, validated
on exact ideals in PARI/gp: the identity
$\prod_jP_j=(1+2i)\mathcal O_F$ holds, direct
\texttt{bnfisprincipal} on the folded ideal of a non-constant pattern
with $15$ sampled rescue weights finds no principal instance, and the
constant pattern is rescued at weight $(0,0)$.
Lemma~\ref{lem:stabilizer-rigidity} applies since $4q<t^2$; the six
amplified levels are the orders $w'$ with $132\mid w'\mid31152$, each
checked empty by complete enumeration of the exact predicate.
\end{proof}

\begin{remark}[The two named gaps, and the wall witnesses]
\label{rem:fold-gaps}
(i) \emph{The Stickelberger gap.}  The firing condition cannot yet be
stated as arithmetic of $(q,t,w,p')$: both instances reduce to ``the
$p'$-side degree-one prime classes share a nonzero image in a small
cyclic quotient of $\Cl^-(k_e(i))$ and the rescue set degenerates,''
but predicting either fact currently requires the class group per
instance.  The missing lemma is an explicit Stickelberger/reflection
description of the subgroup of $\Cl^-(k_e(i))$ generated by the
degree-one primes above $p'$, computable from $(t,e,p')$ alone; until
it exists, each instance costs one \texttt{bnfinit} and its
certification.  (ii) \emph{Wall witnesses.}  The dichotomy cannot
close the other two census survivors.  At $q=37825=5^2\cdot17\cdot89$
every nontrivial-$U$ fold field is infeasible (degrees $48$ and
$2364$), $p'=5$ has $U=1$, and $5^2\,\|\,q$ needs the valuation-vector
generalization of the pattern normal form.  At
$q=93721=17\cdot37\cdot149$ ($w=2$, the minimal possible image) the
only near-feasible fold target $p'=37$ has amplified orders
$\{10,20,30,60\}$ that \emph{survive} the exact predicate --- so even
a perfect class-group cut cannot close it: the dichotomy is provably
insufficient there, and a refinement (partial folds at intermediate
levels $k_{e'}(i)$, forcing fold weights to extremes) is the recorded
direction.  These two orders are the correct stress tests for any
strengthened small-image criterion.
\end{remark}

\begin{proposition}[Prime sub-torus component criterion]\label{prop:prime-subtorus-component}
Assume the composite multiplier signs are all $+1$, let $\ell\mid t$ be prime,
and let $H$ be the image of $M(q)$ in $(\Z/\ell\Z)^\times$.  Put
$K=\Q(\zeta_\ell)^H$.  Then
\[
  \Q[C_\ell]^H\simeq \Q\times K,
\]
where the first factor is the trivial component and the second is expressed in
the Gaussian-period orbit basis attached to the $H$-orbits on
$(\Z/\ell\Z)^\times$.  A prime sub-torus marginal solution exists if and only if
there are component pairs
\[
  X_1^2+Y_1^2=q,\qquad X_K^2+Y_K^2=q\quad(X_K,Y_K\in\mathcal O_K)
\]
whose inverse period transform gives integral orbit values satisfying the
fiber boxes and parities.  If this finite component join is empty for this
prime divisor $\ell$, then no Turyn pair exists at $q$.
\end{proposition}

\begin{proof}
For prime $\ell$ there is the rational decomposition
$\Q[C_\ell]\simeq \Q\times\Q(\zeta_\ell)$.  Taking $H$-fixed points gives
$\Q[C_\ell]^H\simeq\Q\times\Q(\zeta_\ell)^H$.  The nontrivial fixed field has
degree equal to the number of nonzero $H$-orbits, and the corresponding period
sums give the orbit-basis component transform.  Proposition~\ref{prop:component-join}
then gives the two component equations, the embedding bounds, the inverse
integrality condition, and the fiber boxes/parities.  An empty component join
therefore contradicts Proposition~\ref{prop:t2-marginal}.
\end{proof}

\begin{proposition}[Prime sub-torus period transform]\label{prop:prime-subtorus-periods}
In the setting of Proposition~\ref{prop:prime-subtorus-component}, write
$m=|H|$ and $d=(\ell-1)/m$, fix a primitive root $g$ modulo $\ell$, and let
\[
  \eta_k=\sum_{r\in g^kH}\zeta_\ell^{\,r}\qquad(0\le k<d)
\]
be the Gaussian periods.  Then $\eta_0,\dots,\eta_{d-1}$ is an integral basis
of $\mathcal O_K$, each $\eta_k$ is real, the Galois group of $K$ is cyclic
with generator $\sigma\colon\eta_k\mapsto\eta_{k+1\bmod d}$, and
$\eta_0+\cdots+\eta_{d-1}=-1$.  For orbit values
$X=x_0E_{\{0\}}+\sum_kx_kE_{g^kH}$ the two components of
Proposition~\ref{prop:prime-subtorus-component} are
\[
  \lambda_1(X)=x_0+m\sum_kx_k,
  \qquad
  \lambda_\ell(X)=\sum_ka_k\eta_k
  \quad\text{with }a_k=x_k-x_0 .
\]
Conversely, a pair $(\lambda_1,(a_k))\in\Z\times\Z^d$ arises from integer
orbit values if and only if
\[
  \lambda_1\equiv m\sum_ka_k\pmod\ell,
\]
in which case $x_0=(\lambda_1-m\sum_ka_k)/\ell$ and $x_k=x_0+a_k$.  The fiber
parities are equivalent to: on the $A$ side every $a_k$ is odd and
$\lambda_1(A)$ is even; on the $B$ side every $a_k$ is even and $\lambda_1(B)$
is odd.  The fiber boxes are $|x_0|\le h-1$ and $|x_0+a_k|\le h$.
\end{proposition}

\begin{proof}
Since $\ell$ is an odd prime, $\Q(\zeta_\ell)/\Q$ is tamely ramified, and the
Galois conjugates $\eta_k$ of the trace $\eta_0=\mathrm{Tr}_{\Q(\zeta_\ell)/K}
(\zeta_\ell)$ form a normal integral basis of $K$; they are real because
$-1\in H$ makes every coset negation-closed, and their sum is
$\sum_{r\ne0}\zeta_\ell^r=-1$.  Evaluating $X$ at the trivial character gives
$\lambda_1$; evaluating at $\zeta_\ell$ gives $x_0+\sum_kx_k\eta_k=
\sum_k(x_k-x_0)\eta_k$ after substituting $1=-\sum_k\eta_k$.  From
$\lambda_1=x_0+m\sum_k(x_0+a_k)=\ell x_0+m\sum_ka_k$ the inverse formula and
the congruence follow.  For the parities: on the $A$ side $x_0$ is even and
every $x_k$ is odd, so every $a_k$ is odd, and
$m\sum_kx_k\equiv md\equiv0\pmod2$ because $md=\ell-1$ is even, so
$\lambda_1\equiv x_0\pmod 2$ is even; the $B$ side is the same computation
with every value odd.  The boxes restate the fiber boxes of
Lemma~\ref{lem:marginal-algebra}.
\end{proof}

\begin{corollary}[Prime sub-torus firing criterion]\label{cor:prime-subtorus-firing}
Assume the hypotheses of Proposition~\ref{prop:prime-subtorus-periods}.  Let
$Z(q)=\{(L,M)\in\Z^2:L^2+M^2=q\}$ and let
\[
  K_\ell(q)=\Bigl\{(\alpha,\beta)\in\bigl(\textstyle\bigoplus_k\Z\eta_k\bigr)^2:
  \alpha^2+\beta^2=q\Bigr\}.
\]
Applying each real place to $\alpha^2+\beta^2=q$ gives
$|\sigma^r(\alpha)|,|\sigma^r(\beta)|\le\sqrt q$ for $0\le r<d$, so
$K_\ell(q)$ is finite and computable inside the period coordinate box.  Let
$\mathcal S_A(q,\ell,h)$ be the set of pairs $(L,\alpha)$ with $L$ a first
coordinate of $Z(q)$ and $\alpha=\sum_ka_k\eta_k$ a first coordinate of
$K_\ell(q)$ with every $a_k$ odd and
$L\equiv m\sum_ka_k\pmod\ell$, whose inverse transform under
Proposition~\ref{prop:prime-subtorus-periods} satisfies the $A$ boxes and
parities; define $\mathcal S_B(q,\ell,h)$ likewise with every $a_k$ even and
the $B$ boxes and parities.  If $\mathcal S_A=\varnothing$, or
$\mathcal S_B=\varnothing$, or no $(L,\alpha)\in\mathcal S_A$ has a mate
$(M,\beta)\in\mathcal S_B$ with $(L,M)\in Z(q)$ and
$(\alpha,\beta)\in K_\ell(q)$, then the prime sub-torus marginal at $\ell$ is
impossible, and hence no Turyn pair exists at $q$.
\end{corollary}

\begin{proof}
By Proposition~\ref{prop:prime-subtorus-periods}, a marginal solution is
exactly a matched pair of component halves whose inverse transforms are
integral, boxed, and correctly oriented; the component values of a marginal
lie in the period lattice because their coordinates are the integers
$x_k-x_0$.  An empty join is therefore an empty marginal, and the obstruction
is Proposition~\ref{prop:t2-marginal}.
\end{proof}

For quadratic lanes the period solution set has an explicit divisor
parametrization.  Degree $d=2$ forces $\ell\equiv1\pmod4$, because $H$ is then
the subgroup of squares and $-1\in H$; the component field is
$K=\Q(\sqrt\ell)$ with $\eta_0,\eta_1=(-1\pm\sqrt\ell)/2$.

\begin{lemma}[Quadratic two-square parametrization]\label{lem:quadratic-two-square}
Let $\ell\equiv1\pmod4$ be prime and let $q\equiv1\pmod4$.  Period coordinate
vectors of constant parity are exactly the elements of $\Z[\sqrt\ell]$: if
$a_0\equiv a_1\pmod2$, then $\alpha=a_0\eta_0+a_1\eta_1=U+V\sqrt\ell$ with
$U=-(a_0+a_1)/2$ and $V=(a_0-a_1)/2$, and conversely $a_0=V-U$, $a_1=-U-V$.
Under this change of coordinates the $A$-side parity of
Corollary~\ref{cor:prime-subtorus-firing} is $U\not\equiv V\pmod2$, the
$B$-side parity is $U\equiv V\pmod2$, and the congruence
$L\equiv m\sum_ka_k\pmod\ell$ becomes $L\equiv U\pmod\ell$.  Moreover the
pairs $(\alpha,\beta)=(U_A+V_A\sqrt\ell,\;U_B+V_B\sqrt\ell)$ in
$\Z[\sqrt\ell]^2$ with $\alpha^2+\beta^2=q$ are exactly:
\begin{itemize}
\item the rational pairs $(V_A,V_B)=(0,0)$ with $(U_A,U_B)\in Z(q)$; and
\item for every factorization $q=ef$ in which $e=p^2+r^2$ is a primitive
  two-square representation and $f=k^2+\ell c^2$ with $c\ne0$, the pairs
  \[
    (U_A,U_B)=k\,(-r,p),\qquad(V_A,V_B)=c\,(p,r),
  \]
  where $(p,r)$ runs over all primitive representations of $e$ and
  $k\in\Z$, $c\in\Z\setminus\{0\}$ run over all representations of $f$.
\end{itemize}
\end{lemma}

\begin{proof}
The coordinate identities are direct computation, using
$m\sum_ka_k=\tfrac{\ell-1}2(-2U)=-(\ell-1)U\equiv U\pmod\ell$.  Expanding
$\alpha^2+\beta^2=q$ in $\Z[\sqrt\ell]$ gives the pair of equations
\[
  U_A^2+U_B^2+\ell(V_A^2+V_B^2)=q,
  \qquad
  U_AV_A+U_BV_B=0 .
\]
Write $z=(U_A,U_B)$ and $w=(V_A,V_B)$ in $\Z^2$.  If $w=0$, the equations say
$z\in Z(q)$.  If $w\ne0$, write $w=cw_0$ with $w_0=(p,r)$ primitive and
$c\ne0$; the orthogonality $z\cdot w_0=0$ forces $z=k(-r,p)$ for some
$k\in\Z$, and substituting gives $(p^2+r^2)(k^2+\ell c^2)=q$.  Primitivity of
$w_0$ makes $e=p^2+r^2$ a primitive representation, and every displayed pair
conversely satisfies both equations.
\end{proof}

In particular a quadratic prime sub-torus marginal has a component half with
$V\ne0$ on either side only if some divisor of $q$ is represented by
$x^2+\ell y^2$ with $y\ne0$ and cofactor a primitive sum of two squares;
otherwise both side sets of Corollary~\ref{cor:prime-subtorus-firing} consist
of rational halves $(L,U)$ with $L$ even and $U$ odd on the $A$ side, $L$ odd
and $U$ even on the $B$ side, and $L\equiv U\pmod\ell$ on each side.  The
firing question then reduces to a residue-matching condition modulo $\ell$
between the even and odd first coordinates of $Z(q)$, together with the
fiber boxes.  For example, at $q=1189=29\cdot41$, $\ell=17$: no divisor of
$q$ is represented by $x^2+17y^2$ with $y\ne0$, the even first coordinates
$\{\pm10,\pm30\}$ reduce to $\{10,7,13,4\}$ modulo $17$, the odd first
coordinates $\{\pm33,\pm17\}$ reduce to $\{16,1,0\}$, no residue matches, and
both side sets are empty: the $t'=17$ marginal fires, reproducing the
exhaustive-MITM kill of $q=1189$ by a quantified criterion.  The same
criterion with the cubic period basis fires at $q=1105$, $\ell=79$.  SAT
controls exist throughout the lane, both with rational-only halves matched
modulo $\ell$ (for example $q=985$, $\ell=17$) and with genuinely irrational
period witnesses (for example $q=441$, $\ell=13$, where
$441=(3\sqrt{13})^2+18^2$ enters the join), so the criterion is a firing
predicate, not an automatic orbit-shape obstruction.

\begin{proposition}[$C_{27}$ full-unit component criterion]\label{prop:c27-component}
Assume $27\mid t$, put $h=t/27$ and $q=54h-1$, and suppose the image $V'$ of the
composite multiplier group on $C_{27}$ is $(\Z/27\Z)^\times$.  Let the four
orbits be
\[
  O_0=\{0\},\quad O_1=(\Z/27\Z)^\times,\quad
  O_2=3(\Z/9\Z)^\times,\quad O_3=9(\Z/3\Z)^\times .
\]
For $X=x_0E_{O_0}+x_1E_{O_1}+x_2E_{O_2}+x_3E_{O_3}$ the rational components at
conductors $1,3,9,27$ are
\[
\begin{aligned}
  \lambda_1(X)&=x_0+18x_1+6x_2+2x_3,\\
  \lambda_3(X)&=x_0-9x_1+6x_2+2x_3,\\
  \lambda_9(X)&=x_0-3x_2+2x_3,\\
  \lambda_{27}(X)&=x_0-x_3 .
\end{aligned}
\]
Hence a $C_{27}$ marginal solution exists if and only if there are four
Gaussian representations
\[
  r_d^2+s_d^2=q\qquad(d=1,3,9,27)
\]
such that applying the inverse transform
\[
\begin{aligned}
  x_0&=(u_1+2u_3+6u_9+18u_{27})/27,\\
  x_1&=(u_1-u_3)/27,\\
  x_2&=(u_1+2u_3-3u_9)/27,\\
  x_3&=(u_1+2u_3+6u_9-9u_{27})/27
\end{aligned}
\]
to $(u_1,u_3,u_9,u_{27})=(r_1,r_3,r_9,r_{27})$ gives the $A$-orbit values and
applying it to $(s_1,s_3,s_9,s_{27})$ gives the $B$-orbit values, with both
inverse vectors integral and satisfying the fiber boxes and parities
$x_0(A)$ even with $|x_0(A)|\le h-1$, all other $A$-values odd of absolute value
at most $h$, and all $B$-values odd of absolute value at most $h$.  If this
Gaussian component join is empty, then no Turyn pair exists at $q$.
\end{proposition}

\begin{proof}
For $V'=(\Z/27\Z)^\times$, the fixed algebra has the four rational conductor
components displayed above; the period sums are the usual Ramanujan sums for
moduli $1,3,9,27$.  The displayed inverse matrix is the inverse of this
component transform.  Proposition~\ref{prop:component-join} gives the
component equations $r_d^2+s_d^2=q$ and says that precisely the integral inverse
vectors satisfying the boxes/parities are marginal solutions.  The obstruction
statement is Proposition~\ref{prop:t2-marginal}.
\end{proof}

\begin{proposition}[$C_{27}$ nested congruence criterion]\label{prop:c27-nested}
In the setting of Proposition~\ref{prop:c27-component}, let
\[
  G(q)=\{z=a+bi\in\Z[i]: a^2+b^2=q\},
\]
and let $G_{eo}(q)$ (respectively $G_{oe}(q)$) denote the elements of $G(q)$
with even real part and odd imaginary part (respectively odd real part and even
imaginary part).  A $C_{27}$ full-unit marginal solution exists if and only if
there are
\[
  z_1\in G_{eo}(q),\qquad z_3,z_9,z_{27}\in G_{oe}(q)
\]
such that, in $\Z[i]$,
\[
  z_1\equiv z_3\pmod{27},\qquad
  z_1\equiv z_9\pmod 9,\qquad
  z_1\equiv z_{27}\pmod 3,
\]
and the inverse transform
\[
\begin{aligned}
  w_0&=(z_1+2z_3+6z_9+18z_{27})/27,\\
  w_1&=(z_1-z_3)/27,\\
  w_2&=(z_1+2z_3-3z_9)/27,\\
  w_3&=(z_1+2z_3+6z_9-9z_{27})/27
\end{aligned}
\]
has real parts satisfying the $A$-boxes and imaginary parts satisfying the
$B$-boxes.  Equivalently, the $C_{27}$ marginal fires precisely when this
oriented nested congruence join is empty after the remaining box cut.
\end{proposition}

\begin{proof}
Write $z_d=r_d+is_d$.  Proposition~\ref{prop:c27-component} says that the
component criterion is the inverse transform applied separately to
$(r_1,r_3,r_9,r_{27})$ and $(s_1,s_3,s_9,s_{27})$.  Applying the inverse
transform in $\Z[i]$ packages these two real inversions into the displayed
$w_j$.

The inverse is integral if and only if the displayed nested congruences hold.
Indeed, $w_1\in\Z[i]$ is equivalent to $z_1\equiv z_3\pmod{27}$; then
$w_2\in\Z[i]$ is equivalent to $z_1\equiv z_9\pmod9$; with those two congruences,
$w_3\in\Z[i]$ is equivalent to $z_1\equiv z_{27}\pmod3$.  The formula for $w_0$
is then automatically integral.  The orientation condition is exactly the
fiber parity condition: $z_1\in G_{eo}$ and $z_3,z_9,z_{27}\in G_{oe}$ make
$\Re(w_0)$ even, the other real parts odd, and all imaginary parts odd; the
converse follows from the same parity calculation.  The remaining conditions
are precisely the fiber boxes of Proposition~\ref{prop:c27-component}.
\end{proof}

\begin{corollary}[Automatic boxes for large $C_{27}$ marginals]\label{cor:c27-auto-box}
In Proposition~\ref{prop:c27-nested}, if $h\ge57$, then the box inequalities are
automatic.  Hence for $h\ge57$ the $C_{27}$ full-unit marginal exists if and
only if the oriented nested congruence join is nonempty.
\end{corollary}

\begin{proof}
For every $z_d\in G(q)$ we have $|z_d|=\sqrt q=\sqrt{54h-1}$.  The displayed
inverse transform gives
\[
  |w_0|\le\sqrt q,\quad |w_1|\le\frac{2}{27}\sqrt q,\quad
  |w_2|\le\frac{2}{9}\sqrt q,\quad |w_3|\le\frac{2}{3}\sqrt q .
\]
When $h\ge57$, $\sqrt{54h-1}<h$, so every coordinate of each $w_j$ has absolute
value at most $h$.  In the non-vacuous case $q$ is a sum of two squares, so
$q\equiv1\pmod4$ and hence $h$ is odd; the even coordinate $\Re(w_0)$ is then
strictly smaller than $h$ in absolute value, so $|\Re(w_0)|\le h-1$.  These are
exactly the $A$- and $B$-boxes.
\end{proof}

\begin{corollary}[$C_{27}$ residue-pair firing criterion]\label{cor:c27-residue-pair}
Let $\overline G_{eo}(q)$ and $\overline G_{oe}(q)$ be the images of
$G_{eo}(q)$ and $G_{oe}(q)$ in $\Z[i]/27\Z[i]$.  If
\[
  \overline G_{eo}(q)\cap\overline G_{oe}(q)=\varnothing,
\]
then the $C_{27}$ full-unit marginal is impossible, hence no Turyn pair exists
whenever this marginal is forced by the multiplier group.  Conversely, if
$h\ge57$, then the $C_{27}$ full-unit marginal exists if and only if this
intersection is nonempty.
\end{corollary}

\begin{proof}
Any oriented nested congruence tuple in Proposition~\ref{prop:c27-nested}
satisfies $z_1\equiv z_3\pmod{27}$ with $z_1\in G_{eo}(q)$ and
$z_3\in G_{oe}(q)$, so it gives a residue in the displayed intersection.  Thus
an empty intersection kills the marginal before the box cut.  Conversely, if
$z\in G_{eo}(q)$ and $w\in G_{oe}(q)$ have the same residue modulo $27$, set
$z_1=z$ and $z_3=z_9=z_{27}=w$.  This satisfies all nested congruences and
orientations.  For $h\ge57$, Corollary~\ref{cor:c27-auto-box} makes the box
conditions automatic.
\end{proof}

\begin{corollary}[$C_{27}$ arithmetic firing alternatives]\label{cor:c27-arithmetic-firing}
Assume the hypotheses of Proposition~\ref{prop:c27-component}.  Let
$\mathcal J_{27}(q)$ be the set of oriented nested congruence tuples of
Proposition~\ref{prop:c27-nested}, and let $\mathcal B_{27}(q,h)\subset
\mathcal J_{27}(q)$ be the subset whose inverse orbit values satisfy the
fiber boxes.  The $C_{27}$ marginal is impossible, and hence no Turyn pair
exists, in either of the following quantified situations:
\[
  \overline G_{eo}(q)\cap\overline G_{oe}(q)=\varnothing,
  \qquad\text{or}\qquad
  h<57\ \text{ and }\ \mathcal B_{27}(q,h)=\varnothing .
\]
For $h\ge57$, the first condition is also necessary for firing within this
$C_{27}$ lane: the marginal exists if and only if the displayed residue-pair
intersection is nonempty.
\end{corollary}

\begin{proof}
The empty residue-pair condition is Corollary~\ref{cor:c27-residue-pair}.  If
$h<57$ and $\mathcal B_{27}(q,h)$ is empty, then
Proposition~\ref{prop:c27-nested} has no boxed inverse tuple and therefore no
marginal solution.  The large-$h$ equivalence is exactly
Corollary~\ref{cor:c27-auto-box} applied to
Corollary~\ref{cor:c27-residue-pair}.
\end{proof}

This is the first obstruction in the paper that sees genuinely composite
quotients between the prime quotients and the full torus. It closes cases where
each prime quotient passes but the data cannot coexist on a larger sub-torus;
for example $q=549$ has lifting prime quotients but is killed at the
prime-power marginal $t'=25$.

The $C_{27}$ criteria are the $(p,e)=(3,3)$ instance of one general
statement: at a prime-power node with \emph{full-unit} multiplier image,
all conductor components are rational, and the marginal is a join of
$e{+}1$ Gaussian representations glued by nested congruences.

\begin{proposition}[Full-unit prime-power nested criterion]
\label{prop:pe-full-unit-nested}
Let $p$ be an odd prime, $e\ge1$, $t'=p^e$, and suppose the multiplier
image at $t'$ is the full unit group $V'=(\Z/p^e)^\times$, with all
signs $+1$ and $q=2p^eh-1$, $h$ odd.  The $V'$-orbits are $\{0\}$ and
$O_j=p^j\,(\Z/p^{e-j})^\times$ for $0\le j\le e-1$; write $y_j$ for the
value on $O_j$ and $y_e$ for the origin value.  For $0\le m\le e$, the
conductor-$p^m$ component is rational, and
\[
  \lambda_{p^m}
  =y_e+\sum_{j\ge m}\varphi(p^{e-j})\,y_j-p^{e-m}\,y_{m-1}
  \qquad(\text{the last term present only for }m\ge1).
\]
The marginal system at $t'$ is solvable if and only if there are
Gaussian integers $z_1,z_p,\dots,z_{p^e}$ with $z_d\bar z_d=q$
satisfying the nested congruences
\[
  z_1\equiv z_{p^m}\pmod{p^{\,e+1-m}}\qquad(1\le m\le e),
\]
the orientation $z_1\in G_{eo}(q)$, $z_{p^m}\in G_{oe}(q)$, and whose
inverse orbit values (real parts on the $A$ side, imaginary parts on
the $B$ side) satisfy the fiber boxes.  The inverse is automatically
integral under the congruences, and the orientation is equivalent to
the fiber parities.
\end{proposition}

\begin{proof}
For a character $\chi$ of conductor $p^m$ and $u$ ranging over
$(\Z/p^e)^\times$, the sum of $\chi(p^ju)$-values is
$\varphi(p^{e-j})/\varphi(p^{m-j})\cdot\mu(p^{m-j})$ for $m\ge j$ and
$\varphi(p^{e-j})$ for $m\le j$: nonzero only for $m-j\le1$, giving the
displayed evaluation.  Taking successive differences,
$\lambda_{p^m}-\lambda_{p^{m+1}}=p^{e-m}(y_m-y_{m-1})$ for
$0\le m\le e-1$ (with $y_{-1}:=0$), so integrality of all orbit
values is equivalent to the adjacent congruences
$z_{p^m}\equiv z_{p^{m+1}}\pmod{p^{e-m}}$, which by transitivity along
the chain of moduli are equivalent to the displayed anchored form.  For
the parities: every $\varphi(p^{e-j})$ is even and every $p^{e-m}$ is
odd, so $\lambda_1\equiv y_e$ and $\lambda_{p^m}\equiv y_e+y_{m-1}
\pmod2$; on the $A$ side ($y_e$ even, other $y_j$ odd) this is the
orientation $z_1\in G_{eo}$, $z_{p^m}\in G_{oe}$, and on the $B$ side
(all odd) the complementary parities, exactly as in
Proposition~\ref{prop:c27-nested}.  Feasibility is then
Proposition~\ref{prop:component-join} for the rational conductor
decomposition.
\end{proof}

\begin{corollary}[Residue-pair collapse at $p^e$]\label{cor:pe-residue-pair}
If the residue sets of $G_{eo}(q)$ and $G_{oe}(q)$ modulo $p^e$ in
$\Z[i]$ are disjoint, the full-unit $p^e$ marginal fires.  Conversely,
for $h\ge2p^e+3$ the boxes are automatic from $|z_d|=\sqrt q\le h-1$,
and a common residue gives a marginal witness by taking
$z_p=\dots=z_{p^e}$ in the matching $G_{oe}$ class: the oriented
residue join alone decides the node.
\end{corollary}

\begin{proof}
A common residue is forced by $z_1\equiv z_p\pmod{p^e}$, giving the
firing direction.  For the converse, the binding box is the $A$-side
origin: $q=2p^eh-1\le(h-1)^2$ holds iff $h^2-(2p^e+2)h+2\ge0$, true
for $h\ge2p^e+3$ (odd $h$); the remaining coordinates are bounded by
$2\sqrt q\cdot p^{\,j-e}\le h$ already for small $h$.  The construction
$z_p=\dots=z_{p^e}$ satisfies all adjacent congruences.
\end{proof}

\begin{corollary}[$C_{49}$ full-unit nested criterion]\label{cor:c49-nested}
For $(p,e)=(7,2)$: with $q=98h-1$ and full-unit image at $t'=49$, the
marginal is solvable iff there are $z_1,z_7,z_{49}$ of norm $q$ with
$z_1\equiv z_7\pmod{49}$, $z_1\equiv z_{49}\pmod7$, the orientation
$z_1\in G_{eo}(q)$, $z_7,z_{49}\in G_{oe}(q)$, and boxed inverse
\[
  x_0=\tfrac1{49}(\lambda_1+6\lambda_7+42\lambda_{49}),\quad
  x_1=\tfrac1{49}(\lambda_1-\lambda_7),\quad
  x_2=\tfrac1{49}(\lambda_1+6\lambda_7-7\lambda_{49}).
\]
The residue-pair collapse (Corollary~\ref{cor:pe-residue-pair}) holds
with modulus $49$ and automatic boxes for $h\ge101$.
\end{corollary}

The $(7,2)$ instance closes the two remaining prime-power secondary
orders: $q=685$ ($h=7$) and $q=1665$ ($h=17$) both fire by disjoint
residue pairs modulo $49$, each confirmed by the exact orbit-value MITM
($448\times512$ and $5{,}508\times5{,}832$ assignments); the control
sweep over $q=98k-1$, $k$ odd, $k\le301$ finds $22$ shape-matching
orders of which $21$ fire and $q=23225$ ($h=237$) is a genuine SAT
control in the automatic-box regime.  The transform and inverse were
machine-derived by exact linear algebra and round-trip verified against
the constructed orbit algebras at both orders.  Note the full-unit
shape does not collide with the $p=7$ lift lane of
Proposition~\ref{prop:p7-cubic}: $q=1469$ at $t'=49$ has $|V'|=14$, not
$42$.

\begin{theorem}[Transitive two-orbit marginal obstruction]\label{thm:t2-twoorbit}
When $V'$ is transitive on the nonzero residues of $C_m$ for a prime
$m=t'$, the marginal algebra has the two orbits $\{0\}$ and $C_m\setminus\{0\}$.
Writing
$A=a_0+c\sum_{j\ne0}y^j$ and $B=b_0+d\sum_{j\ne0}y^j$, the nonzero PAF
constraint is
\[
  2a_0c+(m-2)c^2+2b_0d+(m-2)d^2=0,
\]
and the counting equation is
\[
  a_0^2+(m-1)c^2+b_0^2+(m-1)d^2=q.
\]
If these two equations have no solution with $a_0$ even, $|a_0|\le h-1$, and
$c,b_0,d$ odd with absolute value at most $h$, then no Turyn pair of order $t$
exists.
\end{theorem}

\begin{proof}
Transitivity gives $A=a_0+c\sum_{j\ne0}y^j$ and
$B=b_0+d\sum_{j\ne0}y^j$. For any nonzero shift, the PAF of the $A$ row has two
terms involving the origin and $m-2$ off-origin/off-origin terms, hence
$2a_0c+(m-2)c^2$; the $B$ row is identical with $(b_0,d)$. The zero-shift
coefficient is the displayed sum of squares. Lemma~\ref{lem:marginal-algebra}
supplies the boxes and parities, and Proposition~\ref{prop:t2-marginal} gives
the obstruction.
\end{proof}

Seven of the $q\le2000$ T2 closures below are of this transitive type; they are
certified by the two scalar equations above rather than by a general-purpose
solver.

The scalar system is small enough to admit a complete arithmetic
characterization: feasibility is equivalent to the existence of a pair of
integer representations of $q$ with opposite parity orientations, congruent
coordinatewise modulo $m$.  This converts the \texttt{scalar\_unsat}
mechanism of the branch accounting from an enumeration outcome into
closed-form arithmetic in $(q,M(q))$, and it isolates exactly which
two-orbit nodes can never fire.

\begin{theorem}[Two-orbit representation criterion]\label{thm:two-orbit-reps}
Assume the hypotheses of Theorem~\ref{thm:t2-twoorbit}: $m=t'$ is prime,
$V'$ is transitive on the nonzero residues, $h=t/m$, and $q=2mh-1$.  The
scalar system of Theorem~\ref{thm:t2-twoorbit} has a solution satisfying the
boxes and parities if and only if there are integers $X,Y,X',Y'$ with
\[
  q=X^2+Y^2=X'^2+Y'^2,\qquad
  X,\,Y'\ \text{even},\quad Y,\,X'\ \text{odd},
\]
\[
  X\equiv X'\pmod m,\qquad Y\equiv Y'\pmod m,
\]
whose inverse values
\[
  a_0=\frac{X+(m-1)X'}{m},\quad
  c=\frac{X-X'}{m},\quad
  b_0=\frac{Y+(m-1)Y'}{m},\quad
  d=\frac{Y-Y'}{m}
\]
satisfy $|a_0|\le h-1$ and $|c|,|b_0|,|d|\le h$.  The parities of
$(a_0,c,b_0,d)$ are automatic.  Moreover, if $h\ge2m+2$, the four box
conditions are automatic, so the marginal fires if and only if no congruent,
oppositely oriented pair of representations exists.
\end{theorem}

\begin{proof}
For prime $m$ with transitive $V'$, the fixed algebra is
$\Q[C_m]^{V'}\simeq\Q\times\Q$, with component maps
$\lambda_1(A)=a_0+(m-1)c$ (trivial character) and
$\lambda_m(A)=a_0-c$ (any nontrivial character sends the nonzero-orbit sum
to $\mu(m)=-1$), and likewise for $B$.  By
Proposition~\ref{prop:component-join}, the scalar system with its boxes and
parities is equivalent to the two component equations
$\lambda_\kappa(A)^2+\lambda_\kappa(B)^2=q$ together with integrality of the
inverse transform, the parities, and the boxes.  Set
$(X,Y)=(\lambda_1(A),\lambda_1(B))$ and
$(X',Y')=(\lambda_m(A),\lambda_m(B))$.  Since $a_0$ is even and $c,b_0,d$
are odd while $m-1$ is even, $X$ and $Y'$ are even and $Y$ and $X'$ are odd.
The displayed inverse is the linear inversion of the component transform,
and its integrality is exactly the two congruences.  Conversely, given the
congruences and the orientation, the parities of the inverse are automatic
because $m$ is odd: $ma_0=X+(m-1)X'$ is even and $mc=X-X'$ is odd, so $a_0$
is even and $c$ odd, and likewise $b_0,d$ are odd.

For the automatic boxes with $h\ge2m+2$: each component obeys
$|\,\cdot\,|\le\sqrt q$ by the two-square equations, and $a_0$ is the convex
combination $\frac1mX+\frac{m-1}mX'$, so $|a_0|\le\sqrt q$; the inequality
$\sqrt q\le h-1$ is $(h-1)^2-q=h^2-2(m+1)h+2\ge0$, which holds for
$h\ge2m+2$.  Similarly $|b_0|\le\sqrt q\le h$, and
$|c|,|d|\le2\sqrt q/m\le2(h-1)/m\le h$ for $m\ge3$.
\end{proof}

\begin{corollary}[Parity firing for large nodes]\label{cor:two-orbit-parity}
In the setting of Theorem~\ref{thm:two-orbit-reps}, if $m^2>4q$ ---
equivalently, using $q=2mh-1$, whenever $m\ge8h$ --- then no congruent pair
exists and the two-orbit marginal fires unconditionally: any
$X\equiv X'\pmod m$ with $|X|,|X'|\le\sqrt q<m/2$ forces $X=X'$,
contradicting the opposite parities.  No enumeration is involved.
\end{corollary}

\begin{corollary}[Transitive full torus]\label{cor:transitive-full-torus}
Let $t\ge5$ be prime with $V$ transitive on $(\Z/t)^\times$, all multiplier
signs $+1$.  Then no Turyn pair of order $t$ exists.  Indeed at the full
torus $h=1$, so $a_0=0$ and $|\lambda_1(A)|=(t-1)|c|=t-1$, while
$\lambda_1(A)^2+\lambda_1(B)^2=q$ forces $(t-1)^2\le2t-1$, false for
$t\ge5$.  (For $t\ge11$ this is also the parity case $t^2>4q=8t-4$ of
Corollary~\ref{cor:two-orbit-parity}.)
\end{corollary}

\begin{corollary}[The $C_3$ node never fires]\label{cor:c3-feasible}
Let $3\mid t$, $h=t/3\ge8$, and let $q=6h-1$ be a sum of two squares.  Then
the $C_3$ marginal is feasible.  Consequently the guaranteed node $t'=3$ of
Theorem~\ref{thm:divisor-selection}(iii) is never a certifying node for
$q\ge47$; orders with $3\mid t$ must fire elsewhere, for instance at the
$h=3$ node $t'=t/3$, as at $q=425$, $1445$, $1625$.
\end{corollary}

\begin{proof}
Here $q\equiv2\pmod3$, so any representation $q=e^2+o^2$ ($e$ even, $o$
odd) has $e^2\equiv o^2\equiv1\pmod3$, hence $3\nmid e,o$.  Use the same
representation on both components: choose signs so that
$X=\pm e\equiv X'=\pm o\pmod 3$ and $Y=\pm o\equiv Y'=\pm e\pmod 3$, which
is possible because $e,o\equiv\pm1\pmod3$ and the four signs are
independent.  The boxes are automatic since $h\ge8=2\cdot3+2$, so
Theorem~\ref{thm:two-orbit-reps} gives feasibility.  (The excluded cases
$h\in\{3,5,7\}$ have $q\in\{17,29,41\}$, all prime.)
\end{proof}

\begin{remark}[Arithmetic form, and the audit]\label{rem:two-orbit-arithmetic}
In $\Z[i]$ the criterion reads: with $D(q)=\{z:z\bar z=q\}$, the node
$(m,h)$ in the automatic-box regime is feasible if and only if
$z_1\equiv i\bar z_2\pmod{m\Z[i]}$ for some $z_1,z_2\in D(q)$ in the
even-real orientation --- a finite condition on the residues of the Gaussian
divisors of $q$ modulo $m$, computable from the factorization of $q$ in
$\Z[i]$.  It yields closed forms per $m$: for $m=3$ the condition is always
satisfiable (Corollary~\ref{cor:c3-feasible}); for $m=5$, since
$q\equiv4\pmod5$ forces exactly one of $e,o$ to be divisible by $5$ in each
representation, the node is feasible if and only if representations of both
types occur.  The criterion is machine-verified against the exact scalar
enumeration on every two-orbit node of every panel order
(\texttt{--two-orbit-rep-audit}): at $q\le2000$, $123$ nodes split into
$64$ feasible, $24$ parity fires (Corollary~\ref{cor:two-orbit-parity};
$15$ of them full-torus), $32$ residue fires, and $3$ box cuts ($65@11$,
$265@19$, $493@19$, each with $h<2m+2$); at $q\le10000$, $652$ nodes split
$346$/$128$/$169$/$9$ the same way, with zero disagreements.  The feasible
nodes are exactly the scalar local-pass rows: the complement that the
structural dichotomy must route to other divisors.
\end{remark}

\begin{definition}[Prime-square lift marginal]\label{def:prime-square-lift}
Let $p$ be an odd prime and suppose $t'=p^2$. A marginal is called a
\emph{prime-square $\{\pm1\}$-lift marginal} if the multiplier image is
\[
  V'=\{u\in(\Z/p^2)^\times:\ u\equiv \pm1 \pmod p\}.
\]
For $1\le i\le(p-1)/2$, put
\[
  U_i=\{a\in\Z/p^2:\ p\nmid a,\ a\equiv\pm i\pmod p\},\qquad
  P_i=\{\pm ip\}.
\]
The $V'$-orbits are $\{0\}$, the unit orbits $U_i$ of size $2p$, and the
$p$-multiple orbits $P_i$ of size $2$.
\end{definition}

\begin{theorem}[Prime-square lift marginal certificate]\label{thm:prime-square-lift}
Assume $q=2t-1$, $p^2\mid t$, $h=t/p^2$, and the marginal at $t'=p^2$ is a
prime-square $\{\pm1\}$-lift marginal. Let $E_0,E^U_i,E^P_i$ be the orbit-sum
basis from Definition~\ref{def:prime-square-lift}. For an orbit-value vector
$x$, let
\[
  X=x_0E_0+\sum_i x^U_iE^U_i+\sum_i x^P_iE^P_i
\]
and define
\[
  \Phi_p(x)=\bigl((XX^*)_O\bigr)_O\oplus X(1)^2,
\]
where $O$ ranges over the $p$ orbit classes. If a Turyn pair exists, then
there are orbit-value vectors $a,b$ such that
\[
  \Phi_p(a)+\Phi_p(b)=(q,0,\ldots,0,q),
\]
with $a_0$ even and $|a_0|\le h-1$, and all other entries odd with absolute
value at most $h$. Consequently, if this finite quadratic system has no
solution, then no Turyn pair exists at $q$.
\end{theorem}

\begin{proof}
Definition~\ref{def:prime-square-lift} gives the orbit-sum basis of
$\mathcal A(p^2,V')$. Lemma~\ref{lem:marginal-algebra} supplies the equation
$AA^*+BB^*=q\cdot1$, the trivial-character equation, and the same fiber
boxes/parities. These are exactly the displayed finite system.
\end{proof}

For $p=3$, this theorem is already a three-orbit scalar system. With
$x=(x_0,x_U,x_P)$,
\[
  \Phi_3(x)=
  \begin{pmatrix}
  x_0^2+6x_U^2+2x_P^2\\
  2x_0x_U+3x_U^2+4x_Ux_P\\
  2x_0x_P+6x_U^2+x_P^2\\
  (x_0+6x_U+2x_P)^2
  \end{pmatrix}.
\]
The orbit pattern alone is not an obstruction: some $p=3$ lift marginals have
solutions. The theorem is therefore a reusable exact certificate criterion; the
next closed-form step is to find arithmetic hypotheses on $q,h,p$ forcing the
join above to be empty.

\begin{proposition}[The $p=3$ Gaussian component criterion]\label{prop:p3-gaussian}
In the prime-square lift marginal with $t'=9$, write an orbit-value vector as
$x=(x_0,x_U,x_P)$, where $x_U$ is the value on the units and $x_P$ is the value
on $\{3,6\}$. Put
\[
  \lambda_1(x)=x_0+6x_U+2x_P,\quad
  \lambda_3(x)=x_0-3x_U+2x_P,\quad
  \lambda_9(x)=x_0-x_P .
\]
Then a $V'$-invariant marginal solution exists if and only if there are three
ordered representations
\[
  X_i^2+Y_i^2=q\qquad (i=1,3,9)
\]
such that
\[
\begin{aligned}
a_0&=\frac{X_1+2X_3+6X_9}{9},&
a_U&=\frac{X_1-X_3}{9},&
a_P&=\frac{X_1+2X_3-3X_9}{9},\\
b_0&=\frac{Y_1+2Y_3+6Y_9}{9},&
b_U&=\frac{Y_1-Y_3}{9},&
b_P&=\frac{Y_1+2Y_3-3Y_9}{9}
\end{aligned}
\]
are integral and satisfy the fiber boxes/parities. In particular, absence of
such a compatible triple of Gaussian representations is a field-free firing
criterion for the $p=3$ lift marginal.
\end{proposition}

\begin{proof}
The fixed algebra $\Q[C_9]^{(\Z/9)^\times}$ has three rational character
components, represented by the trivial character, a character of conductor $3$,
and a primitive character of conductor $9$. Evaluating
$x_0E_0+x_UE_U+x_PE_P$ at these three components gives the displayed
$\lambda_1,\lambda_3,\lambda_9$. Hence $AA^*+BB^*=q\cdot1$ is equivalent to
$\lambda_i(A)^2+\lambda_i(B)^2=q$ for $i=1,3,9$. The displayed inverse is the
inverse linear transform from component values back to orbit values. The boxes
and parities are those of Lemma~\ref{lem:marginal-algebra}.
\end{proof}

This criterion reproduces the $p=3$ prime-square lift split in the
$q\le2000$ blind-spot panel: it is infeasible for
$q=485,1385,1421,1745$, while the known SAT controls admit compatible triples.

\begin{proposition}[The $p=5$ quadratic-component criterion]\label{prop:p5-sqrt5}
In the prime-square lift marginal with $t'=25$, write an orbit-value vector as
$x=(x_0,x_1,x_2,x_3,x_4)$, where $x_1,x_2$ are the two unit-orbit values
modulo $5$ and $x_3,x_4$ are the values on $\{\pm5\}$ and $\{\pm10\}$.
Let $u$ satisfy $u^2+u-1=0$. Put
\[
\begin{aligned}
\lambda_1(x)&=x_0+10x_1+10x_2+2x_3+2x_4,\\
\lambda_5(x)&=(x_0+2x_3+2x_4-5x_2)+5(x_1-x_2)u,\\
\lambda_{25}(x)&=(x_0-x_4)+(x_3-x_4)u .
\end{aligned}
\]
Then a $V'$-invariant marginal solution exists if and only if there are
representations
\[
  X_1^2+Y_1^2=q,\qquad
  \alpha_5^2+\beta_5^2=q,\qquad
  \alpha_{25}^2+\beta_{25}^2=q
\]
with $\alpha_5,\beta_5,\alpha_{25},\beta_{25}\in\Z[u]$, such that the inverse
transform of $(X_1,\alpha_5,\alpha_{25})$ and
$(Y_1,\beta_5,\beta_{25})$ gives integral orbit values satisfying the fiber
boxes/parities. Explicitly, if $\lambda_5=A+Bu$ and $\lambda_{25}=C+Du$, set
\[
  \Delta_U=\frac{B}{5},\qquad \Delta_P=\frac{A-C-2D}{5},\qquad
  x_4=\frac{\lambda_1-C-2D+20\Delta_P-10\Delta_U}{25},
\]
and then
\[
  x_3=D+x_4,\qquad x_0=C+x_4,\qquad
  x_2=x_4-\Delta_P,\qquad x_1=x_2+\Delta_U .
\]
The divisibility and box/parity conditions in this inverse transform are
therefore an exact field-free firing criterion for the $p=5$ lift marginal.
\end{proposition}

\begin{proof}
The fixed algebra decomposes into the trivial component and two quadratic
components isomorphic to $\Q(u)$: the conductor-$5$ component and the
conductor-$25$ component. Evaluating the five orbit sums in these three
components gives the displayed $\lambda_1,\lambda_5,\lambda_{25}$. Since the
involution is trivial on the fixed quadratic components, the marginal equation
is equivalent to the three displayed two-square equations. The inverse formulas
solve the displayed linear transform. Finally, Lemma~\ref{lem:marginal-algebra}
supplies the fiber boxes and parity restrictions.
\end{proof}

The criterion reproduces the exact MITM obstructions at $q=549$ and $q=949$,
and it also closes the formerly unresolved $p=5$ stress instance
$q=1649,t'=25$. The next member, $p=7$, is cubic rather than quadratic but
has the same component shape.

\begin{proposition}[The $p=7$ cubic-component criterion]\label{prop:p7-cubic}
In the prime-square lift marginal with $t'=49$, write an orbit-value vector as
\[
  x=(x_0,x^U_1,x^U_2,x^U_3,x^P_1,x^P_2,x^P_3),
\]
where $x^U_i$ is the value on residues congruent to $\pm i\bmod 7$ and
$x^P_i$ is the value on $\{\pm7i\}$.  Let $u$ satisfy
$u^3+u^2-2u-1=0$.  Put
\[
\begin{aligned}
\lambda_1(x)
  &=x_0+14(x^U_1+x^U_2+x^U_3)+2(x^P_1+x^P_2+x^P_3),\\
\lambda_7(x)
  &=(x_0+2(x^P_1+x^P_2+x^P_3)+7(-2x^U_2+x^U_3))\\
  &\quad +7(x^U_1-x^U_3)u+7(x^U_2-x^U_3)u^2,\\
\lambda_{49}(x)
  &=(x_0-2x^P_2+x^P_3)+(x^P_1-x^P_3)u+(x^P_2-x^P_3)u^2 .
\end{aligned}
\]
Then a $V'$-invariant marginal solution exists if and only if there are
representations
\[
  L^2+M^2=q,\qquad
  \alpha_7^2+\beta_7^2=q,\qquad
  \alpha_{49}^2+\beta_{49}^2=q
\]
with $\alpha_7,\beta_7,\alpha_{49},\beta_{49}\in\Z[u]$, such that the inverse
transform of $(L,\alpha_7,\alpha_{49})$ and $(M,\beta_7,\beta_{49})$ gives
integral orbit values satisfying the fiber boxes/parities.  Explicitly, if
\[
  \alpha_7=A+Bu+Cu^2,\qquad \alpha_{49}=D+Eu+Fu^2,
\]
set
\[
\begin{gathered}
  \Delta_1=B/7,\qquad \Delta_2=C/7,\\
  R=(A-D-2E-4F+2C)/7,\qquad
  S=(L-D-2E-4F-2B-2C)/7,
\end{gathered}
\]
then require these quantities and $x^U_3=(S-R)/7$ to be integral, and put
\[
\begin{aligned}
x^P_3&=R+x^U_3,&
x^P_1&=E+x^P_3,&
x^P_2&=F+x^P_3,\\
x_0&=D+2F+x^P_3,&
x^U_1&=x^U_3+\Delta_1,&
x^U_2&=x^U_3+\Delta_2 .
\end{aligned}
\]
The same inverse applied to $(M,\beta_7,\beta_{49})$ gives the $B$ row.  If
no component halves survive these inverse boxes and parities, or no surviving
halves have componentwise two-square mates, then the $C_{49}$ marginal is
impossible and no Turyn pair exists.
\end{proposition}

\begin{proof}
Let $s_i=\zeta_7^i+\zeta_7^{-i}$.  With $u=s_1$, one has
$s_2=u^2-2$ and $s_3=1-u-u^2$.  In the conductor-$7$ component, a unit orbit
contributes $7s_i$ and a $7$-multiple orbit contributes $2$; in the
conductor-$49$ component, the unit-orbit sums vanish and the $7$-multiple
orbits contribute $s_i$.  This gives the displayed
$\lambda_1,\lambda_7,\lambda_{49}$.  Since the fixed algebra is real, the
marginal equation is equivalent to the three displayed two-square component
equations.  Solving the displayed linear transform gives the inverse formulas,
and Lemma~\ref{lem:marginal-algebra} supplies the boxes and parities.
\end{proof}

\begin{proposition}[Closure of $q=1469$ at $t'=49$]\label{prop:q1469-p7-cubic}
Let $q=1469$, so $t=735=3\cdot5\cdot7^2$, and take the prime-square marginal
$t'=49$ with $h=15$.  The $p=7$ cubic-component criterion of
Proposition~\ref{prop:p7-cubic} has no valid side vector.  Hence no Turyn pair
exists at $q=1469$.
\end{proposition}

\begin{proof}
The multiplier image is
$V'=\{u\bmod49:\ u\equiv\pm1\bmod7\}$, so the marginal is exactly the
$p=7$ prime-square lift.  The component enumeration has $16$ ordered integer
representations of $1469$ and $256$ ordered representations
$\alpha^2+\beta^2=1469$ in $\Z[u]$.  Applying the inverse transform of
Proposition~\ref{prop:p7-cubic} to all component halves yields no admissible
$A$ side vector and no admissible $B$ side vector.  Thus the component join is
empty, and Proposition~\ref{prop:t2-marginal} applies.
\end{proof}

\begin{proposition}[$C_{21}$ quadratic-component criterion]\label{prop:c21-sqrt21}
Assume $21\mid t$, put $h=t/21$, and suppose the image of the composite
multiplier group on $C_{21}$ is
\[
  H=\{1,4,5,16,17,20\}\subset(\Z/21\Z)^\times .
\]
Let
\[
\begin{aligned}
O_0&=\{0\},&
O_1&=H,&
O_2&=2H=\{2,8,10,11,13,19\},\\
O_3&=3(\Z/7\Z)^\times=\{3,6,9,12,15,18\},&
O_4&=\{7,14\}.
\end{aligned}
\]
Write an orbit-value vector as $x=(x_0,x_1,x_2,x_3,x_4)$ in this order and set
$\eta=(1+\sqrt{21})/2$, so $\eta^2=\eta+5$. The conductor components are
\[
\begin{aligned}
  \lambda_1(x)&=x_0+6x_1+6x_2+6x_3+2x_4,\\
  \lambda_3(x)&=x_0-3x_1-3x_2+6x_3-x_4,\\
  \lambda_7(x)&=x_0-x_1-x_2-x_3+2x_4,\\
  \lambda_{21}(x)&=(x_0+x_2-x_3-x_4)+(x_1-x_2)\eta .
\end{aligned}
\]
Consequently a $C_{21}$ marginal solution exists if and only if there are
three Gaussian representations
\[
  L_d^2+M_d^2=q\qquad(d=1,3,7)
\]
and one representation
\[
  \alpha^2+\beta^2=q,\qquad \alpha,\beta\in\Z[\eta],
\]
such that, writing $\alpha=P+Q\eta$, the inverse transform
\[
\begin{aligned}
  x_0&=(L_1+2L_3+6L_7+12P+6Q)/21,\\
  x_1&=(L_1-L_3-L_7+P+11Q)/21,\\
  x_2&=(L_1-L_3-L_7+P-10Q)/21,\\
  x_3&=(L_1+2L_3-L_7-2P-Q)/21,\\
  x_4&=(L_1-L_3+6L_7-6P-3Q)/21
\end{aligned}
\]
gives integral $A$-orbit values satisfying the $A$ fiber boxes/parities, and
the same inverse applied to $(M_1,M_3,M_7,R,S)$, where $\beta=R+S\eta$, gives
integral $B$-orbit values satisfying the $B$ boxes/parities. If this component
join is empty, then no Turyn pair exists at $q$.
\end{proposition}

\begin{proof}
The conductor-$1$, $3$, and $7$ components are rational. The displayed
$\lambda_1,\lambda_3,\lambda_7$ are the corresponding Ramanujan sums on the
five $H$-orbits. For conductor $21$, the primitive residues split into $H$ and
its complementary coset. If
\[
  \eta=\sum_{u\in H}\zeta_{21}^u,
\]
then the complementary period is $1-\eta$, and
$\eta(1-\eta)=-5$; hence $\Z[\eta]=\Z[(1+\sqrt{21})/2]$ and
$\eta^2=\eta+5$. Evaluating the five orbit sums at a primitive $21$st root
therefore gives the displayed $\lambda_{21}$.

Since $-1\in H$, inversion is trivial on the fixed algebra. Thus
Proposition~\ref{prop:component-join} converts the marginal identity into the
three rational two-square equations and the one $\Z[\eta]$ two-square equation.
The displayed formulas are the inverse of the component transform. Integral
inverse values satisfying the fiber boxes/parities are exactly the marginal
solutions of Lemma~\ref{lem:marginal-algebra}; if no such join exists,
Proposition~\ref{prop:t2-marginal} gives the obstruction.
\end{proof}

\begin{corollary}[$C_{21}$ arithmetic firing criterion]\label{cor:c21-arithmetic-firing}
Assume the hypotheses of Proposition~\ref{prop:c21-sqrt21}.  Let
$Z(q)=\{(L,M)\in\Z^2:L^2+M^2=q\}$ and let
$K_{21}(q)=\{(\alpha,\beta)\in\Z[\eta]^2:\alpha^2+\beta^2=q\}$, with
$\eta=(1+\sqrt{21})/2$.  Define $\mathcal S_A(q,h)$ to be the set of
component halves $(L_1,L_3,L_7,\alpha)$ with each $L_i$ appearing as the first
coordinate of some element of $Z(q)$ and $\alpha$ appearing as the first
coordinate of some element of $K_{21}(q)$, such that the inverse transform of
Proposition~\ref{prop:c21-sqrt21} is integral and satisfies the $A$ boxes and
parities.  Define $\mathcal S_B(q,h)$ in the same way using the $B$ boxes.
If $\mathcal S_A(q,h)=\varnothing$, or $\mathcal S_B(q,h)=\varnothing$, or no
element $(L_1,L_3,L_7,\alpha)\in\mathcal S_A(q,h)$ has a mate
$(M_1,M_3,M_7,\beta)\in\mathcal S_B(q,h)$ with
\[
  (L_i,M_i)\in Z(q)\quad(i=1,3,7),
  \qquad
  (\alpha,\beta)\in K_{21}(q),
\]
then the $C_{21}$ marginal is impossible, and hence no Turyn pair exists.
\end{corollary}

\begin{proof}
Proposition~\ref{prop:c21-sqrt21} says that every marginal solution is exactly
such a matched pair of component halves whose inverse transforms satisfy the
fiber boxes and parities.  If one side set is empty, or if the two side sets
have no componentwise two-square mate, the component join is empty.  The
obstruction is then Proposition~\ref{prop:t2-marginal}.
\end{proof}

This criterion reproduces the compact-middle $C_{21}$ obstructions at
$545@21$ and $629@21$. It also shows why the signature is a criterion rather
than an automatic obstruction: examples such as $2225@21$ admit boxed
component witnesses.

The same template fires at $t'=57$, where it retires the last recorded
exhaustive certificate of the finite panel.

\begin{proposition}[$C_{57}$ quadratic-component criterion]\label{prop:c57-sqrt57}
Let $q=114h-1$ with $h$ odd, and suppose the multiplier image at
$t'=57$ is $V'=\ker(\chi_3\chi_{19})$, the unique index-two subgroup of
$(\Z/57)^\times$ containing $-1$ that surjects onto both
$(\Z/3)^\times$ and $(\Z/19)^\times$; the orbits are
$\{0\},H,3\,(\Z/19)^\times,gH,\{19,38\}$ of sizes $1,18,18,18,2$.  The
fixed algebra is $\Q\times\Q\times\Q\times\Q(\sqrt{57})$ with
conductors $1,3,19,57$; the conductor-$57$ component takes values in
$\Z[\eta]$, $\eta=(1+\sqrt{57})/2$, $\eta^2=\eta+14$.  In orbit order
$(x_0,\dots,x_4)$ the transform is
\[
\begin{aligned}
  \lambda_1&=x_0+18x_1+18x_2+18x_3+2x_4,\\
  \lambda_3&=x_0-9x_1+18x_2-9x_3-x_4,\\
  \lambda_{19}&=x_0-x_1-x_2-x_3+2x_4,\\
  \lambda_{57}&=(x_0-x_2+x_3-x_4)+(x_1-x_3)\,\eta,
\end{aligned}
\]
with inverse of denominator $57$ (writing $\lambda_{57}=a+b\eta$)
\[
\begin{aligned}
  x_0&=\tfrac1{57}(\lambda_1+2\lambda_3+18\lambda_{19}+36a+18b),&
  x_1&=\tfrac1{57}(\lambda_1-\lambda_3-\lambda_{19}+a+29b),\\
  x_2&=\tfrac1{57}(\lambda_1+2\lambda_3-\lambda_{19}-2a-b),&
  x_3&=\tfrac1{57}(\lambda_1-\lambda_3-\lambda_{19}+a-28b),\\
  x_4&=\tfrac1{57}(\lambda_1-\lambda_3+18\lambda_{19}-18a-9b).&&
\end{aligned}
\]
The marginal at $t'=57$ is solvable if and only if there are three
Gaussian representations $z_1,z_3,z_{19}$ of $q$ and one representation
$\alpha^2+\beta^2=q$ in $\Z[\eta]$ --- equivalently, with
$\alpha=a+b\eta$, $\beta=c+d\eta$: $a^2+c^2+14(b^2+d^2)=q$ and
$2ab+2cd+b^2+d^2=0$, completely enumerable inside
$|b|,|d|\le\sqrt{q/14}$ --- whose inverse transform per side is
integral and lands in the fiber boxes, with parities
$(\lambda_1,\lambda_3,\lambda_{19},a,b)\equiv(0,1,1,1,0)$ on the $A$
side and $(1,0,0,0,0)$ on the $B$ side modulo $2$.
\end{proposition}

\begin{proof}
Identical to Proposition~\ref{prop:c21-sqrt21} with the conductor data
$(1,3,7,21)$, $\eta^2=\eta+5$ replaced by $(1,3,19,57)$,
$\eta^2=\eta+14$: the character evaluations give the displayed
transform ($\eta_H+\eta_{gH}=\mu(57)=1$,
$\eta_H-\eta_{gH}=\sqrt{57}$ by the even-character Gauss sum), the
parity equivalences are exact linear algebra modulo $2$, and
feasibility is Proposition~\ref{prop:component-join}.  The transform
and inverse were machine-derived and round-trip verified against the
constructed orbit algebra.
\end{proof}

At $q=1937$ ($h=17$) the criterion fires by an empty component join:
two valid side vectors on each side and no compatible pair, confirmed
by the exact orbit-value MITM ($1{,}784{,}592\times1{,}889{,}568$
assignments), matching the previously recorded exhaustive kill.  The
control sweep over $q=114k-1$, $k$ odd, $k\le301$ finds $48$
shape-matching orders: $42$ fire and six are SAT controls (first
$q=5585$ --- which carries both shapes: its $C_{49}$ node fires by
Corollary~\ref{cor:pe-residue-pair} while its $C_{57}$ node is
solvable).  With this promotion every closure in the $21$-order
T3-unreachable panel is a quantified firing route: no finite
certificate remains.

\begin{proposition}[The $p=5$ prime-square lift certificate at
\texorpdfstring{$q=549$}{q=549}]\label{prop:q549-fiveorbit}
Let $q=549$, so $t=275$, and take the prime-square marginal $t'=25$ with
$h=11$. The induced multiplier group $V'$ has five orbits
\[
\begin{aligned}
O_0&=\{0\},\\
O_1&=\{1,4,6,9,11,14,16,19,21,24\},\\
O_2&=\{2,3,7,8,12,13,17,18,22,23\},\\
O_3&=\{5,20\},\qquad O_4=\{10,15\}.
\end{aligned}
\]
Write $A=\sum_{i=0}^4 a_iE_i$ and $B=\sum_{i=0}^4 b_iE_i$. Put
\[
\begin{aligned}
F_0(x)&=x_0^2+10x_1^2+10x_2^2+2x_3^2+2x_4^2,\\
F_1(x)&=2x_0x_1+10x_1x_2+4x_1x_3+4x_1x_4+5x_2^2,\\
F_2(x)&=2x_0x_2+5x_1^2+10x_1x_2+4x_2x_3+4x_2x_4,\\
F_3(x)&=2x_0x_3+10x_1^2+10x_2^2+2x_3x_4+x_4^2,\\
F_4(x)&=2x_0x_4+10x_1^2+10x_2^2+x_3^2+2x_3x_4,
\end{aligned}
\]
and $L(x)=x_0+10x_1+10x_2+2x_3+2x_4$. Then the marginal system is
\[
  (F_0,\ldots,F_4,L^2)(a)+(F_0,\ldots,F_4,L^2)(b)
  =(549,0,0,0,0,549),
\]
with $a_0$ even, $|a_0|\le10$, and $a_i,b_i$ odd with absolute value at most
$11$ otherwise. This finite system has no solution. Hence no Turyn pair exists
at $q=549$.
\end{proposition}

\begin{proof}
The displayed forms are the coefficient equations of
Lemma~\ref{lem:marginal-algebra} in the five orbit-sum basis. The stated boxes
are the $h=11$ fiber boxes. The exact certificate is a side-separated
meet-in-the-middle over the integer orbit values: enumerate all
$11\cdot12^4=228096$ possible $a$-vectors and all $12^5=248832$ possible
$b$-vectors, storing the exact vector $(F_0,\ldots,F_4,L^2)$. No $a$-side vector
and $b$-side vector sum to $(549,0,0,0,0,549)$. By
Proposition~\ref{prop:t2-marginal}, this infeasible marginal rules out a Turyn
pair at $q=549$.
\end{proof}

This five-orbit certificate is the first prime-power marginal model beyond the
two-orbit theorem and is the $p=5$ instance of
Theorem~\ref{thm:prime-square-lift}. It shows the shape of the next structural
task: prime quotients can lift separately, while a prime-power marginal may
compress the problem to a finite algebraic join whose infeasibility must then be
proved by additional arithmetic, modular, energy, or box/parity input.

The same component strategy extends beyond prime-power sub-tori. The last open
case of the $q\le2000$ harvest, $q=1325$, is decided by the analogous criterion
at the composite marginal $t'=51$, where the two nonrational components live in
the real quartic subfield of $\Q(\zeta_{17})$.

\begin{proposition}[The quartic-component criterion at $t'=51$]\label{prop:c51-quartic}
Let $q=2t-1$ with $51\mid t$, $h=t/51$, and suppose the multiplier image in
$(\Z/51)^\times$ is the order-$8$ product group
\[
  V'=\{\pm1\bmod3\}\times\{\pm1,\pm4\bmod17\}
    =\{1,4,13,16,35,38,47,50\}.
\]
The $V'$-orbits on $\Z/51$ are $\{0\}$; four unit orbits $U_0,\dots,U_3$ of
size $8$; four $3$-multiple orbits $T_0,\dots,T_3$ of size $4$; and the
$17$-multiple orbit $\{17,34\}$ of size $2$. Here $k$ indexes the cosets
$3^k\langle4\rangle$ of $\langle4\rangle=\{1,4,13,16\}$ in $(\Z/17)^\times$:
the mod-$17$ image of $U_k$ is the coset $3^k\langle4\rangle$ with each residue
hit twice, and that of $T_k$ is the same coset with each residue hit once. Let
$\eta_k=\sum_{r\in3^k\langle4\rangle}\zeta_{17}^r$ be the quartic Gaussian
periods and $F\subset\Q(\zeta_{17})$ the real quartic subfield, so that
$\mathcal O_F=\Z\eta_0\oplus\cdots\oplus\Z\eta_3$ and $\Gal(F/\Q)$ shifts the
index $k$ cyclically. For an orbit-value vector
$x=(x_0;u_0,\dots,u_3;p_0,\dots,p_3;x_{17})$ put (indices modulo $4$)
\[
\begin{aligned}
\lambda_1(x)&=x_0+2x_{17}+8\sum\nolimits_ku_k+4\sum\nolimits_kp_k,\\
\lambda_3(x)&=x_0-x_{17}-4\sum\nolimits_ku_k+4\sum\nolimits_kp_k,\\
\lambda_{17}(x)&=(x_0+2x_{17})+\sum\nolimits_k(2u_k+p_k)\,\eta_k,\\
\lambda_{51}(x)&=(x_0-x_{17})+\sum\nolimits_k(p_k-u_k)\,\eta_{k+3}.
\end{aligned}
\]
Then the marginal equation is equivalent to the four two-square equations
\[
  \lambda_1(a)^2+\lambda_1(b)^2=q,\qquad
  \lambda_3(a)^2+\lambda_3(b)^2=q\qquad\text{in }\Z,
\]
\[
  \lambda_{17}(a)^2+\lambda_{17}(b)^2=q,\qquad
  \lambda_{51}(a)^2+\lambda_{51}(b)^2=q\qquad\text{in }\mathcal O_F .
\]
Every element occurring in the $\mathcal O_F$ equations has all four real
embeddings bounded by $\sqrt q$, so both $\mathcal O_F$ equations have finitely
many, exactly enumerable solutions. Writing $\lambda_{17}=\sum_ka_k\eta_k$ and
$\lambda_{51}=\sum_kb_k\eta_k$ in the period basis (via
$1=-\sum_k\eta_k$), the inverse transform is
\[
  P=\frac{\lambda_1-4\sum_ka_k}{17},\qquad
  C=\frac{\lambda_3-4\sum_kb_k}{17},\qquad
  x_{17}=\frac{P-C}{3},\qquad x_0=x_{17}+C,
\]
\[
  u_k=\frac{(a_k+P)-(b_{k+3}+C)}{3},\qquad
  p_k=u_k+b_{k+3}+C .
\]
A $V'$-invariant marginal solution exists if and only if some choice of the
four representations makes both inverse transforms integral and satisfy the
fiber boxes/parities of Lemma~\ref{lem:marginal-algebra}.
\end{proposition}

\begin{proof}
Since $V'$ is the full product of $\{\pm1\}\subset(\Z/3)^\times$ and
$\langle4\rangle\subset(\Z/17)^\times$, the fixed algebra $\Q[C_{51}]^{V'}$
decomposes by conductor as $\Q\times\Q\times F\times F$: the conductor-$3$
component is fixed by all of $(\Z/3)^\times$, hence rational; the
conductor-$17$ and conductor-$51$ components are fixed by $\langle4\rangle$,
hence land in $F$ (for conductor $51$, the $\zeta_3$-part of each character
sums to $\zeta_3+\zeta_3^2=-1$ over the $\{\pm1\bmod3\}$ factor). The
dimensions $1+1+4+4$ match the ten orbits, and the periods are an integral
basis of $\mathcal O_F$ permuted cyclically by $\Gal(F/\Q)$.

Evaluating orbit sums: under the conductor-$17$ character
$j\mapsto\zeta_{17}^{\,j\bmod17}$, the orbit $U_k$ evaluates to $2\eta_k$,
$T_k$ to $\eta_k$, and $\{17,34\}$ to $2$; under the conductor-$51$ character
$j\mapsto\zeta_{51}^{\,j}=\zeta_3^{2j}\zeta_{17}^{6j}$, the $\zeta_3$-factor
sums to $-1$ on $U_k$ and on $\{17,34\}$ and to $+1$ on $T_k$ (whose elements
are $0\bmod3$), while multiplication by $6\in3^3\langle4\rangle$ shifts the
coset index by $3$; hence $U_k\mapsto-\eta_{k+3}$, $T_k\mapsto\eta_{k+3}$,
$\{17,34\}\mapsto-1$. This gives the displayed evaluations. Since $-1\in V'$,
every component value is fixed by the involution, so the marginal equation
$AA^*+BB^*=q\cdot1$ is equivalent to the componentwise two-square equations;
the embedding bound follows because
$\sigma(\lambda(A))^2\le\sigma(\lambda(A))^2+\sigma(\lambda(B))^2=q$ at each
real place. For the inverse: with $P=x_0+2x_{17}$ and $C=x_0-x_{17}$, the
period coordinates are $a_k=(2u_k+p_k)-P$ and $b_{k+3}=(p_k-u_k)-C$, so
$\sum_ka_k=2\sum u_k+\sum p_k-4P$ and $\lambda_1-4\sum_ka_k=17P$; similarly
$\lambda_3-4\sum_kb_k=17C$; the per-index $2\times2$ blocks
$(a_k+P,\;b_{k+3}+C)=(2u_k+p_k,\;p_k-u_k)$ invert as displayed. The boxes and
parities are those of Lemma~\ref{lem:marginal-algebra}.
\end{proof}

\begin{proposition}[Closure of $q=1325$ at the $t'=51$ marginal]\label{prop:q1325-quartic}
Let $q=1325=5^2\cdot53$, $t=663=3\cdot13\cdot17$, and take the marginal
$t'=51$, $h=13$, whose multiplier image is exactly the $V'$ of
Proposition~\ref{prop:c51-quartic}. Then:
\emph{(i)} the ordered integer representations of $1325$ as a sum of two
squares are the $24$ pairs arising from
$35^2+10^2=34^2+13^2=29^2+22^2=1325$;
\emph{(ii)} the equation $\alpha^2+\beta^2=1325$ has exactly $48$ ordered
solutions in $\mathcal O_F$, all lying in the quadratic subfield
$\Q(\sqrt{17})\subset F$: the $24$ rational pairs together with the orbit of
\[
  1325=30^2+\bigl(5\sqrt{17}\bigr)^2
      =\bigl(24-3\sqrt{17}\bigr)^2+\bigl(18+4\sqrt{17}\bigr)^2
\]
under conjugation, negation, and swap;
\emph{(iii)} exactly two $A$-side and two $B$-side orbit-value vectors pass
the inverse transform with the $h=13$ boxes and parities, namely
\[
  \pm a^\dagger:\quad x_0=0,\ u_k=p_k=-1,\ x_{17}=13;\qquad
  \pm b^\dagger:\quad x_0=13,\ u_k=p_k=-1,\ x_{17}=3,
\]
with components
$(\lambda_1,\lambda_3,\lambda_{17},\lambda_{51})(a^\dagger)=(-22,-13,29,-13)$
and $(\lambda_1,\lambda_3,\lambda_{17},\lambda_{51})(b^\dagger)=(-29,10,22,10)$;
\emph{(iv)} no combination survives: for every sign choice,
\[
  \lambda_3(\pm a^\dagger)^2+\lambda_3(\pm b^\dagger)^2=13^2+10^2=269\ne1325 .
\]
Hence the $t'=51$ marginal system is infeasible, and by
Proposition~\ref{prop:t2-marginal} no Turyn pair exists at $q=1325$.
\end{proposition}

\begin{proof}
(i) is elementary. For (ii), every solution has all four real embeddings in
$[-\sqrt q,\sqrt q]$, hence lies in a finite lattice box in the Minkowski
embedding of $\mathcal O_F$ ($400{,}753$ integer vectors); an exact
meet-in-the-middle on the period coordinates of $\alpha^2$ over this box
produces exactly the $48$ stated pairs. (The enumeration was reproduced by a
structurally independent square-root--recovery method, and the solution set is
closed under the Galois action, negation, and swap, as it must be.) For (iii),
run the complete join of Proposition~\ref{prop:c51-quartic}: over the
$24\times24\times48\times48$ component choices, the divisibility conditions
($17\mid\lambda_1-4\sum a_k$, $17\mid\lambda_3-4\sum b_k$, the mod-$3$ block
conditions) and the $h=13$ boxes/parities leave exactly the displayed rows;
two independent implementations of the reconstruction agree. Item (iv) is a
$2\times2$ hand check on the $\lambda_3$ component. Everything is integer
arithmetic in the fixed order $\Z\eta_0\oplus\cdots\oplus\Z\eta_3$; no number
field beyond the explicit quartic periods, no class-group input, and no solver
is involved.
\end{proof}

The same criterion applied out of panel at $q=6425=5^2\cdot257$ (the next
$q\equiv1\bmod4$ with $51\mid t$, the same $V'$, and a two-square
representation) is again infeasible at $t'=51$, with zero surviving side
vectors. As a positive control, relaxing the $q=1325$ fiber boxes to
$|\,\cdot\,|\le29$ produces eight exact joins, each verified as an identity in
the orbit algebra; the $q=1325$ obstruction is therefore genuinely a fiber-box
obstruction, not an artifact of the component compatibility search.

\subsection{The divisor lattice and existential selection}
\label{sec:divisor-lattice}

The criteria above are stated one divisor at a time.  This subsection
organizes them over the whole divisor lattice of $t$ and proves the
structural layer of the existential divisor-selection problem: marginal
feasibility is monotone along the lattice, the orbit model at every divisor
is computed by $M(q)$ alone, criterion-equipped divisors always exist, and
the top node is not a relaxation but the exact multiplier-reduced decision.
The consequence is that, for composite all-plus $q$, the converse statement
``no Turyn pair exists at $q$'' is \emph{equivalent} to the existential
statement ``some divisor's marginal system is empty.''  The selection
quantifier itself is therefore no longer an obligation; what remains open
for the full converse is exactly the per-family firing content, made precise
in Remark~\ref{rem:selection-open} below.

Throughout this subsection $q=2t-1$ is composite with $t$ odd, the composite
multiplier signs are all $+1$, and for $t'\mid t$ we write $S(t')$ for the
finite (possibly empty) solution set of the marginal system of
Lemma~\ref{lem:marginal-algebra} at $t'$: pairs of $V'$-invariant vectors in
$\mathcal A(t',V')$ satisfying the orbit-algebra equations, the
trivial-character equation, and the fiber boxes and parities for $h=t/t'$.

\begin{lemma}[Lattice functoriality]\label{lem:lattice-functoriality}
Let $t''\mid t'\mid t$ and let $\pi:\Z[C_{t'}]\to\Z[C_{t''}]$ be the natural
surjection.  Then $\pi(S(t'))\subseteq S(t'')$.  Consequently
\[
  U(q)=\{\,t'\mid t:\ S(t')=\emptyset\,\}
\]
is an up-set in the divisor lattice of $t$: if $t''\in U(q)$ and
$t''\mid t'\mid t$, then $t'\in U(q)$.
\end{lemma}

\begin{proof}
Write $h=t/t'$ and $h'=t'/t''$, so the fiber parameter at $t''$ is
$h'h=t/t''$; both $h$ and $h'$ are odd because $t$ is odd.  The map $\pi$ is
a ring homomorphism, so it preserves the equation $AA^*+BB^*=q\cdot1$, and
evaluation at the trivial character is unchanged, so the trivial-character
equation is preserved.  The image of $V'$ under
$(\Z/t')^\times\to(\Z/t'')^\times$ is $V''$ (both are images of $V$), so
$\pi$ carries $V'$-invariance to $V''$-invariance.  For the boxes and
parities: an off-origin cell of $\pi(A)$ is a sum of $h'$ cells of $A$, each
odd with absolute value at most $h$; a sum of $h'$ odd integers is odd, and
the bound is $h'h$.  The origin cell of $\pi(A)$ is the origin cell of $A$
(even, absolute value at most $h-1$) plus $h'-1$ odd cells, hence even with
absolute value at most $(h-1)+(h'-1)h=h'h-1$.  The $B$ side is the same
computation without the origin exception.  These are exactly the boxes and
parities of Lemma~\ref{lem:marginal-algebra} at $t''$ with fiber parameter
$h'h$.  The up-set property is the contrapositive: if
$S(t')\ne\emptyset$ then $S(t'')\supseteq\pi(S(t'))\ne\emptyset$.
\end{proof}

\begin{lemma}[Orbit census]\label{lem:orbit-census}
For $f\mid t$ let $V_f$ denote the image of $V$ in $(\Z/f)^\times$.  Then for
every divisor $t'\mid t$:
\emph{(i)} the $V'$-orbit of an element $j\in\Z/t'$ with $\gcd(j,t')=t'/f$
has size $|V_f|$, the number of such orbits is $\varphi(f)/|V_f|$, and
\[
  \#\Omega(t',V')=\sum_{f\mid t'}\frac{\varphi(f)}{|V_f|}\,;
\]
\emph{(ii)} the rational Wedderburn decomposition of
Proposition~\ref{prop:component-join} is indexed by conductors,
\[
  \Q[C_{t'}]^{V'}\;\simeq\;\prod_{f\mid t'}\Q(\zeta_f)^{V_f},
\]
and the conductor-$f$ component is a totally real field of degree
$\varphi(f)/|V_f|$.
In particular the fiber parameter, orbit sizes, orbit count, and component
degree profile at every divisor of $t$ are determined by $t$ and $M(q)$
alone, before any enumeration.
\end{lemma}

\begin{proof}
(i) The elements of $\Z/t'$ with $\gcd(j,t')=t'/f$ are the $(t'/f)j_0$ with
$j_0\in(\Z/f)^\times$, and $u\in V'$ acts by $j_0\mapsto uj_0\bmod f$, i.e.,
through $V_f$.  The multiplication action of $V_f$ on $(\Z/f)^\times$ is
free, so every orbit on this stratum has size $|V_f|$ and there are
$\varphi(f)/|V_f|$ of them; summing over $f\mid t'$ exhausts $\Z/t'$.
(ii) The characters of $C_{t'}$ of conductor $f$ span a block of
$\Q[C_{t'}]$ isomorphic to $\Q(\zeta_f)$, on which $V'$ acts through the
Galois action of $V_f$; the fixed part is $\Q(\zeta_f)^{V_f}$, of degree
$\varphi(f)/|V_f|$, totally real because $-1\in V_f$.  The degrees sum to
the orbit count of (i), as both equal $\dim_\Q\Q[C_{t'}]^{V'}$.
\end{proof}

\begin{lemma}[Exactness at the top node]\label{lem:top-exact}
$S(t)\ne\emptyset$ if and only if a Turyn pair of order $t$ exists.
\end{lemma}

\begin{proof}
At $t'=t$ the fiber parameter is $h=1$: the boxes force the $A$-origin to be
$0$ and every other cell of $A$ and every cell of $B$ to be a sign.  If a
Turyn pair exists, it is itself $V$-invariant by
Theorem~\ref{thm:composite-multiplier} in the all-plus case (no translate is
needed) and symmetric, hence defines an element of $S(t)$.  Conversely, an
element of $S(t)$ is a pair $(A,B)$ of $V$-invariant sign rows with zero
$A$-origin; since $-1\in V$, both rows are symmetric, so $A^*=A$ and
$B^*=B$, and the marginal equation reads $A^2+B^2=q\cdot1$: a Turyn pair.
Compare Proposition~\ref{prop:t3-decision}, of which this is the
solution-set form.
\end{proof}

\begin{theorem}[Existential divisor selection]\label{thm:divisor-selection}
Let $q=2t-1$ be composite with $t$ odd and all composite multiplier signs
$+1$.  Then:
\emph{(i)} no Turyn pair exists at $q$ if and only if $S(t')=\emptyset$ for
some divisor $t'\mid t$, if and only if $t\in U(q)$, if and only if
$U(q)\ne\emptyset$;
\emph{(ii)} $U(q)$ is an up-set, so its minimal elements form an antichain,
and certifying $S(t')=\emptyset$ at any single node closes $q$;
\emph{(iii)} the divisors carrying a written finite criterion are nonempty
for every $t$ and are located by Lemma~\ref{lem:orbit-census}: every prime
$\ell\mid t$ carries the period component criterion
(Proposition~\ref{prop:prime-subtorus-component},
Corollary~\ref{cor:prime-subtorus-firing}) of degree
$(\ell-1)/|V_\ell|$; if $3\mid t$, the node $t'=3$ is transitive two-orbit
(Theorem~\ref{thm:t2-twoorbit}, since $V_3=(\Z/3)^\times$) and the node
$t'=t/3$ carries the $h=3$ support hierarchy
(Proposition~\ref{prop:h3-support-augmentation},
Corollary~\ref{cor:h3-support-paf}); and the top node $t'=t$ carries the
exact sign decision (Lemma~\ref{lem:top-exact}; for prime $t$,
Corollary~\ref{cor:full-torus-sign});
\emph{(iv)} if $t$ is composite and every proper divisor node is feasible,
then $U(q)\subseteq\{t\}$, and nonexistence at $q$ is equivalent to
emptiness of the top node alone.
\end{theorem}

\begin{proof}
(i) If $S(t')=\emptyset$ for some $t'$, then no Turyn pair exists by
Proposition~\ref{prop:t2-marginal}.  If no Turyn pair exists, then
$S(t)=\emptyset$ by Lemma~\ref{lem:top-exact}, so $t\in U(q)$ and
$U(q)\ne\emptyset$.  Finally $U(q)\ne\emptyset$ implies, by the up-set
property of Lemma~\ref{lem:lattice-functoriality}, that $t\in U(q)$, hence
$S(t)=\emptyset$ and no pair exists.  (ii) is
Lemma~\ref{lem:lattice-functoriality}.  (iii) $V_3\supseteq\{\pm1\}=
(\Z/3)^\times$, so the $t'=3$ orbit census is one zero orbit plus one
transitive nonzero orbit; the remaining memberships are the stated
hypotheses of the cited criteria, all decided by
Lemma~\ref{lem:orbit-census}.  (iv) is (i) together with the up-set
property.
\end{proof}

\begin{corollary}[Projection-image criterion]\label{cor:projection-image}
Let $t''_1,\dots,t''_k$ and $t'$ be divisors of $t$ with $t''_i\mid t'$ for
each $i$, and let $R(t')\supseteq S(t')$ be any finite relaxation defined by
a subset of the marginal constraints at $t'$ (for instance the boxes,
parities, and $V'$-invariance alone).  If
\[
  \bigl\{\,x\in R(t'):\ \pi_i(x)\in S(t''_i)\ \text{for }i=1,\dots,k\,\bigr\}
  =\emptyset,
\]
then $S(t')=\emptyset$ and no Turyn pair exists at $q$.
\end{corollary}

\begin{proof}
$S(t')\subseteq R(t')$, and $\pi_i(S(t'))\subseteq S(t''_i)$ by
Lemma~\ref{lem:lattice-functoriality}; so $S(t')$ is contained in the
displayed set.  Emptiness then gives $S(t')=\emptyset$, and
Proposition~\ref{prop:t2-marginal} applies.
\end{proof}

The $q=441$ obstruction of \S\ref{sec:gluing} is the instance
$t'=221$, $\{t''_1,t''_2\}=\{13,17\}$, with $R$ the per-side boxed
$V$-invariant rows: the $C_{13}$-projection of any boxed cover row is
$\equiv\pm1\pmod 8$ on nonzero orbits, while every element of $S(13)$ has
nonzero side values $\pm3$, so already the $A$-side compatible set is
empty.  Corollary~\ref{cor:projection-image} is the general form of that
obstruction, available at every node of the lattice against any family of
exactly solved quotient nodes.

\begin{remark}[What the selection theorem does and does not provide]
\label{rem:selection-open}
Theorem~\ref{thm:divisor-selection} closes the quantification layer of the
divisor-selection problem: the correct quantifier is existential over all
divisors of $t$; it is complete (part (i)); the search space it ranges over
is monotone (part (ii)), so a nonexistence proof may stop at any certified
node; family membership at every node is computed from $M(q)$ by
Lemma~\ref{lem:orbit-census} before any enumeration; and criterion-equipped
nodes always exist (part (iii)).  It does \emph{not} prove that a written
criterion fires at some node for every composite $q$: by part (i), that
statement is equivalent to Conjecture~\ref{conj:converse} itself for the
all-plus branch.  The open content of the full converse is therefore
isolated in per-family firing theorems --- when is the period join of
Corollary~\ref{cor:prime-subtorus-firing} empty, when does the $h=3$
hierarchy fire, and their analogues --- together with the top-node analysis
of part (iv).  The finite-panel branch accounting of
\S\ref{sec:finite-panel-accounting} is the instantiation of the theorem at
$q\le2000$: each closure records one certified minimal element of $U(q)$.
Two panel facts are worth recording in lattice terms.  First, only $7$ of
the $21$ recent panel closures occur at the divisor preferred by the
shape-ranking heuristic: minimal elements of $U(q)$ are not located by
orbit shape alone, which is why the quantifier must remain existential.
Second, every composite-$t$ case in the panel has a certified \emph{proper}
minimal node --- the top node was needed only for prime $t$, where no proper
node exists.  Thus the family of part (iv), composite $t$ with all proper
nodes feasible, has no known member with $q\le2000$; whether it is empty is
precisely the gluing question, now stated as a property of the lattice
rather than of one example.
\end{remark}

\paragraph{An explicit lane where the converse holds.}
The selection theorem and the $h=3$ closed forms combine into the
program's first infinite-family-shaped partial converse: on the lane
$q=6\ell-1$ with $\ell\equiv3\pmod4$ prime, the firing conditions hold
\emph{deterministically} for every composite $q$ with two prime
factors --- no density or equidistribution input is needed.  Three
elementary facts collapse the lane.

\begin{lemma}[Anchor: no $\pm1$ factors at composite shifted primes]
\label{lem:anchor-pm1}
Let $\ell\ge3$ be odd and let $q=6\ell-1$ be composite.  Then no prime
divisor of $q$ is congruent to $\pm1\pmod\ell$.  If moreover $\ell$ is
prime, every prime divisor $p$ of $q$ has $\ord_\ell(p)\ge3$; and if
$\ell\equiv3\pmod4$, the odd part of $\ord_\ell(p)$ is at least $3$.
\end{lemma}

\begin{proof}
An odd $p>1$ with $p\equiv\pm1\pmod\ell$ satisfies $p\ge2\ell-1$,
because $\ell\pm1$ are even.  Then $c=q/p<(6\ell-1)/(2\ell-1)<7/2$, and
$c$ is odd, so $c\in\{1,3\}$; $c=3$ contradicts $q\equiv-1\pmod3$, and
$c=1$ contradicts compositeness.  The order statements follow since
$\ord_\ell(p)\le2$ means $p\equiv\pm1\pmod\ell$, and
$v_2(\ord_\ell(p))\le v_2(\ell-1)=1$ when $\ell\equiv3\pmod4$.
\end{proof}

\begin{lemma}[Semiprime multiplier image]\label{lem:semiprime-image}
Let $\ell\equiv3\pmod4$ be prime and $q=6\ell-1=p_1p_2$ with
$p_1\ne p_2$ prime, both $\equiv1\pmod4$.  Let $x=p_1\bmod\ell$,
$n_0$ the odd part of $\ord_\ell(x)$.  Then
$\{p_1,p_2\}\equiv\{1,5\}\pmod{12}$, the multiplier image satisfies
$M(q)|_\ell=\langle x^2\rangle$, cyclic of odd order $n_0\ge3$
(independent of which factor is used), all multiplier signs are $+1$,
and
\[
  w=\bigl|V_\ell\bigr|=2n_0\ge6,\qquad d=\frac{\ell-1}{w}\ \text{odd}.
\]
\end{lemma}

\begin{proof}
$q\equiv5\pmod{12}$ forces $\{p_1,p_2\}\bmod12\in\{\{1,5\},\{7,11\}\}$,
and $p_i\equiv1\pmod4$ selects $\{1,5\}$; all powers of the $p_i$ are
$\equiv1\pmod4$, so every multiplier sign is $+1$.  Write elements of
$(\Z/12\ell)^\times\cong(\Z/4)^\times\times(\Z/3)^\times\times(\Z/\ell)^\times$
as triples and let $y=p_2\bmod\ell$.  With $p_1\equiv1$,
$p_2\equiv5\pmod{12}$ (the mirror case is symmetric),
$\langle p_1\rangle=\{(1,1,x^a)\}$ and
$\langle p_2\rangle=\{(1,2^b,y^b)\}$, so intersecting forces $b$ even
and $M(q)|_\ell=\langle x\rangle\cap\langle y^2\rangle$.  From
$q\equiv-1\pmod\ell$ we get $y=-x^{-1}$, hence
$\langle y^2\rangle=\langle x^2\rangle\subseteq\langle x\rangle$ and
$M(q)|_\ell=\langle x^2\rangle$, of order the odd part $n_0$ of
$\ord_\ell(x)$ (as $v_2(\ord_\ell(x))\le1$); $n_0\ge3$ by
Lemma~\ref{lem:anchor-pm1}, and $y=-x^{-1}$ gives the same odd part on
the other side.  Finally $-1\notin\langle x^2\rangle$ (odd order), so
$w=|\langle x^2,-1\rangle|=2n_0$, and $v_2(w)=1=v_2(\ell-1)$ makes $d$
odd.
\end{proof}

\begin{theorem}[Semiprime partial converse on the lane
$q\equiv5\pmod6$]\label{thm:semiprime-converse}
Let $\ell\equiv3\pmod4$ be prime and let $q=6\ell-1$ be composite with
$\Omega(q)=2$.  Then no Turyn pair of order $t=3\ell$ exists.
\end{theorem}

\begin{proof}
$q$ is not a square: $q\equiv-1\pmod\ell$ and $-1$ is not a quadratic
residue for $\ell\equiv3\pmod4$.  So $q=p_1p_2$ with distinct primes.
If some $p_i\equiv3\pmod4$, both are (mod-$12$ classification), each to
the first power, and Proposition~\ref{prop:twosquares} applies.
Otherwise Lemma~\ref{lem:semiprime-image} gives all-plus signs,
$w=2n_0\ge6$ and $d$ odd at the divisor $t'=\ell$ ($h=3$, the node of
Theorem~\ref{thm:divisor-selection}(iii)); then
$(\ell-1)/2=(w/2)d\equiv w/2\pmod w$ and $w/2\ge3\notin\{0,1\}$, so the
$h=3$ support stage fires by
Corollary~\ref{cor:h3-support-sum-prime}, and
Proposition~\ref{prop:t2-marginal} concludes.
\end{proof}

\begin{theorem}[Two-prime-support extension]\label{thm:omega2-converse}
Let $\ell\equiv3\pmod4$ be prime and
$q=6\ell-1=p_1^{e_1}p_2^{e_2}$ composite, $p_1\ne p_2$ prime.  Let
$n_i$ be the odd part of $\ord_\ell(p_i)$ and
$m_i=n_i/\gcd(n_i,\operatorname{oddpart}(e_i))$.  If some
$p_i\equiv3\pmod4$ has $e_i$ odd, no Turyn pair exists by two squares.
Otherwise, if $\max(m_1,m_2)\ge3$ --- in particular whenever some $e_i$
is a power of two --- no Turyn pair of order $t=3\ell$ exists.
\end{theorem}

\begin{proof}
The signed branch would force $q$ to be a square
(Lemma~\ref{lem:signed-branch-reduction}), excluded as above; so the
signs are all $+1$.  Squaring $x^{e_1}y^{e_2}\equiv-1\pmod\ell$ gives
$\langle x^{2e_1}\rangle=\langle y^{2e_2}\rangle=:H$; even powers of
$p_1$ and $p_2$ are trivial mod $12$, so $H\subseteq M(q)|_\ell$.  The
order of $H$ has odd part $m_1$ (equivalently $m_2$), so
$w\ge2\max(m_1,m_2)\ge6$ with $v_2(w)=1$, and the firing argument of
Theorem~\ref{thm:semiprime-converse} repeats.
\end{proof}

\begin{theorem}[Conditional infinitude]\label{thm:dickson-family}
Assume Dickson's conjecture~\cite{Dickson04} for the admissible pair
$\ell(c)=20c+11$, $v(c)=24c+13$.  Then there are infinitely many $c$
with both values prime, and for each, $q=5v(c)=6\ell(c)-1$ satisfies
the hypotheses of Theorem~\ref{thm:semiprime-converse}.  Hence,
conditionally, Turyn pairs fail to exist for infinitely many composite
$q$; the first members are $q=65,185,785$.
\end{theorem}

\begin{theorem}[Dichotomy under a congruence-restricted Chen
input]\label{thm:chen-dichotomy}
Assume the following analytic hypothesis \textup{(C)}: there are
infinitely many primes $\ell\equiv3\pmod4$ with
$\Omega(6\ell-1)\le2$.  For every such $\ell$, the converse
Conjecture~\ref{conj:converse} holds at $q=6\ell-1$: if $q$ is prime a
Turyn pair exists and the conjecture is vacuous there; if $q$ is
composite, Theorem~\ref{thm:semiprime-converse} denies a Turyn pair.
Consequently, under \textup{(C)}, at least one of the following is
true:
(i) $6\ell-1$ is prime for infinitely many primes $\ell\equiv3\pmod4$;
(ii) there are infinitely many composite $q$ admitting no Turyn pair.
\end{theorem}

\begin{remark}[Status of hypothesis (C)]\label{rem:chen-input}
The 2026-07-06 literature audit found no retrievable statement covering
exactly \textup{(C)}.  What is in print: Chen's $P_2$ bound~\cite{Chen73} for the
degree-one case $p+2$; Richert's weighted-sieve theorem giving
$f(p)=P_{2\deg f+1}$ for irreducible $f$, whose degree-one
specialization covers $6p-1=P_2$ over \emph{unrestricted} primes $p$
(attribution as summarized by Irving~\cite{Irving14}); and the abundance of Chen
primes in fixed arithmetic progressions for the form $p+2$ (Lewulis~\cite{Lewulis16}).
The missing step is only the residue restriction
$\ell\equiv3\pmod4$ on the sieving variable, which the standard
weighted-sieve machinery with fixed modulus is expected to absorb
routinely; until a derivation appendix or a primary citation is in
place, the dichotomy is stated conditionally on \textup{(C)} rather
than attributed to Chen.
\end{remark}

\begin{remark}[The parity wall, and the lane's calibration]
\label{rem:parity-wall}
Unconditional infinitude of the composite members is a parity-problem
instance: ``$\Omega(6\ell-1)=2$ for infinitely many primes $\ell$'' is
of the same class as ``$p+2$ is a semiprime infinitely often,'' open
despite Chen's theorem, and the standard $E_2$/GPY technology does not
apply because $\ell$ itself must be prime.  For $\omega(q)\ge3$ the
deterministic collapse fails structurally: the product relation
$\prod p_i^{e_i}\equiv-1\pmod\ell$ constrains the join of the cyclic
groups $\langle p_i\rangle$, never their intersection.  Calibration at
$q\le10^7$ ($62{,}929$ lane primes $\ell$): $18.7\%$ of lane members
are Theorem~\ref{thm:semiprime-converse} kills, $18.5\%$ die
classically, $17.9\%$ have $q$ prime, and $41.4\%$ sit behind the
parity wall --- of which $99.98\%$ pass the gcd-of-orders proxy for
$w\ge4$, so the obstruction to closing them is proof technology, not
truth.  The first Theorem-\ref{thm:semiprime-converse} members in the blind-spot
$\{1,5\}\bmod12$ class (the two-squares members, e.g.\ $q=473,497$, fall to
the classical obstruction) are $q=65$, $185$, $785$, $905$, $1073$, all
verified against the exact
multiplier computation, and the framework provably cannot collide with
Turyn's construction: prime $q$ has its own factor $\equiv-1\pmod\ell$
($w=2$, every criterion silent), and odd prime powers reach a prime
$\ell$ only through the $p=5$ tower, which lands $\ell\equiv1\pmod4$.
\end{remark}

\subsection{The lane at general \texorpdfstring{$\omega$}{omega}: an
averaged-equidistribution reduction}\label{sec:lane-density}

Remark~\ref{rem:parity-wall} localizes the failure of
Theorem~\ref{thm:semiprime-converse} at $\omega(q)\ge3$ and prices
unconditional infinitude at a parity problem.  This subsection removes
the first limitation and reprices the second.  At general $\omega$ the
multiplier image on the lane is still computed by a gcd of orders
(Lemma~\ref{lem:lane-image}), the criteria still fire deterministically
whenever that image exceeds $\{\pm1\}$
(Proposition~\ref{prop:lane-kill}), and a character-pinning identity
(Lemma~\ref{lem:character-pinning}) determines the residue status of the
one possible prime factor of $q$ above $\sqrt q$ from the small factors.
The open analytic content is then isolated in an \emph{averaged
equidistribution} input of Bombieri--Vinogradov type for cubic Kummer
splitting (Hypothesis $\mathrm{EQ}_3$ below) together with a structural
control of the factors of $6\ell-1$ above the sieve level; under the
hypothesis alone we prove the small-factor obstruction count
(Proposition~\ref{prop:small-factor-sieve}).  Everything before the
hypothesis is unconditional.

\begin{lemma}[Multiplier image on the lane at general $\omega$]
\label{lem:lane-image}
Let $\ell\equiv3\pmod4$ be prime and let
$q=6\ell-1=p_1^{e_1}\cdots p_k^{e_k}$ be composite.  Put
$n_i=\ord_\ell(p_i)$, $g=\gcd(n_1,\dots,n_k)$, and let $g_0$ be the odd
part of $g$.  Then the image $M(q)|_\ell$ of $M(q)$ in
$(\Z/\ell)^\times$ is cyclic of order $g_0$ or $2g_0$, and
\[
  V_\ell=\bigl\langle M(q)|_\ell,\,-1\bigr\rangle,\qquad
  w=|V_\ell|=2g_0,\qquad v_2(w)=1,\qquad
  d=\frac{\ell-1}{w}\ \text{odd},
\]
independent of the exponents $e_i$.  Moreover every sign branch reduces
to the all-plus branch: $q\equiv5\pmod{12}$ is never a square, so
Lemma~\ref{lem:signed-branch-reduction} excludes the sign-alternating
case.
\end{lemma}

\begin{proof}
All prime factors of $q$ are coprime to $12$, and
$(\Z/12\ell)^\times\cong(\Z/12)^\times\times(\Z/\ell)^\times$; write
$x_i=p_i\bmod\ell$.  Since $v_2(\ell-1)=1$, every $n_i$ has
$v_2(n_i)\le1$, and $-1$ is the unique involution of the cyclic group
$(\Z/\ell)^\times$, contained in every even-order subgroup.

\emph{Lower bound.}  Every odd square is $1\bmod4$ and every square
prime to $3$ is $1\bmod3$, so $p_i^2\equiv1\pmod{12}$ and
$\langle p_i^2\rangle=\{1\}\times\langle x_i^2\rangle$.  Hence
$K:=\bigcap_i\langle p_i^2\rangle\subseteq M(q)$ projects isomorphically
to $\bigcap_i\langle x_i^2\rangle$.  The order of $x_i^2$ is the odd
part of $n_i$ (as $v_2(n_i)\le1$); in a cyclic group an intersection of
subgroups has order the gcd of their orders, and odd part commutes with
gcd, so $|K|=\gcd_i\mathrm{oddpart}(n_i)=g_0$.  As $g_0$ is odd,
$-1\notin K|_\ell$, and $w\ge2g_0$.

\emph{Upper bound.}  Projecting,
$M(q)|_\ell\subseteq\bigcap_i\langle x_i\rangle=C_g$, the unique
subgroup of order $g$.  If $g$ is odd then
$|\langle C_g,-1\rangle|=2g=2g_0$; if $g$ is even then $v_2(g)=1$, so
$C_g$ already contains the involution and
$\langle C_g,-1\rangle=C_g$ has order $g=2g_0$.  Either way
$w\le2g_0$; the $2$-part of $M(q)|_\ell$ itself never needs to be
resolved.

Exponent-independence is definitional
($M(q)=\bigcap_{p\mid q}\langle p\rangle$ sees only the distinct
primes), and $w=2g_0\mid\ell-1$ with $g_0$ odd makes
$d=(\ell-1)/(2g_0)$ odd.
\end{proof}

For $k=2$ the formula specializes to
Lemma~\ref{lem:semiprime-image} ($g_0=n_0$).  Notably, the proof never
uses the product relation $\prod_i p_i^{e_i}\equiv-1\pmod\ell$ that
drove the semiprime argument: both bounds are relation-free.  The
relation returns as the central actor in
Lemma~\ref{lem:character-pinning}.

\begin{proposition}[Deterministic kill at general $\omega$]
\label{prop:lane-kill}
Let $\ell\equiv3\pmod4$ be prime and $q=6\ell-1$ composite.  If
$g_0=\mathrm{oddpart}\bigl(\gcd_{p\mid q}\ord_\ell(p)\bigr)\ge3$, then
no Turyn pair exists at $q$.  Equivalently: a Turyn pair at composite
$q=6\ell-1$ forces $w=2$, i.e.\ the multiplier image in
$(\Z/\ell)^\times$ collapses to $\{\pm1\}$.
\end{proposition}

\begin{proof}
Lemma~\ref{lem:lane-image} gives $w=2g_0\ge6$ and $d$ odd, so
$(\ell-1)/2=(w/2)\,d\equiv w/2\pmod w$ with $w/2\ge3$, hence
$(\ell-1)/2\bmod w\notin\{0,1\}$ and the prime support-sum closed form
(Corollary~\ref{cor:h3-support-sum-prime}) fires at the node
$t'=\ell$ of $t=3\ell$; Proposition~\ref{prop:t2-marginal} concludes.
\end{proof}

\begin{lemma}[Residue reformulation of the bad set]
\label{lem:lane-residue}
With $\ell,q$ as above and, for an odd prime $r\mid\ell-1$, write
$S_r\le(\Z/\ell)^\times$ for the subgroup of $r$th powers.
\begin{enumerate}
\item If some odd prime $r\mid\ell-1$ has $p\bmod\ell\notin S_r$ for
  every $p\mid q$, then $g_0\ge r\ge3$ and
  Proposition~\ref{prop:lane-kill} kills $q$.
\item If $g_0=1$, then for every odd prime $r\mid\ell-1$ some
  $p\mid q$ has $p\bmod\ell\in S_r$.
\end{enumerate}
\end{lemma}

\begin{proof}
$S_r$ is the unique subgroup of order $(\ell-1)/r$, so
$x\in S_r\iff\ord(x)\mid(\ell-1)/r$.  If $x\notin S_r$ then
$r\mid\ord(x)$, giving (i); if $r\nmid\ord(x)$ then
$\ord(x)$ divides $(\ell-1)/r^{v_r(\ell-1)}$, which divides
$(\ell-1)/r$, so $x\in S_r$, giving (ii).
\end{proof}

\begin{lemma}[Character pinning]\label{lem:character-pinning}
Let $\ell\equiv3\pmod4$ be prime, $q=6\ell-1$ composite,
$r\mid\ell-1$ an odd prime, and $\chi$ a character of
$(\Z/\ell)^\times$ of order $r$.  Then $\chi(-1)=1$, so
$q\equiv-1\pmod\ell$ gives $\prod_{p^e\parallel q}\chi(p)^{e}=1$.
Consequently, writing $q=mP^{\epsilon}$ with $P$ the (at most one,
automatically first-power) prime factor exceeding $\sqrt q$ and $m$ the
$\sqrt q$-smooth cofactor,
\[
  \epsilon=1\ \Longrightarrow\
  \chi(P)=\prod_{p^{e}\parallel m}\chi(p)^{-e}:
\]
the residue status of $P$ at every odd prime $r\mid\ell-1$ is
determined by the factors of $q$ below $\sqrt q$.
\end{lemma}

\begin{proof}
$\chi(-1)^2=1$ and $\chi(-1)\in\mu_r$ with $r$ odd force $\chi(-1)=1$;
apply $\chi$ to $q\equiv-1\pmod\ell$.  Two prime factors $>\sqrt q$
would multiply past $q$.  Finally $P\notin S_r\iff\chi(P)\ne1$, since
$S_r=\ker\chi$.
\end{proof}

The lane's $\ell$-dependent prime factors are the reason no fixed
Chebotarev statement applies directly: the splitting field of ``$p$ is
an $r$th power residue mod $\ell$'' varies with $p\mid6\ell-1$.
Lemma~\ref{lem:character-pinning} removes the one factor that cannot be
averaged over (the large one); for the small factors, the right input
is equidistribution \emph{on average over the base}.  We restrict to
the sub-lane $\ell\equiv7\pmod{12}$ (equivalently $3\mid\ell-1$), where
$r=3$ is always available.

\begin{definition}[Cubic Kummer average]\label{def:eq3}
For squarefree $d=p_1\cdots p_r$ with $(d,6)=1$ let
\[
  M_d=\Q\bigl(\zeta_3,\;p_1^{1/3},\dots,p_r^{1/3},\;\zeta_{12d}\bigr),
  \qquad G_d=\Gal(M_d/\Q),
\]
and let $C_d\subseteq G_d$ be the union of conjugacy classes cut out by
the conditions: $\ell\equiv7\pmod{12}$, and for each $i$,
$\ell\equiv6^{-1}\pmod{p_i}$ and $p_i$ a cubic residue mod $\ell$ (for
$\ell\equiv1\bmod3$ the latter says $\ell$ splits completely in
$\Q(\zeta_3,p_i^{1/3})$).  Its density is an \emph{exact} product,
\[
  \delta_d:=\frac{|C_d|}{|G_d|}
  =\frac14\prod_{i=1}^{r}\frac1{3\,\varphi(p_i)},
\]
because the Kummer layers are independent
($\Gal(\Q(\zeta_3,p_1^{1/3},\dots,p_r^{1/3})/\Q(\zeta_3))\cong(\Z/3)^r$:
no product $\prod p_i^{a_i}$, $a_i\in\{0,1,2\}$ not all zero, is a
cube) and no cubic Kummer subfield is abelian over $\Q$, so the Kummer
part is linearly disjoint from $\Q(\zeta_{12d})$ over $\Q(\zeta_3)$.
For fixed $\theta\in(0,\tfrac12)$, \emph{Hypothesis}
$\mathrm{EQ}_3(\theta)$ asserts: for every $A>0$,
\[
  \sum_{\substack{d\le x^{2\theta}\\ d\ \text{squarefree},\ (d,6)=1}}
  \Bigl|\,\pi_{C_d}(x;M_d)-\delta_d\operatorname{Li}(x)\Bigr|
  \ \ll_A\ \frac{x}{(\log x)^{A}}\,.
\]
(The variant with prescribed cubic-symbol values at each $p_i$, needed
for character-twisted sums, is defined analogously; the two nontrivial
symbol values are exchanged by conjugation in each $S_3$ layer, so
only their real combinations are class functions, which is all such
sums use.)
\end{definition}

\begin{proposition}[Small-factor obstruction count]
\label{prop:small-factor-sieve}
Fix $\theta\in(0,\tfrac12)$ and assume $\mathrm{EQ}_3(\theta)$.  There
is $c(\theta)>0$ such that, for all large $x$, the number of primes
$\ell\le x$ with $\ell\equiv7\pmod{12}$, $q=6\ell-1$ composite, and
every prime $p\le x^{\theta}$ dividing $q$ a cubic non-residue mod
$\ell$, is at least $c(\theta)\,x/(\log x)^{4/3}$.
\end{proposition}

\begin{proof}
Let $z=x^\theta$, $D=z^2=x^{2\theta}$, and let $\mathcal A$ be the set
of primes $\ell\le x$ with $\ell\equiv7\pmod{12}$, of size
$X=(\tfrac14+o(1))\,x/\log x$ (prime number theorem in the fixed
progression).  For a prime $p\le z$, $(p,6)=1$, put
$E_p=\{\ell\in\mathcal A:\ p\mid6\ell-1$ and $p$ a cubic residue mod
$\ell\}$; note $2,3\nmid q$ automatically.  For squarefree $d$ composed
of such primes,
\[
  A_d:=\Bigl|\bigcap_{p\mid d}E_p\Bigr|
      =\pi_{C_d}(x;M_d)+O\bigl(\omega(d)\bigr)
      =g(d)\,X+r_d,\qquad
  g(d)=\prod_{p\mid d}\frac{1+o(1)}{3\varphi(p)},
\]
with the product density exact (Definition~\ref{def:eq3}) and
$\sum_{d\le D}|r_d|\ll_A x(\log x)^{-A}$ by $\mathrm{EQ}_3(\theta)$
(the $O(\omega(d))$ terms contribute $O(D\log D)$, absorbed).  The
density has sieve dimension $\kappa=\tfrac13$, since
$g(p)=\tfrac13\cdot\tfrac{1+o(1)}{p-1}$ on the sieving primes.  The
beta-sieve lower bound of dimension $\tfrac13$ at
$s=\log D/\log z=2$ applies --- the sieving limit satisfies
$\beta(\tfrac13)\le\beta(\tfrac12)=1<2$, the half-dimensional value
(see \cite[Ch.~11]{FI10}, \cite{HR74}) --- and gives
\[
  \#\{\ell\in\mathcal A:\ \ell\notin E_p\ \forall p\le z\}
  \ \ge\ \bigl(f_{1/3}(2)-o(1)\bigr)\,X\prod_{p\le z}\bigl(1-g(p)\bigr)
  \ -\ \sum_{d\le D}|r_d|,
\]
with $f_{1/3}(2)>0$.  By the Mertens-type theorem for sieve densities
of dimension $\kappa$ (\cite{HR74}),
$\prod_{p\le z}(1-g(p))\asymp_\theta(\log x)^{-1/3}$, so the main term
is $\gg_\theta x/(\log x)^{4/3}$ and the remainder is negligible.  A
survivor $\ell$ has every prime factor $p\le z$ of $q$ a cubic
non-residue; survivors with $q$ prime number
$O(x/\log^2x)=o\bigl(x/(\log x)^{4/3}\bigr)$ (upper-bound sieve for
the pair $(\ell,6\ell-1)$, \cite{HR74}) and are removed.
\end{proof}

\begin{remark}[From the small factors to a family: status]
\label{rem:lane-density-status}
Three calibrations.  \emph{(i) What is not yet proved.}  Members of the
set in Proposition~\ref{prop:small-factor-sieve} are not certified
nonexistence orders: $q$ may have prime factors in
$(x^\theta,\sqrt q\,]$ about which the proposition says nothing, and
only the at-most-one factor above $\sqrt q$ is pinned by
Lemma~\ref{lem:character-pinning}.  Converting the count into an
infinite family of Proposition~\ref{prop:lane-kill} kills requires
exactly (a) a structure count --- sub-lane primes with
$6\ell-1=(x^\theta\text{-smooth})\cdot(\text{prime})$ --- which the
linear sieve at Bombieri--Vinogradov level does not reach (sifting the
middle range to $\sqrt{6x}$ at level $x^{1/2-\varepsilon}$ has $s=1$,
below the linear sieving limit; per fixed smooth cofactor the count is
a two-linear-forms prime pair, a parity-type problem), and (b) a joint
(vector-sieve) average running the Kummer conditions relative to the
structured set, which $\mathrm{EQ}_3$ as stated does not supply.  A
2026-07-06 bookkeeping attempt at (a)+(b) was retracted on
verification, gapped at both points; the almost-all version of the
statement (nonexistence at all but $o(x/\log x)$ of the composite lane
members, via Lemma~\ref{lem:lane-residue}(ii) at moderately large
$r\mid\ell-1$) reduces to $r$th-power analogues of $\mathrm{EQ}_3$
uniformly for odd $r\le(\log x)^{1+\epsilon}$ \emph{plus the same two
open points}.  \emph{(ii) GRH tier.}  Under GRH,
$\mathrm{EQ}_3(\theta)$ holds for every $\theta<\tfrac18$: the
effective Chebotarev theorem of Lagarias--Odlyzko~\cite{LO77} bounds
each term by $x^{1/2}\bigl(\log|d_{M_d}|+n_{M_d}\log x\bigr)\ll
x^{1/2}d^{1+o(1)}\log x$ (as
$n_{M_d}\le6^{\omega(d)}\varphi(12d)=d^{1+o(1)}$), and summing over
$d\le x^{2\theta}$ gives $x^{1/2+4\theta+o(1)}$; with the
class-proportional form of the error term the same computation gives
$\theta<\tfrac14$.  Note what GRH delivers is the hypothesis, not the
theorem it feeds: points (a) and (b) of (i) remain open under GRH ---
the per-modulus $x^{1/2}$ barrier is the same --- so the lane density
program is \emph{not} GRH-complete as of this draft.
\emph{(iii) Placement.}  $\mathrm{EQ}_3$ is a level-of-distribution
statement --- Bombieri--Vinogradov for cubic Kummer splitting, averaged
over the base --- with no parity obstruction; the nearest technology is
the average-Artin theorem of Stephens~\cite{Stephens69} and the cubic
Kummer-sum equidistribution (Kummer-on-average) results of
Heath-Brown--Patterson~\cite{HBP79}, neither directly sufficient.
Calibration through $q\le2\cdot10^5$: the lane has exactly two members
with $w=2$ ($q=9185$ and $q=117185$), and the $\omega$-stratified
frequencies of the cubic-non-residue conditions match the exact
Chebotarev densities of Definition~\ref{def:eq3}.  Chapter-and-verse
verification of the classical sieve inputs cited above (beta-sieve
limits, dimension-$\kappa$ Mertens) is pending, per the 2026-07-06
audit discipline.
\end{remark}

\subsection{Quotient descent and the rigidity criterion}\label{sec:rigidity}

Fix a prime $m\mid t$ and write $t=mh$. The ring surjection $\Z[C_t]\to\Z[C_m]$
(induced by $x\mapsto g$, $C_m=\langle g\rangle$) sends a Turyn pair $(r,s)$ to
symmetric $A,B\in\Z[C_m]$ with
\begin{equation}\label{eq:quotient}
  A^2+B^2=q\cdot1\qquad\text{in }\Z[C_m].
\end{equation}
Each coefficient of $A,B$ is a sum of $h$ (or $h-1$, at the identity of $A$)
original entries over a residue class; so all coefficients are bounded by $h$
with fixed parities. A zero solution count for~\eqref{eq:quotient} proves
nonexistence upstairs; a positive count carries no information.

Write $n=(m-1)/2$, $u_j=g^j+g^{-j}$, and $K=\Q(\zeta_m+\zeta_m^{-1})$, the real
cyclotomic field of degree $n$, with $\Z$-basis $1,u_1,\dots,u_n$ of $\Z[C_m]^+$
and $A=A_0+\sum_j A_j u_j$. Rationally $\Q[C_m]^+\cong\Q\times K$ (trivial
character and $g\mapsto\zeta_m$), and~\eqref{eq:quotient} becomes an ordinary
two-squares equation in $\Z$ together with
\begin{equation}\label{eq:Klevel}
  \alpha^2+\beta^2=q,\qquad
  \alpha=A_0+\sum_j A_j\bigl(\zeta_m^j+\zeta_m^{-j}\bigr)\in\OK,
\end{equation}
and likewise $\beta$ for $B$. Since $1,u_1,\dots,u_{n-1}$ is an integral basis
of $\OK$ and $\sum_j u_j=-1$ at $g\mapsto\zeta_m$, we have
$\alpha=(A_0-A_n)+\sum_{j<n}(A_j-A_n)u_j(\zeta_m)$, so $\alpha\in\Q$ iff
$A_1=\dots=A_n$, i.e.\ iff $A$ is constant off the identity. Call solutions with
$A,B$ both constant off the identity the \emph{constant branch}: writing
$A=a_0+s\sum_j u_j$, $B=b_0+c\sum_j u_j$, it is equivalent to a pair of Gaussian
representations $X^2+Y^2=q$, $U^2+V^2=q$ with $X=a_0+(m-1)s$, $U=a_0-s$,
$Y=b_0+(m-1)c$, $V=b_0-c$, subject to $X\equiv U$, $Y\equiv V\pmod m$ and the
parity/box conditions on $a_0,s,b_0,c$. Checking for a constant-branch witness
is thus a finite loop over the representations of $q$ as a sum of two squares.

\begin{definition}\label{def:rigid}
Let $p\nmid2m$ and let $f=f_K(p)$ be the residue degree of $p$ in $K$ (the order
of $p$ in $(\Z/m)^\times/\{\pm1\}$). Call $p$ \emph{rigid for $m$} if either
\emph{(i)} $p^f\equiv3\pmod4$ (equivalently $p\equiv3\bmod4$ and $f$ odd), or
\emph{(ii)} $f=n$ and $p\equiv1\pmod4$.
\end{definition}

\begin{theorem}[Rigidity]\label{thm:rigidity}
Let $m$ be an odd prime and $\gcd(q,2m)=1$. If every prime $p\mid q$ is rigid
for $m$, then every solution of~\eqref{eq:Klevel} has $\alpha,\beta\in\Z$;
equivalently every symmetric solution of~\eqref{eq:quotient} is constant off the
identity. No class-number or unit hypothesis is used.
\end{theorem}

\begin{proof}
Let $L=K(i)$ (abelian of degree $2n$, conductor $\mid4m$), $c$ the nontrivial
element of $\Gal(L/K)$ (complex conjugation in every embedding), and
$\gamma=\alpha+i\beta\in\OL$, so $\gamma\,c(\gamma)=q$ and $|\tau(\gamma)|^2=q$
at every archimedean $\tau$.

\emph{Step 1 ($\mu_L=\{\pm1,\pm i\}$).} If $\zeta_w\in L$ with $m\mid w$ then
$\Q(\zeta_m)\subseteq L$; both have degree $2n$, so they are equal and
$i\in\Q(\zeta_m)$, impossible for odd conductor $m$. Hence $w\mid4$, and $w=4$.

\emph{Step 2 (primes above $q$).} Each $p\mid q$ is unramified in $L$; with
$k_p=v_p(q)$, every $P\mid p$ satisfies
$v_P(\gamma)+v_{c(P)}(\gamma)=k_p$.

\emph{Step 3 (rigid (i)).} If $p^f\equiv3\pmod4$, the residue field of $K$ at
$p$ has order $\equiv3\bmod4$, so $x^2+1$ is irreducible there and every $P\mid
p$ is inert in $L/K$ ($c(P)=P$); then $2v_P(\gamma)=k_p$, $k_p$ is even, and the
$p$-part of $(\gamma)$ is $\bigl(\prod_{P\mid p}P\bigr)^{k_p/2}=(p^{k_p/2})$,
generated by a rational integer. (If $k_p$ is odd there is no solution and the
theorem holds vacuously.)

\emph{Step 4 (rigid (ii)).} If $f=n$ and $p\equiv1\pmod4$, write
$p=\pi\overline\pi$ in $\Z[i]$. As $K$ is totally real,
$\Gal(L/\Q)\cong\Gal(K/\Q)\times\Gal(\Q(i)/\Q)$ and the Frobenius of $p$ is
$(\sigma_p,1)$ of order $f=n$, so there are exactly $2$ primes of $L$ above $p$,
namely $\pi\OL,\overline\pi\OL$ (comaximal, swapped by $c$, product $p\OL$).
The $p$-part of $(\gamma)$ is $(\pi^a\overline\pi^{\,k_p-a})$, generated by a
Gaussian integer.

\emph{Step 5 (Kronecker).} Let $w\in\Z[i]$ be the product of the generators from
Steps 3--4; then $(\gamma)=(w)$ and $w\overline w=q$, so $u=\gamma/w\in\OL^\times$
has $|\tau(u)|=1$ everywhere, hence $u\in\mu_L=\{\pm1,\pm i\}$ (Step 1) and
$\gamma\in\Z[i]$. Writing $\gamma=x+iy$ ($x,y\in\Z$): if $\beta\ne y$ then
$i=(\alpha-x)/(y-\beta)\in K$, contradicting $K$ real; so $\alpha=x,\beta=y\in\Z$.
The group-ring statement follows by Section~\ref{sec:rigidity}'s dictionary
(rational $K$-components $\Leftrightarrow$ constant off the identity).
\end{proof}

The content is case (ii), where $p\equiv1\pmod4$ and every self-conjugacy tool
is silent.

\begin{remark}\label{rem:escapes}
The two rigid cases are precisely those in which the valuation vector of
$(\gamma)$ above $p$ is \emph{forced} to be that of a Gaussian integer. If
$p\equiv1\bmod4$ has $f<n$ there are $2n/f\ge4$ primes above $p$ and non-uniform
exponent vectors of non-Gaussian shape appear; if $p\equiv3\bmod4$ has $f$ even
(possible only for $m\equiv1\bmod4$) the primes of $K$ split in $L$ and the same
freedom appears. Whether such a non-Gaussian ideal solution has a generator of
absolute value $\sqrt q$ everywhere is a class-group and unit question---the
subject of \S\S\ref{sec:nonrigid}--\ref{sec:defect}.
\end{remark}

Two corollaries supply the obstruction.

\begin{corollary}[All-rigid obstruction]\label{cor:allrigid}
If every $p\mid q$ is rigid for some prime $m\mid t$ and the constant branch
of~\eqref{eq:quotient} has no witness, then no Turyn pair of order $t$ exists.
The witness check is the finite loop over Gaussian representations of $q$; no
row enumeration is involved.
\end{corollary}

\begin{corollary}[Prime order]\label{cor:primet}
Let $t=m>3$ be prime ($h=1$). The folding boxes force $A_0=0$ and all other
coefficients in $\{\pm1\}$, so a constant solution has $A=d(\sum_j u_j)$,
$d=\pm1$, and $B=b_0+c\sum_j u_j$ with $b_0,c\in\{\pm1\}$, whence $\alpha=-d$,
$\beta=b_0-c\in\{0,\pm2\}$ and $\alpha^2+\beta^2\le5<q$. Hence if every $p\mid q$
is rigid for $t$, no Turyn pair of prime order $t$ exists. At $t=3$,
$q=5=1^2+2^2$ shows the bound is sharp.
\end{corollary}

\begin{theorem}\label{thm:4373}
No Turyn pair of order $43$ or $73$ exists. Equivalently, there is no negacyclic
$W(86,85)$ or $W(146,145)$ of Turyn-pair type, and Turyn's construction has no
analogue at $q=85$ or $q=145$.
\end{theorem}

\begin{proof}
For $t=43$ ($n=21$): modulo $43$, $5^{21}\equiv-1$ and $\ord(17)=21$, so both
divisors of $q=85=5\cdot17$ have $f=21=n$ and are $\equiv1\bmod4$ (rigid (ii));
apply Corollary~\ref{cor:primet}. For $t=73$ ($n=36$): $5^{36}\equiv29^{36}
\equiv-1\pmod{73}$, so both divisors of $q=145=5\cdot29$ have $f=36=n$; apply
Corollary~\ref{cor:primet}.
\end{proof}

These are the first prime-order cases in the blind spot of
Remark~\ref{rem:blindspot}. The case $t=43$ was independently confirmed by an
exhaustive meet-in-the-middle sweep ($2^{21}\times2^{22}$ symmetric rows, zero
solutions); $t=73$ ($2^{36}$ rows per side) rests on the theorem. A scan
(Section~\ref{sec:results}) applies Corollary~\ref{cor:allrigid} over every
prime $m\mid t$ and closes $62$ of the $111$ blind-spot cases.

\begin{remark}[Consistency]\label{rem:consistency}
For prime-power $q=p^e$ the hypothesis of Corollary~\ref{cor:primet} can never
hold---rigid (i) contradicts $q\equiv1\bmod4$, and rigid (ii) forces $p$ to
generate $(\Z/t)^\times/\{\pm1\}$ with $e\ge n$, giving $q\ge3^n>2t-1$ for
$n\ge3$---exactly as required, since Turyn's construction realizes all prime
powers.
\end{remark}

\begin{remark}[Sharpness]\label{rem:sharp}
When some $p\mid q$ is nonrigid, nonconstant solutions of~\eqref{eq:Klevel}
occur and lift to nonconstant symmetric solutions of~\eqref{eq:quotient}
(coordinates $(A_0,\dots,A_n)$):
\begin{itemize}
\item $q=1105=5\cdot13\cdot17$, $m=7$ ($13\equiv-1\bmod7$ nonrigid, $f=1<3$):
  $A=(-6,3,-15,-1)$, $B=(-7,-5,-9,13)$ satisfy $A^2+B^2=(1105,0,0,0)$ ($96$
  nonconstant $K$-level pairs exist here).
\item $q=1989=3^2\cdot13\cdot17$, $m=5$ ($3\equiv3\bmod4$ has $f=2$ even, the
  split escape): $A=(-14,9,19)$, $B=(5,-19,9)$ satisfy $A^2+B^2=(1989,0,0)$.
\end{itemize}
Conversely, on cases closed by Corollary~\ref{cor:allrigid} the enumeration
finds no nonconstant $K$-level pair (checked at $q=153,265,405,545,657,685$),
while the control $q=601$ (prime) shows $24$. These certificates settle only the
quotient's constancy behaviour; they neither construct nor exclude Turyn pairs
upstairs. That lift question is the next subject.
\end{remark}

\subsection{The nonrigid case: three exact conditions}\label{sec:nonrigid}

When some $p\mid q$ is nonrigid for $m$, we give an exact criterion for a
symmetric quotient solution, replacing the earlier ad hoc layer terminology of
the working notes by the standard trichotomy for a norm-form realization
problem: \emph{existence} of a solution to the relevant norm equation, the
\emph{congruences} it must satisfy, and the \emph{size} (coefficient bound) it
must respect.

Keep $L=K_m(i)$, $c$ complex conjugation, and $p_m$ the unique prime of $\OK$
above $m$ ($\OK/p_m=\F_m$, $u_j\mapsto2$). Set the congruence modulus
\[
  M_B=2\,p_m\OL .
\]
For $\theta=c_0+c_1u_1+\cdots+c_{n-1}u_{n-1}\in\OK$ and $X\in\Z$ define
$\Recon_{m,h}(X,\theta;\epsilon)$, $\epsilon\in\{R,S\}$, by
\[
  b(\theta)=c_0+2(c_1+\cdots+c_{n-1}),\quad w=\tfrac{X-b(\theta)}{m},\quad
  \Recon=(c_0+w,c_1+w,\dots,c_{n-1}+w,w),
\]
defined only when $w\in\Z$ and the row satisfies the Turyn boxes (off-identity
entries odd, $|\cdot|\le h$; the $R$ identity entry with zero-diagonal parity
and $|a_0|\le h-1$; the $S$ identity entry odd, $|a_0|\le h$). Write
$\mu_4=\{\pm1,\pm i\}$, $\Gamma(q)=\{\gamma\in\OL:\gamma\,c(\gamma)=q\}/\mu_4$,
and $\Sq$ for the compressed square signature in $\Z[C_m]^+$. Say an ideal $I$
is \emph{unit-admissible} if $I\,c(I)=(q)$ and $I=(\delta)$ has a generator with
$\delta\,c(\delta)=q$ (equivalently $\delta c(\delta)/q\in\OK^\times$ is a
relative unit norm from $\OL^\times$). Define the source/target ray sets
\[
  S_B(m,q)=\{\mu\,\sigma(\gamma)\bmod M_B:\gamma\in\Gamma(q),\mu\in\mu_4,
    \sigma\in\{1,c\}\},
\]
\[
  T_B(m,q)=\{r\bmod M_B:\ r\equiv1\ (\mathrm{mod}\ 2\OL),\
    r\equiv x+iy\ (\mathrm{mod}\ p_m),\ x^2+y^2=q,\ x\ \text{even},\ y\ \text{odd}\}.
\]

\begin{lemma}[Reconstruction]\label{lem:recon}
Under $\Q[C_m]^+\cong\Q\times K$ a symmetric row maps to $(X,\theta)$; it is
integral iff $X\equiv\theta\pmod{p_m}$, and is then the unique row
$\Recon_{m,h}(X,\theta;\epsilon)$, with identity coefficient
$a_0=(X+2\Tr_{K/\Q}(\theta))/m$.
\end{lemma}

\begin{proof}
Inverse DFT gives the displayed row; integrality is $m\mid X-b(\theta)$, exactly
$X\equiv\theta\pmod{p_m}$ since $u_j\mapsto2$. Using $\Tr(1)=n$, $\Tr(u_j)=-1$,
$m=2n+1$: $X+2\Tr(\theta)=X+2(nc_0-c_1-\cdots-c_{n-1})=m(c_0+w)$.
\end{proof}

\begin{lemma}[Ideal/unit layer]\label{lem:idealunit}
Solutions of $\alpha^2+\beta^2=q$ in $\OK$ are exactly $\gamma=\alpha+i\beta\in
\OL$ with $\gamma\,c(\gamma)=q$; their ideals satisfy $I\,c(I)=(q)$, and such an
ideal contributes a norm-$q$ generator iff it is unit-admissible. Two such
generators differ by an element of $\mu_4$.
\end{lemma}

\begin{proof}
$(\alpha+i\beta)c(\alpha+i\beta)=\alpha^2+\beta^2$ gives the first claim. If
$I=(\delta)$ with $Ic(I)=(q)$ then $\delta c(\delta)=\varepsilon q$,
$\varepsilon\in\OK^\times$; $u\delta$ has norm $q$ iff $uc(u)=\varepsilon^{-1}$,
i.e.\ unit-admissibility. Two norm-$q$ generators differ by a relative
norm-one unit; such a unit has all absolute values $1$, so by Kronecker is a
root of unity, and $\mu_L=\mu_4$ (Step 1 of Theorem~\ref{thm:rigidity}).
\end{proof}

\begin{lemma}[Parity and ray gluing]\label{lem:raygluing}
For a Turyn-compatible pair the components satisfy $\theta_R\equiv1$,
$\theta_S\equiv0\pmod{2\OK}$; the trivial values have $X$ even, $Y$ odd; and
$\theta_R\equiv X$, $\theta_S\equiv Y\pmod{p_m}$. Hence $\gamma=\theta_R+i\theta_S
\in T_B(m,q)\bmod M_B$, and conversely membership in $T_B$ yields exactly these
conditions.
\end{lemma}

\begin{proof}
The parities are the compressed-coordinate form of the Turyn boxes (for $R$ the
zero-diagonal makes the constant coordinate odd and the rest even; for $S$ all
basis coordinates even); summing gives the parity of $X,Y$.
Lemma~\ref{lem:recon} gives the $p_m$ gluing. As $2\OL$ and $p_m\OL$ are
coprime, CRT combines them into the single modulus $M_B$.
\end{proof}

We may now state the three conditions and the criterion.
\begin{description}
\item[\normalfont(N) Norm equation.] There is a Weil generator
  $\gamma\in\Gamma(q)$ (equivalently, a unit-admissible ideal).
\item[\normalfont(C) Congruence.] $S_B(m,q)\cap T_B(m,q)\ne\varnothing$: some
  oriented conjugate of $\gamma$ meets the target ray set modulo $M_B$.
\item[\normalfont(R) Realization.] The rows $\Recon_{m,h}(x,\alpha;R)$,
  $\Recon_{m,h}(y,\beta;S)$ are defined inside their coefficient boxes.
\end{description}

\begin{theorem}[Uniform realization criterion]\label{thm:realization}
Fix a congruence-admissible generator $\gamma=\alpha+i\beta$ and a signed
two-square pair $(x,y)$, so that $\alpha^2+\beta^2=q$ in $\OK$ and
$x^2+y^2=q$ in $\Z$. Then condition \emph{(R)} holds if and only if the two
rows
\[
  R=\Recon_{m,h}(x,\alpha;R),\qquad
  S=\Recon_{m,h}(y,\beta;S)
\]
are defined. When they are defined,
\[
  \Sq(R)+\Sq(S)=(q,0,\dots,0)
\]
automatically.

More explicitly, for $\theta=c_0+\cdots+c_{n-1}u_{n-1}$ and
$X\in\Z$, put $w=(X-b(\theta))/m$ and
$a_0=c_0+w$. The row $\Recon_{m,h}(X,\theta;\epsilon)$ is defined exactly when
$w\in\Z$, all tail entries
\[
  c_1+w,\dots,c_{n-1}+w,w
\]
are odd and have absolute value at most $h$, and the head entry satisfies the
corresponding identity box:
\[
  \epsilon=R:\ |a_0|\le h-1,\ a_0\equiv h-1\pmod 4;
  \qquad
  \epsilon=S:\ |a_0|\le h,\ a_0\equiv1\pmod 2.
\]
\end{theorem}

\begin{proof}
The explicit head and tail tests are just the definition of
$\Recon_{m,h}$, with Lemma~\ref{lem:recon} identifying the head coefficient as
$a_0=(X+2\Tr_{K/\Q}(\theta))/m$. Thus the two reconstructed rows are unique
whenever they exist.

Let $\Phi:\Q[C_m]^+\to\Q\times K$ be the usual product map, sending a row to its
trivial value and its $K$-component. This map is injective. If the displayed
rows are defined, then
\[
  \Phi(R^2+S^2)=(x^2+y^2,\alpha^2+\beta^2)=(q,q)=\Phi(q),
\]
so $R^2+S^2=q$ in $\Q[C_m]^+$, hence in $\Z[C_m]^+$ because the rows are
integral. The compressed square-signature equation is this equality in
coordinates. Conversely, any Turyn-compatible rows with the fixed data must be
the two reconstructions by Lemma~\ref{lem:recon}; hence realization for that
data is precisely their definedness.
\end{proof}

\begin{theorem}[Prime-quotient descent]\label{thm:descent}
Turyn-compatible symmetric rows $R,S\in\Z[C_m]^+$ with $R^2+S^2=q$ exist if and
only if conditions \emph{(N)}, \emph{(C)}, \emph{(R)} hold simultaneously for
some $\gamma=\alpha+i\beta\in\Gamma(q)$, $\mu_4$-orientation, optional
conjugation, and signed two-square pair $(x,y)$. If \emph{(C)} fails the
quotient is closed with no coordinate enumeration; if \emph{(C)} holds the
question reduces to the finite check \emph{(R)}.
\end{theorem}

\begin{proof}
Given rows $R,S$ with trivial values $X,Y$ and components $\theta_R,\theta_S$,
set $\gamma=\theta_R+i\theta_S$; Lemma~\ref{lem:idealunit} gives
$\gamma\in\Gamma(q)$, Lemma~\ref{lem:raygluing} puts its residue in $T_B$, and
Lemma~\ref{lem:recon} identifies the rows as the stated reconstructions with
$\Sq(R)+\Sq(S)=(q,0,\dots,0)$; the conditions are necessary. Conversely, given
the data, the ray condition supplies the parity and gluing
(Lemma~\ref{lem:raygluing}), and condition (R) says the two unique
reconstructions are inside the boxes. Theorem~\ref{thm:realization} then gives
the exact square signature, so these are Turyn-compatible rows with
$R^2+S^2=q$.
\end{proof}

Two examples calibrate the conditions.
\begin{itemize}
\item $q=1377=3^4\cdot17$, $m=13$: condition (C) fails---no unit-admissible
  generator meets $T_B(13,1377)$, so the $C_{13}$ quotient is closed. This
  closure is unconditional ($\mathrm{bnfcertify}$ certifies the class data of
  $K_{13}(i)$) and independent of coefficient enumeration.
\item $q=185=5\cdot37$, $m=31$, $h=3$: condition (C) \emph{passes} but (R)
  fails. All $24$ unit-admissible generators reconstruct the $R$-row inside the
  box, but the $S$-row identity coefficient is forced to $|a_0|=5>h=3$ (the
  surviving two-square eigenvalue is $y\in\{\pm13\}$, giving
  $a_0(S)=(y+2\Tr(\theta_S))/31=\pm5$). Thus (R) is a genuinely separate
  obstruction, a thin-fiber ($h=3$) phenomenon at the trivial character. This
  quotient-lane certificate depends on the $m=31$ class list; the unconditional
  closure of $q=185$ used in Section~\ref{sec:results} comes instead from the
  field-free reduced decision of Section~\ref{sec:t3-decision}.
\end{itemize}

On the field-tractable subpanel where (C) passes, the uniform realization
criterion gives the following compact certificate. Rows marked ``yes'' have
$\mathrm{bnfcertify}$-verified class data; rows marked ``cond.'' are not used as
certified closures. The rows are emitted by
\texttt{condition\_r\_realization.py} (a diagnostic-table generator
retained in the companion source repository, not among the shipped
\texttt{code/} scripts --- it pulls in the \textsc{pari/gp}-driven
layer-B subtree) and are provided as the machine-readable cache
\texttt{code/data/condition\_r\_realization\_qmax2000\_degree40.jsonl}.

\begin{center}
\small
\begin{tabular}{r r r l r r r l l}
$m$ & $q$ & $h$ & (R) & lifts & $\min|a_0(R)|$ & $\min|a_0(S)|$ & cert. & block\\
\hline
31 & 185  & 3   & WIPEOUT & 0   & 2  & 5  & cond. & $S$-head\\
5  & 369  & 37  & PASS    & 64  & 4  & 3  & yes   & lift\\
13 & 441  & 17  & PASS    & 8   & 0  & 15 & yes   & lift\\
17 & 441  & 13  & PASS    & 8   & 0  & 5  & yes   & lift\\
5  & 549  & 55  & PASS    & 64  & 2  & 3  & yes   & lift\\
5  & 1089 & 109 & PASS    & 8   & 0  & 11 & yes   & lift\\
7  & 1469 & 105 & PASS    & 384 & 4  & 1  & yes   & lift\\
13 & 1585 & 61  & PASS    & 48  & 16 & 5  & yes   & lift\\
\end{tabular}
\end{center}

\begin{remark}[The identity-coefficient bound]\label{rem:box}
The (R) failure at $q=185$ is forced, not an artifact. By Lemma~\ref{lem:recon}
and $|\sigma(\theta)|\le\sqrt q$ at each embedding, $|a_0|\le(\sqrt q+2n\sqrt
q)/m=\sqrt q$; on the $S$-side $a_0$ is a nonzero odd integer $\le\sqrt q$ in
absolute value. Moreover a general constant-off-identity solution is excluded
whenever $2\cdot(\text{largest two-square component of }q)<m$ and $m>h$: then
$X\equiv U\pmod m$ with $|X-U|<m$ forces $s=0$ (likewise $t=0$), collapsing the
row to a scalar $a_0^2+b_0^2=q$ that cannot fit $|a_0|,|b_0|\le h$ since
$q=2mh-1>2h^2$. For $q=185$, $m=31$: $2\cdot13<31$, so the constant branch is
empty and only the nonconstant $|a_0(S)|=5$ obstruction remains.
\end{remark}

\subsection{Where the remaining obstruction lives}\label{sec:defect}

Condition (C) is the only step whose evaluation, as stated, requires
constructing $K_m(i)$: the target $T_B$ is local (two-square residues and
parities), but the source residues $\gamma\bmod M_B$ are those of the Weil
generators. Whether they can be read from local splitting data is governed by
the class group.

Let $C_{\mathrm{loc}}(m,q)$ be the finite set of ideals $I$ with $Ic(I)=(q)$
built from local prime data: for each $P\mid p$ ($p\mid q$), if $P=c(P)$ the
valuation is forced to $e_p/2$ (existing only for $e_p$ even), and if $P\ne c(P)$
choose exponents $a$ and $e_p-a$ on the pair. Let $\Cl_B(L)$ be the ray class
group modulo $M_B$.

\begin{definition}[Ray defect]\label{def:defect}
$D_B(m,q)=\{\operatorname{cl}_B(I):I\in C_{\mathrm{loc}}(m,q),\ I\text{
unit-admissible}\}\subseteq\Cl_B(L)$.
\end{definition}

\begin{lemma}[Local factor lemma]\label{lem:localfactor}
The ideals $I$ with $Ic(I)=(q)$ and $(I,M_B)=1$ are exactly
$C_{\mathrm{loc}}(m,q)$.
\end{lemma}

\begin{proof}
Since $(q,2m)=1$, every prime dividing $I$ lies above a rational $p\mid q$ and is
prime to $M_B$; $v_P(I)+v_{c(P)}(I)=e_p$ forces $2v_P(I)=e_p$ when $P=c(P)$ and
allows independent choices on distinct conjugate pairs otherwise.
\end{proof}

\begin{proposition}[Ray criterion up to the defect]\label{prop:fieldfree}
If condition (C) holds for $(m,q)$, then there are $I\in C_{\mathrm{loc}}(m,q)$,
$r\in T_B(m,q)$, $\mu\in\mu_4$, $\sigma\in\{1,c\}$ with $I$ unit-admissible and
$\operatorname{cl}_B(\sigma I)=[\mu^{-1}r]_B$ in $\Cl_B(L)$. The converse is not
immediate: the ray-class equality only records that $\sigma I$ has \emph{some}
generator congruent to $\mu^{-1}r$, whereas (C) requires this specifically of a
generator of norm $q$, and by Lemma~\ref{lem:idealunit} such generators form only
a $\mu_4$-orbit inside the full (unit-blind) generator set recorded by
$\operatorname{cl}_B$. When $K_m$ has units of infinite order this orbit can be
a proper subset, so the displayed ray-class data is necessary for (C) but not
in general known to be sufficient.
\end{proposition}

\begin{proof}
By Lemma~\ref{lem:localfactor} the ideal of any norm-$q$ generator lies in
$C_{\mathrm{loc}}(m,q)$; Theorem~\ref{thm:descent} identifies (C) with
$\mu\sigma(\gamma)\equiv r\pmod{M_B}$ for some Weil generator $\gamma$, which in
particular exhibits a generator of $\sigma I$ congruent to $\mu^{-1}r$, i.e.\ the
displayed ray-class equality.
\end{proof}

\begin{corollary}[Trivial-defect product formula]\label{cor:trivialdefect}
If the relevant unit-admissible local factors are principal with known residues
modulo $M_B$, then $S_B(m,q)=\mu_4\{1,c\}\cdot(\text{forced Gaussian residue})
\cdot\prod_{\text{split pairs}}(\text{local prime-pair residues})$, and (C) holds
iff this product meets $T_B(m,q)$.
\end{corollary}

The Corollary's hypothesis is exactly principality of the primes above $q$, and
it fails often.

\begin{proposition}\label{prop:imq}
Let $\mathrm{Im}_q=\langle[P]:P\mid q\rangle\subseteq\Cl(L)$. The source residues
in \emph{(C)} are a product of explicit local (principal-prime-factor) residues
---computable without building $K_m(i)$---if and only if $\mathrm{Im}_q=0$. If
$\mathrm{Im}_q\ne0$, a Weil generator is produced only through the global class
group, and $D_B(m,q)$ is a genuine, nonzero global datum.
\end{proposition}

On the field-tractable panel of nonrigid survivors ($q\le2000$, field degree
$\le40$), $\mathrm{Im}_q\ne0$ for half the cases, including the frontier
$q=185$, where the offending classes have odd order $41$ in $\Cl(K_{31}(i))$
($h_L=5084$). Moreover quadratic genus theory does not decide it: the $2$-part
obstructions lie in $\Cl(L)^2$, below the genus quotient $\Cl(L)/2\Cl(L)$
(at $m=17$ the class is $4=2\cdot(\Z/8)$; at $m=41$ it has order $4$), and are
invisible to genus characters; individual primes above $q$ can even be
genus-visible on the sum-of-two-squares domain (for $m=17$ the first case is
$q=13837=101\cdot137$), but in a unit-admissible (principal) ideal the genus
images cancel. Hence (C) is not, in general, a local test; the residual global
input is the odd part of $\mathrm{Im}_q$ (no reciprocity handle) together with a
$2$-part accessible at best to R\'edei/higher reciprocity. This is a negative
structural result: it locates the obstruction precisely rather than removing it.
A further reduction --- quotienting the ray class group by the field-constant
generator ambiguities (roots of unity, local norm-one residues, global relative
units) to isolate a finite ``absorbed'' image --- was explored in an earlier
draft; the reverse direction of that reduction (that every class in the
quotient is witnessed by an actual Weil generator) turned out not to be
established, so it is omitted here. None of the closures in
Section~\ref{sec:results} depend on it.

\section{Closed orders}\label{sec:results}

We summarize the composite $q\le2000$ closed by the hierarchy. Counts in the
self-conjugacy-blind population are not additive: the rigidity and T3 reduced
decisions overlap.

\begin{center}
\small
\begin{tabular}{p{0.22\textwidth} p{0.34\textwidth} p{0.34\textwidth}}
\hline
Layer & Mechanism & Effect for $q\le2000$\\
\hline
Two squares (\S\ref{sec:twosquares}) & $\ell\equiv3\bmod4$ to odd power &
$211$ closed\\
Self-conjugacy (\S\ref{sec:selfconj}) & inert tower, Prop.~\ref{prop:selfconj} &
$13$ further closed\\
Rigidity scan (\S\ref{sec:rigidity}) & prime-quotient rigidity &
$62$ of the $111$ blind-spot cases\\
T3 reduced decision (\S\ref{sec:t3-decision}) & composite multipliers + MITM &
$85$ of the $111$ blind-spot cases\\
Rigidity $\cup$ T3 & combined unconditional coverage &
$95$ of the $111$ blind-spot cases\\
Marginal decision (\S\ref{sec:t2-marginal}) & composite/full sub-tori &
closes all $16$ remaining cases\\
Full hierarchy & all field-free layers above &
every blind-spot multi-prime case $q\le2000$ closed\\
Nonrigid descent (\S\ref{sec:nonrigid}) & (N)/(C)/(R), field-tractable panel &
structural diagnostics and conditional local certificates\\
\hline
\end{tabular}
\end{center}

The rigidity scan reproduces without enumeration the five cases previously
closed by exact quotient computation ($q=65,153,265,405,657$) and adds $57$
(many at prime $t$ via Corollary~\ref{cor:primet}), among them
\[
  q=85,\ 145,\ 205,\ 445,\ 481,\ 565,\ 745,\ 793,\ 865,\ 925,\ 1261,\ 1345,\
  1537,\ 1717,\ 1813,\ \dots
\]
The T3 reduced decision closes $85$ blind-spot cases. It overlaps the rigidity
scan in $52$ cases and adds the following $33$ unconditional closures:
\[
\begin{gathered}
185,\ 225,\ 245,\ 325,\ 365,\ 369,\ 441,\ 485,\ 585,\ 605,\ 685,\ 697,\ 725,\\
801,\ 985,\ 1017,\ 1089,\ 1145,\ 1157,\ 1233,\ 1285,\ 1305,\ 1377,\ 1417,\\
1465,\ 1513,\ 1521,\ 1525,\ 1665,\ 1685,\ 1765,\ 1885,\ 1989 .
\end{gathered}
\]
This list includes $q=185$, formerly the first frontier for the nonrigid
quotient lane; $q=441$, where both prime quotients pass and the obstruction is
invisible to the quotient criteria; and the all-rigid-with-witness cases
$q=485,685,725$, where quotient descent is permanently silent.

After combining the unconditional rigidity scan with the original T3 harvest,
there were $16$ residual self-conjugacy-blind cases through $2000$. The marginal
decision closes all of them:
\[
\begin{gathered}
377,\ 425,\ 545,\ 549,\ 629,\ 1025,\ 1105,\ 1205,\ 1325,\ 1445,\ 1469,\\
1565,\ 1585,\ 1625,\ 1649,\ 1937 .
\end{gathered}
\]
Here $1205$ and $1565$ are full-torus T3 decisions; $1585$ is killed already at
the $t'=61$ marginal (and also at the full torus by the CP-SAT formulation);
$1325$, the final residual, is killed by the quartic-component criterion at
$t'=51$ (Proposition~\ref{prop:q1325-quartic}); the others are killed at proper
composite sub-tori. Thus the current field-free hierarchy closes every
self-conjugacy-blind multi-prime composite $q\le2000$.
The field-tractable (N)/(C)/(R) tables above remain useful for locating quotient
phenomena, but the T3/T2 field-free layers supersede several old conditional
closure claims in the headline harvest.

\begin{center}
\small
\begin{tabular}{r r r l}
$q$ & killing $t'$ & $h=t/t'$ & certificate\\
\hline
377 & 27 & 7 & $C_{27}$ boxed congruence criterion (Cor.~\ref{cor:c27-arithmetic-firing})\\
425 & 71 & 3 & $h=3$ support-augmentation obstruction (Prop.~\ref{prop:h3-support-augmentation})\\
545 & 21 & 13 & $C_{21}$ component join criterion (Cor.~\ref{cor:c21-arithmetic-firing})\\
549 & 25 & 11 & $p=5$ component criterion (Prop.~\ref{prop:p5-sqrt5})\\
629 & 21 & 15 & $C_{21}$ component join criterion (Cor.~\ref{cor:c21-arithmetic-firing})\\
1025 & 27 & 19 & $C_{27}$ residue-pair criterion (Cor.~\ref{cor:c27-residue-pair})\\
1105 & 79 & 7 & prime sub-torus period criterion (Cor.~\ref{cor:prime-subtorus-firing}); also HiGHS MILP\\
1205 & 603 & 1 & T3 MITM\\
1325 & 51 & 13 & quartic component criterion (Prop.~\ref{prop:q1325-quartic})\\
1445 & 241 & 3 & $h=3$ support-PAF obstruction (Prop.~\ref{prop:q1445-h3-paf})\\
1469 & 49 & 15 & $p=7$ cubic component criterion (Prop.~\ref{prop:q1469-p7-cubic})\\
1565 & 783 & 1 & T3 MITM\\
1585 & 61 & 13 & prime sub-torus period criterion (Cor.~\ref{cor:prime-subtorus-firing}); also exhaustive MITM and HiGHS MILP\\
1625 & 271 & 3 & $h=3$ support-sum obstruction (Prop.~\ref{prop:h3-support-augmentation})\\
1649 & 25 & 33 & $p=5$ component criterion; also exhaustive MITM at $t'=55$\\
1937 & 57 & 17 & $C_{57}$ quadratic-component criterion (Prop.~\ref{prop:c57-sqrt57}); also exhaustive MITM and MILP at $t'=323$\\
\end{tabular}
\end{center}

\begin{remark}[What ``closed'' asserts, restated]\label{rem:closed-restated}
Each entry above is a proof that no Turyn \emph{pair} of the corresponding order
$t$ exists---equivalently, no negacyclic $W(2t,2t-1)$ of Turyn-pair type, no
binary almost-perfect sequence of period $4t$ of the corresponding kind. It is
\emph{not} a proof that no Hadamard matrix of order $4t$ exists. For $q$ within
the negacyclic $v<1000$ classification the nonexistence is a known data point
and the contribution is the proof; the criterion also closes cases (e.g.\ larger
prime $t$ via Corollary~\ref{cor:primet}) at orders beyond exhaustive reach.
\end{remark}

\section{Related work}\label{sec:related}

\emph{The Turyn lineage.} Turyn~\cite{Turyn72}; Whiteman's cyclotomic
re-proof~\cite{Whiteman73}; Yamamoto--Yamada's relative Gauss sums~\cite{YY85};
Lang--Schneider's uniqueness up to equivalence through order $99$~\cite{LS12};
the recent Williamson-type searches of
Kharaghani--Mohammadian--Tayfeh-Rezaie~\cite{KMT26}.

\emph{Equivalent objects.} Negacyclic conference matrices and the $v<1000$
classification~\cite{DGS71,Seberry16}; almost-perfect autocorrelation sequences
and their cyclic-divisible-difference-set description~\cite{Wolfmann92,PB95};
the two-circulant-core construction~\cite{BAOD22}; the classical
$\mathrm{BGW}(q+1,q,q-1)$ from Bose's affine difference set and Jungnickel's
equivalence~\cite{Bose42,Jungnickel82,JK}; almost-perfect polyphase
sequences~\cite{Luke97,Luke01,ZHL05}.

\emph{Method.} Rigid case (i) is Turyn's self-conjugacy mechanism~\cite{Turyn65}
in relative form; no novelty is claimed for it. Rigid case (ii) is the
composite-conductor analogue of the order-driven field descent of Leung and
Schmidt---the anti-field-descent method~\cite[Lemma 4.6, Thm 4.5]{LS16} and its
nonsquare refinement~\cite[Thm 22]{LS19}---which are proved for odd prime-power
conductor $p^a$; the field of relevance here, $L=K_m(i)$, has even composite
conductor $4m$, and Schmidt's composite-conductor descent~\cite{Schmidt99}
cannot reach the conclusion (its $F(M,N)$ is a multiple of $\mathrm{rad}(M)$, so
it never strips $m$; concretely $F(172,85)=172$, no descent, whereas
Theorem~\ref{thm:rigidity} lands solutions in $\Z[i]$). The contribution is the
exact criterion (Definition~\ref{def:rigid}), the corollaries it feeds, and the
sharpness certificates that bound it, not a new descent method. The multiplier
theorems are likewise normal-form refinements of classical multiplier theory for
affine difference sets~\cite{JK}. The composite group
$M(q)=\cap_{p\mid q}\langle p\rangle$ is the natural common-decomposition
subgroup; before any novelty claim it should be checked against the
divisible-difference-set multiplier literature. The marginal obstruction is a
finite quotient consequence of the same multiplier theorem. No such audit is
attempted here. The unconditional closures from
Sections~\ref{sec:t3-decision}--\ref{sec:t2-marginal} do not depend on that
novelty posture.

\emph{Honesty on the closed facts.} As facts, the nonexistence of Turyn pairs of
orders $43,73$ (periods $172,292$), and also $q=185,441$, lie within the
DGS/Eades computational range for negacyclic conference matrices. The
almost-perfect-sequence papers of Wolfmann~\cite{Wolfmann92} and
Bradley--Pott~\cite{PB95} treat these periods, and a divisible-difference-set
multiplier argument there could conceivably cover some of the same closures.
Their full texts were not obtained; even so, for the all-$1\bmod4$ cases any
such argument would rest on multipliers and counting, not self-conjugacy, and
would differ from the quotient-descent mechanism here.

\emph{Multiplier and RDS shadows (2026-07-06 audit).}  The composite
multiplier theorem should be read as a Turyn-normal-form specialization
of classical multiplier theory.  In the group-developed
relative-difference-set setting, multiplier lifting and fixed-translate
theorems are standard (the Pott/Lam and Arasu--Jungnickel--Ma--Pott
lineage, as summarized in Gordon's cyclic-RDS survey), and circulant
weighing-matrix searches use the same multiplier-orbit and
intersection-number constraints (Arasu--Gordon--Zhang).  What is used
here is the common-decomposition subgroup $M(q)$ acting on an arbitrary
Turyn pair, where the zero diagonal removes the translate and the
$R+iS$ normal form pins the signs.  Likewise, the quotient maps and
Gaussian-period component algebras throughout are classical; the
novelty claimed for Theorem~\ref{thm:divisor-selection} is the
marginal-complete organization for Turyn pairs --- the exact top node,
monotone emptiness on the divisor lattice, and finite per-family firing
criteria attached to specified nodes --- not the quotienting of group
rings or the periods themselves.  The fold layer of
Proposition~\ref{prop:fold-dichotomy} uses standard ideal-class
obstructions in CM fields, in the same broad family as self-conjugacy
and Leung--Schmidt field descent; its contribution is the stabilizer
amplification interface, and any closure using it retains the
certification tier of the class-group computation.

\emph{Scope of the ``first determination'' claims.}  The closures at
$q=1469$, $1937$, $1325$ (and, conditionally, $5185$, $62305$) are
apparently first determinations in the negacyclic/Turyn-pair literature
searched in the 2026-07-06 audit: they lie beyond the
DGS/Eades--Seberry $v<1000$ negacyclic classification and the
Lang--Schneider range, and exact-order searches (including the
associated Hadamard/TCC orders and APS periods) found no hit in
accessible sources.  This is not a statement about arbitrary Hadamard
matrices, and it remains subject to the closed-table caveat above.  The
semiprime-lane theorem is new only in the same audited sense: its proof
combines classical shadows --- common-decomposition multipliers and
finite quotient marginals --- with the lane-specific anchor and support
congruence, and no prior statement excluding all semiprime $q=6\ell-1$,
$\ell\equiv3\pmod4$, was found.

\section{What remains}\label{sec:remains}

Conjecture~\ref{conj:converse} is open. The hierarchy above closes every
self-conjugacy-blind multi-prime composite $q\le2000$, but the proof status is
not simply ``the finite search got bigger.''  The present state separates into
three parts: the prime-quotient conceptual layers that are now exact criteria,
the finite-panel branch accounting for the field-free multiplier/marginal
layers, and the uniform theorem obligations still needed for a full converse.

\begin{conjecture}[Marginal-complete obstruction program]\label{conj:program}
Let $q=2t-1$ be composite. Then at least one holds: \emph{(1)} a classical
two-squares or self-conjugacy obstruction fires; \emph{(2)} the full T3
multiplier-reduced system is infeasible; \emph{(3)} for some proper divisor
$t'\mid t$, the $V'$-invariant marginal system of
Proposition~\ref{prop:t2-marginal} is infeasible; \emph{(4)} all such finite
marginals pass, but their data do not glue to one element of $\Z[C_t]$.
\end{conjecture}

\subsection{Closed conceptual layers}\label{sec:closed-conceptual-layers}

The prime-quotient descent clarifies one layer and closes a second. Condition
(C) is not yet an exact criterion in general: Proposition~\ref{prop:imq} shows
that $\mathrm{Im}_q\subseteq\Cl(K_m(i))$ is essential whenever it is nonzero,
and Proposition~\ref{prop:fieldfree} reduces (C) to a necessary ray-class
congruence, but deciding (C) from that data alone remains open when $K_m$ has
units of infinite order. Condition (R), the realization criterion, is exact:
after (N) and (C), a candidate realizes precisely when the two reconstruction
rows lie in their coefficient boxes. The $q=185$ quotient certificate remains
the model local wipeout---(C) passes, but every survivor overflows the
$S$-head box by Remark~\ref{rem:box}---although the unconditional closure of
$q=185$ used in the headline harvest is the field-free T3 decision.

The field-free multiplier layer is likewise no longer an all-plus-only theorem
framework.  The T3 solver still refuses sign-alternating inputs as an
implementation safety rule, but Proposition~\ref{prop:signed-marginal} gives
the signed marginal system whenever an $\epsilon=-1$ branch occurs.
Lemma~\ref{lem:signed-branch-reduction} reduces any post-two-squares signed
case to the square branch $q=n^2$ with all primes of $n$ congruent to
$3\bmod4$.  The first multi-prime classical survivors in the signed audit,
through $q\le50000$, are $29241$, $40401$, and $43681$; all are closed by
Corollaries~\ref{cor:signed-full-torus-sign} and
\ref{cor:signed-full-unit-prime}.  Thus the sign guard is not a foundational
gap.  A full converse must still route any future signed square branch into a
signed firing criterion, just as the all-plus branch must be routed into an
all-plus firing criterion.

\subsection{Finite-panel branch accounting}\label{sec:finite-panel-accounting}

The field-free harvest should be read at the level of branches, not merely as a
list of closed examples.  After the two-squares and self-conjugacy filters,
there are $111$ self-conjugacy-blind multi-prime composites $q\le2000$.  The
rigidity scan closes $62$ of them, T3 closes $85$, their union closes $95$, and
the marginal layer closes the remaining $16$:
\[
\begin{gathered}
377,\ 425,\ 545,\ 549,\ 629,\ 1025,\ 1105,\ 1205,\ 1325,\ 1445,\ 1469,\\
1565,\ 1585,\ 1625,\ 1649,\ 1937 .
\end{gathered}
\]
The first-determination cases $q=1469$, $q=1937$, and $q=1325$ all have
$2t\ge1000$, lie outside the negacyclic classification range, and pass every
prime-quotient test.  The last of them, $q=1325$ with
$t=663=3\cdot13\cdot17$, has solvable prime marginals $t'=3,13,17$; it is
closed only at the $t'=51$ marginal by the quartic-component criterion
(Propositions~\ref{prop:c51-quartic}--\ref{prop:q1325-quartic}).

The executable scaffold for the existential divisor theorem is
\[
\texttt{orbit\_signature\_scan.py}\quad
\texttt{--existential-divisor-report}\quad
\texttt{--summary-only}.
\]
Unlike the older lane census, it evaluates actual firing predicates on every
proper divisor.  By Theorem~\ref{thm:divisor-selection}(i), each closure in
this accounting records one certified minimal element of the empty-node
up-set $U(q)$, and the branch label records the certifying criterion; the
census formula of Lemma~\ref{lem:orbit-census} is machine-verified against
the constructed orbit algebras by
\texttt{--divisor-lattice-audit}.  With the two-orbit, $h=3$ support-augmentation, $p=3$, $p=5$, $p=7$,
$C_{27}$, $C_{21}$, $C_{51}$, prime sub-torus, full-torus sign, and secondary
exact-marginal predicates in place, the $q\le2000$ branch decomposition is:
\[
\begin{array}{rl}
81 & \text{proper-divisor quantified fires},\\
29 & \text{exact full-torus sign fires},\\
1  & \text{projection-image gluing obstruction }(q=441).
\end{array}
\]
Thus neither a finite-certificate branch nor a secondary-certificate
branch remains in this panel: the last four secondary orders were
promoted to written criteria --- $685$ and $1665$ by the full-unit
$C_{49}$ instance of Proposition~\ref{prop:pe-full-unit-nested}, $1937$
by the $C_{57}$ criterion (Proposition~\ref{prop:c57-sqrt57}), and
$1885$ by the degree-$4$ period join of
Corollary~\ref{cor:prime-subtorus-firing} under the measured box gate
--- and every closure in the $21$-order T3-unreachable panel is a
quantified firing route.  The old proper-divisor
component branch is consumed by the prime sub-torus join and the $q=441$
gluing check.  The old prime-order full-torus component branch is consumed
by the row-sum bound (Theorem~\ref{thm:full-torus-rowsum}): all $29$
full-torus branch members, including the higher-degree cases $q=697$,
$1521$, and $1765$ formerly closed by complete integer MITM, satisfy
$w\ge d+2$ and fire with zero computation.  Every \texttt{scalar\_unsat} row is an instance of the
two-orbit representation criterion (Theorem~\ref{thm:two-orbit-reps}) and
its corollaries, certified by closed-form arithmetic rather than
enumeration; Remark~\ref{rem:two-orbit-arithmetic} records the node-level
audit.

\paragraph{Cross-prime gluing.}\label{sec:gluing}
Passing (N),(C),(R) at every prime quotient does \emph{not} produce a global
element of $\Z[C_t]^+$; the quotient data must be compatible under the common
projections and the integral boxes.  At $q=441=21^2$,
$t=221=13\cdot17$, both prime quotients admit valid nonrigid symmetric rows
($C_{13}$ and $C_{17}$ each lift; see the table in \S\ref{sec:defect}), so the
prime-quotient criteria give no obstruction.  The exact side-set gluing check
does: the $C_{13}$ and $C_{17}$ quotient side-solution sets each contain two
$A$-rows and two $B$-rows, but no boxed $V$-orbit row on $C_{221}$ projects to
the corresponding side sets.  Already in the $C_{13}$ projection, the nonzero
quotient coordinates of a cover row have the form
$4(x_1+x_3+x_6+x_7)+x_9$ and $4(x_2+x_4+x_5+x_{10})+x_{12}$, hence are
$8k\pm1$; the local quotient solutions require $\pm3$.  Thus $q=441$ is a
projection-image gluing obstruction before the full T3 PAF join.  At $q=549$,
both prime quotients lift, but the prime-power marginal $t'=25$ is infeasible.
These examples show that composite marginals and projection gluing are
strictly stronger than all prime-quotient tests taken separately.

\subsection{Plan toward the full converse}\label{sec:full-converse-plan}

The next theorem must replace the finite T2/T3 harvest by a structural
statement about orbit algebras.  Orbit shape gives \emph{compression}; it does
not by itself give an obstruction.  The full converse program therefore has two
layers:

\[
\begin{array}{ll}
\text{compression families} &
  \mathcal A(t',V')=\Z[C_{t'}]^{V'}\text{ has a small or rigid orbit model},\\[2mm]
\text{firing criteria} &
  \text{arithmetic, congruence, energy, box/parity, or gluing arguments}\\
  &\text{prove that the compressed finite system has no global witness.}
\end{array}
\]

The marginal algebra lemma, the transitive two-orbit theorem
(Theorem~\ref{thm:t2-twoorbit}), and the prime-square lift certificate
(Theorem~\ref{thm:prime-square-lift}) are now compression frameworks. They
become obstructions only when their associated finite joins are shown
infeasible. This caution is essential: the two-orbit marginal at
$q=1325,t'=3$ is solvable, and prime-square lift marginals through $q\le2000$
split between SAT and UNSAT examples.

The remaining work is therefore the following.

\begin{enumerate}
\item Keep the composite multiplier foundation explicit.
Theorem~\ref{thm:composite-multiplier} and
Lemma~\ref{lem:composite-coherence} are the source of all
$M(q)$-invariance used by T3 and by the marginal system.

\item Unify the finite component criteria rather than multiplying special
cases.  Proposition~\ref{prop:component-join} is the general template: the
fixed algebra decomposes into conductor components, each component gives a
two-square equation over a fixed subfield, the embedding bound
$|\sigma(\lambda)|\le\sqrt q$ makes the component search finite, and the
inverse transform imposes the integral boxes and parities.  The $C_{27}$
Gaussian congruence criteria
(Propositions~\ref{prop:c27-component}--\ref{prop:c27-nested}), the $p=3$
Gaussian criterion (Proposition~\ref{prop:p3-gaussian}), the $p=5$ quadratic
criterion (Proposition~\ref{prop:p5-sqrt5}), the $C_{21}$ quadratic criterion
(Proposition~\ref{prop:c21-sqrt21}), the $q=1325,t'=51$ quartic criterion
(Propositions~\ref{prop:c51-quartic}--\ref{prop:q1325-quartic}), and the
prime sub-torus period criterion
(Proposition~\ref{prop:prime-subtorus-periods} with
Corollary~\ref{cor:prime-subtorus-firing}) are instances of this framework.
The $h=3$ support-augmentation obstruction
(Proposition~\ref{prop:h3-support-augmentation}) is the complementary
small-fiber criterion.

\item The selection layer of the existential divisor-selection problem is
now a theorem.  The group
\[
  M(q)=\bigcap_{p\mid q}\langle p\rangle\subset(\Z/4t)^\times
\]
controls every $V'$-orbit structure in the precise sense of
Lemma~\ref{lem:orbit-census}, and
Theorem~\ref{thm:divisor-selection} shows that the existential quantifier
over all divisors $t'\mid t$ is complete and monotone: nonexistence at $q$
is \emph{equivalent} to emptiness of some marginal node, the empty nodes
form an up-set, and criterion-equipped nodes always exist.  This absorbs
the earlier obligation to ``quantify over all divisors rather than a
shape-ranked divisor.''  What it deliberately does not provide
(Remark~\ref{rem:selection-open}) is a proof that a written criterion
\emph{fires} at some node for every composite $q$ --- by
Theorem~\ref{thm:divisor-selection}(i) that statement is the all-plus
converse itself.  The remaining structural work is therefore the
per-family firing theorems: arithmetic hypotheses on $(q,M(q))$ under
which the period join, the $h=3$ hierarchy, the component joins, or the
top-node sign join is provably empty.  The first of these is now in
place: the two-orbit family is completely characterized by
Theorem~\ref{thm:two-orbit-reps}, with the zero-computation parity case
(Corollary~\ref{cor:two-orbit-parity}), the transitive full torus
(Corollary~\ref{cor:transitive-full-torus}), and the never-firing $C_3$
node (Corollary~\ref{cor:c3-feasible}); the prime $h=3$ support stage
likewise has the closed form of
Corollary~\ref{cor:h3-support-sum-prime}, and the prime-$t$ top node has
the row-sum bound (Theorem~\ref{thm:full-torus-rowsum}), which closes the
whole full-torus branch whenever $(w-1)^2\ge t$ and bounds the multiplier
image of any genuine prime-order pair by $w<1+\sqrt t$.  The $h=3$
hierarchy is now closed-form through its augmentation stage at prime
nodes (Proposition~\ref{prop:h3-augmentation-reps},
Corollary~\ref{cor:h3-augmentation-residue}); this is the routing theorem
for the dominant local-pass class, because
Corollary~\ref{cor:c3-feasible} shows the $C_3$ node never fires while
Theorem~\ref{thm:divisor-selection}(iii) guarantees the $h=3$ node
$t'=t/3$ for exactly those orders.  The small-image prime-$t$ top nodes
($w\le d+1$) now have a stated mechanism with two decided instances:
the fold-amplification dichotomy
(Proposition~\ref{prop:fold-dichotomy}, with
Lemma~\ref{lem:stabilizer-rigidity}) closes $q=5185$ and $q=62305$,
each GRH-conditional through exactly one class-group input, cascading
into the exact trivial-character predicate; the census to $q\le100000$
leaves four survivors of which these are two, and the family's
remaining obligations are the Stickelberger lemma for the firing
condition, unconditionality of the two instances, and a refinement for
the two wall witnesses (Remark~\ref{rem:fold-gaps}).  The secondary branch is
now empty: the full-unit prime-power criterion
(Proposition~\ref{prop:pe-full-unit-nested}), the $C_{57}$ criterion
(Proposition~\ref{prop:c57-sqrt57}), and the measured box gate for
degree-$4$/$5$ period joins promote its last four orders.  The families
still needing firing characterizations are the very wide period joins
(the lattice-reduced enumeration now evaluates every node up to
$\sim2\cdot10^7$ true lattice points; the surviving refusals $8865$,
$8905$, $9265$ sit at $10^9$--$10^{42}$ points and need an argument,
not enumeration), the $h=3$ PAF stage at wide nodes, the composite
component joins beyond the instances above, and composite-$t$ top
nodes.

\item The remaining finite-certificate branch has been replaced in the finite
panel.  The case $q=1469$ is now closed at the $p=7$ prime-square marginal
$t'=49$ by the cubic component criterion
(Proposition~\ref{prop:q1469-p7-cubic}), so the current $q\le2000$
accounting has no HiGHS/MILP-only branch.  The remaining obligation is to
unify these component criteria and the $h=3$ hierarchy into broader firing
theorems rather than isolated low-degree instances.

\item Keep the signed branch as a theorem-facing obstruction branch, not an
implementation guard.  Lemma~\ref{lem:signed-branch-reduction} reduces any
post-two-squares signed case to
$q=n^2$ with every prime divisor of $n$ congruent to $3\bmod4$.
Proposition~\ref{prop:signed-marginal} and
Corollary~\ref{cor:signed-full-torus-sign} give the signed finite system, and
Corollary~\ref{cor:signed-full-unit-prime} closes the full-unit prime-torus
subcase analytically.  The signed audit through $q\le50000$ finds three
multi-prime classical survivors in this square branch:
$29241$, $40401$, and $43681$; all three are closed by the signed full-torus
criteria.  Thus $\epsilon=-1$ is no longer a separate implementation-side
blocker.  The structural divisor theorem must still route any future signed
square branch into a signed firing criterion, just as it must route all-plus
branches into their firing criteria.

\item The $q=441$ projection-image gluing obstruction is generalized:
Corollary~\ref{cor:projection-image} states it for an arbitrary node
against any family of exactly solved quotient nodes, as a consequence of
lattice functoriality.  In lattice terms the remaining gluing question is
sharp (Remark~\ref{rem:selection-open}): is the family of
Theorem~\ref{thm:divisor-selection}(iv) --- composite $t$ with every
proper node feasible --- empty?  No member with $q\le2000$ is known; a
member would be a $q$ whose converse is decided only at the exact
full-torus node.

\item Use wider sweeps only as evidence.  Under the current branch split, the
same scaffold was refreshed through $q\le10000$ with dynamic secondary exact
checks, dynamic full-torus sign joins, and dynamic $h=3$ support joins
disabled.  With the lattice-reduced period enumeration in place (an
LLL-reduced embedding-cube slicer with exact sup-norm filtering,
gated by the true point-count predictor
$(2\sqrt q)^d/\sqrt{\ell^{d-1}}$ rather than the coordinate-box size),
it reports $415$ proper-divisor quantified fires, $135$ full-torus
fires ($82$ by the zero-computation row-sum bound of
Theorem~\ref{thm:full-torus-rowsum}), $1$ remaining full-torus sign
lane --- the small-image case $q=5185$, decided GRH-conditionally by
Proposition~\ref{prop:q5185-closure} but kept in the field-free
accounting pending certification --- $56$ local-pass quantified lanes,
$3$ recorded secondary exact marginal certificates, no recorded finite
certificates, $1$ projection-gluing obstruction, and $3$ proper
component-only lanes: $8865$, $8905$, $9265$, refused honestly by the
predictor at $\approx2.0\cdot10^9$, $2.4\cdot10^{42}$, and
$1.2\cdot10^{11}$ lattice points.  The formerly refused nodes
$7045@271$ ($1{,}815{,}037$ points) and $4705@181$ ($15{,}123{,}779$
points) are now exact \texttt{no\_side\_vector} fires, each
cross-checked against the coordinate-box enumerator where both run.
The run is deliberately evidence-only: the disabled dynamic branches
mark where the structural theorem must look, not a substitute for that
theorem.
\end{enumerate}


\begin{thebibliography}{99}

\bibitem{Turyn72} R.~J. Turyn, An infinite class of Williamson matrices,
\emph{J. Combin. Theory Ser.~A} \textbf{12} (1972), 319--321.

\bibitem{Turyn65} R.~J. Turyn, Character sums and difference sets,
\emph{Pacific J. Math.} \textbf{15} (1965), 319--346.

\bibitem{Turyn68} R.~J. Turyn, Sequences with small correlation, in:
\emph{Error Correcting Codes} (H.~B. Mann, ed.), Wiley, New York, 1968.

\bibitem{Williamson44} J.~Williamson, Hadamard's determinant theorem and the
sum of four squares, \emph{Duke Math.~J.} \textbf{11} (1944), 65--81.

\bibitem{Hall67} M.~Hall, Jr., \emph{Combinatorial Theory}, Blaisdell, Waltham,
Mass., 1967.

\bibitem{Whiteman73} A.~L. Whiteman, An infinite family of Hadamard matrices of
Williamson type, \emph{J. Combin. Theory Ser.~A} \textbf{14} (1973), 334--340.

\bibitem{YY85} K.~Yamamoto and M.~Yamada, Williamson Hadamard matrices and Gauss
sums, \emph{J. Math. Soc. Japan} \textbf{37} (1985), 703--717.

\bibitem{LS12} W.~Lang and E.~Schneider, Turyn type Williamson matrices up to
order 99, \emph{Des. Codes Cryptogr.} \textbf{62} (2012), 79--84.

\bibitem{KMT26} H.~Kharaghani, A.~Mohammadian, and B.~Tayfeh-Rezaie, A search
for Hadamard matrices of Williamson type, arXiv:2605.08661 (2026).

\bibitem{GS67} J.~M. Goethals and J.~J. Seidel, Orthogonal matrices with zero
diagonal, \emph{Canad. J. Math.} \textbf{19} (1967), 1001--1010.

\bibitem{DGS71} P.~Delsarte, J.~M. Goethals, and J.~J. Seidel, Orthogonal
matrices with zero diagonal.~II, \emph{Canad. J. Math.} \textbf{23} (1971),
816--832.

\bibitem{Seberry16} J.~Seberry, Remarks on negacyclic matrices, preprint, 2016.

\bibitem{Wolfmann92} J.~Wolfmann, Almost perfect autocorrelation sequences,
\emph{IEEE Trans. Inform. Theory} \textbf{38} (1992), 1412--1418.

\bibitem{PB95} A.~Pott and S.~P. Bradley, Existence and nonexistence of
almost-perfect autocorrelation sequences, \emph{IEEE Trans. Inform. Theory}
\textbf{41} (1995), 301--304.

\bibitem{BAOD22} S.~Barrera Acevedo, P.~\'O~Cath\'ain, and H.~Dietrich, Cocyclic
two-circulant core Hadamard matrices, \emph{J. Algebraic Combin.} \textbf{55}
(2022), 201--215.

\bibitem{Bose42} R.~C. Bose, An affine analogue of Singer's theorem,
\emph{J. Indian Math. Soc.} \textbf{6} (1942), 1--15.

\bibitem{Jungnickel82} D.~Jungnickel, On automorphism groups of divisible
designs, \emph{Canad. J. Math.} \textbf{34} (1982), 257--297.

\bibitem{JK} D.~Jungnickel and H.~Kharaghani, Balanced generalized weighing
matrices and their applications, \emph{Matematiche (Catania)} \textbf{59}
(2004), 225--261.

\bibitem{Luke97} H.~D. L\"uke, Almost-perfect polyphase sequences with small
phase alphabet, \emph{IEEE Trans. Inform. Theory} \textbf{43} (1997), 361--362.

\bibitem{Luke01} H.~D. L\"uke, Almost-perfect quadriphase sequences,
\emph{IEEE Trans. Inform. Theory} \textbf{47} (2001), 2607--2608.

\bibitem{ZHL05} X.~Zeng, L.~Hu, and Q.~Liu, A novel method for constructing
almost perfect polyphase sequences, in: \emph{Coding and Cryptography --- WCC
2005}, LNCS \textbf{3969}, Springer, 2006, pp.~346--353.

\bibitem{AF19} J.~A. Armario and D.~L. Flannery, Almost supplementary difference
sets and quaternary sequences with optimal autocorrelation, arXiv:1911.08828
(2019).

\bibitem{BD15} N.~A. Balonin and D.~\v{Z}. {\DJ}okovi\'c, Negaperiodic Golay
pairs and Hadamard matrices, arXiv:1508.00640 (2015).

\bibitem{JP24} J.~Jedwab and T.~Pender, Two constructions of quaternary Legendre
pairs of even length, arXiv:2408.08472 (2024).

\bibitem{KKW24} I.~S. Kotsireas, C.~Koutschan, and A.~Winterhof, Quaternary
Legendre pairs II, arXiv:2408.16318 (2024).

\bibitem{Schmidt99} B.~Schmidt, Cyclotomic integers and finite geometry,
\emph{J. Amer. Math. Soc.} \textbf{12} (1999), 929--952.

\bibitem{LS16} K.~H. Leung and B.~Schmidt, The anti-field-descent method,
\emph{J. Combin. Theory Ser.~A} \textbf{139} (2016), 87--131.

\bibitem{LS19} K.~H. Leung and B.~Schmidt, Nonexistence results on generalized
bent functions, \emph{J. Combin. Theory Ser.~A} \textbf{163} (2019), 1--33.

\bibitem{HR74} H.~Halberstam and H.-E. Richert, \emph{Sieve Methods},
Academic Press, London, 1974.

\bibitem{FI10} J.~Friedlander and H.~Iwaniec, \emph{Opera de Cribro},
Amer. Math. Soc. Colloq. Publ. \textbf{57}, Providence, 2010.

\bibitem{LO77} J.~C. Lagarias and A.~M. Odlyzko, Effective versions of
the Chebotarev density theorem, in: \emph{Algebraic Number Fields}
(A.~Fr\"ohlich, ed.), Academic Press, London, 1977, 409--464.

\bibitem{Stephens69} P.~J. Stephens, An average result for Artin's
conjecture, \emph{Mathematika} \textbf{16} (1969), 178--188.

\bibitem{HBP79} D.~R. Heath-Brown and S.~J. Patterson, The distribution
of Kummer sums at prime arguments, \emph{J. Reine Angew. Math.}
\textbf{310} (1979), 111--130.

\bibitem{Chen73} J.~R. Chen, On the representation of a large even
integer as the sum of a prime and the product of at most two primes,
\emph{Sci. Sinica} \textbf{16} (1973), 157--176.

\bibitem{Dickson04} L.~E. Dickson, A new extension of Dirichlet's
theorem on prime numbers, \emph{Messenger of Math.} \textbf{33} (1904),
155--161.

\bibitem{Irving14} A.~J. Irving, Almost-prime values of polynomials at
prime arguments, arXiv:1410.3333 (2014).

\bibitem{Lewulis16} P.~Lewulis, Chen primes in arithmetic progressions,
arXiv:1601.02873 (2016).

\end{thebibliography}

\appendix

\section{Character liftings and a fillability dichotomy}
\label{sec:mary}

The same pair object arises from the construction side, uniformly in the order
$m$ of the developing character. This section is self-contained and independent
of \S\S\ref{sec:hierarchy}--\ref{sec:defect}; its \emph{existence} content is
classical (the balanced generalized weighing matrix
$\mathrm{BGW}(q+1,q,q-1)$ from Bose's affine difference
set~\cite{Bose42,Jungnickel82,JK}), and what is new is the explicit two-coset
pair normal form, the fillability dichotomy, and the arithmetic specializations.

\begin{remark}[Notation for this section]\label{rem:maryscope}
Throughout Section~\ref{sec:mary}, $m$ denotes the order of a multiplicative
character (a divisor of $q-1$), a role \emph{distinct} from the prime $m\mid t$
of \S\S\ref{sec:rigidity}--\ref{sec:defect}; we write $E=\F_{q^2}$ for the
quadratic extension (called $K$ in the source note), reserving $K$ for the real
cyclotomic field above; and $B(x,y)$ denotes an alternating form (not a
Williamson block).
\end{remark}

\subsection{Notation}\label{sec:marynotation}

Let $q\equiv1\pmod4$ be an odd prime power, $t=(q+1)/2$ odd, $\F=\F_q$,
$E=\F_{q^2}$, with $x\mapsto x^q$ generating $\Gal(E/\F)$. Fix a primitive
$\omega\in E^\times$; $g=\omega^{q+1}$ (a primitive element of $\F^\times$);
$\zeta=\omega^{q-1}$ of order $q+1=2t$ with $\zeta^t=-1$; $\theta=\omega^t$
(so $\theta^q=-\theta$, spanning the trace-zero $\F$-line of $E$); the
nondegenerate alternating $\F$-form $B(x,y)=(xy^q-x^qy)/\theta$; a character
$\chi$ of $\F^\times$ of exact order $m\mid q-1$ (extended by $\chi(0)=0$); and
$\xi=\chi(g)$, a primitive $m$-th root of unity, so $\chi(g^m)=1$ and
$\chi(-1)=\xi^{(q-1)/2}\in\{\pm1\}$. Following Turyn's ordering with $4$ replaced
by $2m$, set the integer exponents
\begin{equation}\label{eq:ordering}
  e_j=2mj\ (0\le j\le t-1),\qquad e_j=t+2m(j-t)\ (t\le j\le2t-1),
\end{equation}
points $x_j=\omega^{e_j}$, and the $(q+1)\times(q+1)$ matrix
$M_{jk}=\chi(B(x_j,x_k))$, written in $t\times t$ blocks with $R=M_{11}$,
$S=M_{12}$, first rows $c_k=R_{0k}$, $s_k=S_{0k}$. For a length-$t$ sequence $u$,
$\paf_u(k)=\sum_j u_j\overline{u_{j+k}}$. Two identities recur: the scaling
rules $B(\lambda x,y)=B(x,\lambda y)=\lambda B(x,y)$,
$B(\omega^kx,\omega^ky)=g^kB(x,y)$, $B(y,x)=-B(x,y)$; and the entry formula
$B(\omega^a,\omega^b)=\omega^{a+b}(\zeta^b-\zeta^a)/\theta$.

\subsection{The normal form}

\begin{lemma}[Coverage]\label{lem:coverage}
$(x_j)_{0\le j\le2t-1}$ represents $\PG(1,q)=E^\times/\F^\times$.
\end{lemma}

\begin{proof}
$\omega^a\equiv\omega^b\bmod\F^\times$ iff $(q+1)\mid a-b$. Any common divisor of
$q-1$ and $t$ divides $2t-(q-1)=2$; as $t$ is odd, $\gcd(q-1,t)=1$, hence
$\gcd(m,t)=1$. So $2mj\bmod2t$ ($0\le j\le t-1$) hits each even residue once and
$t+2m(j-t)$ hits each odd residue once.
\end{proof}

\begin{lemma}[Orthogonality]\label{lem:orth}
$MM^*=qI_{q+1}$, for any representatives.
\end{lemma}

\begin{proof}
$B(x_j,x_k)=0$ iff $k=j$, so each row has $q$ unimodular entries and
$(MM^*)_{jj}=q$. For $j\ne k$,
$(MM^*)_{jk}=\sum_{l\ne j,k}\chi(B(x_j,x_l)/B(x_k,x_l))$; the map
$y\mapsto(B(x_j,y),B(x_k,y))$ is an $\F$-linear bijection onto $\F^2$, so the
ratio takes each value of $\F^\times$ once as $[y]$ ranges over the other
points, and $\sum_{u\in\F^\times}\chi(u)=0$.
\end{proof}

\begin{theorem}[Normal form]\label{thm:normalform}
All four blocks of $M$ are circulant, and
\[
  M=\begin{bmatrix} R & S\\ \chi(-1)\,S & \chi(-g)\,R\end{bmatrix},
\]
where: \emph{(1)} $R$ has zero diagonal, $c_k\ne0$ for $k\ne0$, and
$c_{t-k}=\chi(-1)c_k$; \emph{(2)} $S$ is symmetric ($s_{t-k}=s_k$) with all
entries nonzero; \emph{(3)} all nonzero entries are $m$-th roots of unity and
$M^T=\chi(-1)M$; \emph{(4)} $RR^*+SS^*=qI_t$, equivalently
$\paf_r(k)+\paf_s(k)=0$ for $k\not\equiv0\pmod t$.
\end{theorem}

\begin{proof}
\emph{Circulant blocks.} Incrementing both indices within a block pair raises
both exponents in~\eqref{eq:ordering} by $2m$; by the scaling rules the argument
of $\chi$ multiplies by $g^{2m}$ and $\chi(g^{2m})=\xi^{2m}=1$. A wrap replaces
$e+2m$ by $e+2m-2mt$; since $\omega^{2mt}=g^m\in\F^\times$ and $\chi(g^{\mp m})=1$,
every entry depends only on $(k-j)\bmod t$ within its block.
\emph{$M_{22}=\chi(-g)R$.} By~\eqref{eq:ordering}, $\zeta^t=-1$, $\omega^{2t}=g$:
$B(\omega^t,\omega^{t+2mk})=g\,\omega^{2mk}(1-\zeta^{2mk})/\theta=-g\,B(1,\omega^{2mk})$;
apply $\chi$. \emph{$S$ symmetric.} With $\omega^{2mt}=g^m$, $\chi(g^m)=1$:
$B(1,\omega^{t-2mk})/B(1,\omega^{t+2mk})=g^{-2mk}$ (using
$\zeta^{t\pm2mk}=-\zeta^{\pm2mk}$), and $\chi(g^{-2mk})=1$; the common factor
$1+\zeta^{2mk}\ne0$. \emph{$M_{21}=\chi(-1)S$, $M^T=\chi(-1)M$.} $B$ alternating
gives $M_{kj}=\chi(-1)M_{jk}$; with $S=S^T$ this yields the claim. \emph{Diagonal
and zeros.} $c_0=\chi(B(1,1))=0$; $c_k=0$ iff $t\mid k$ (as $\gcd(m,t)=1$);
$s_k\ne0$ by parity; and $c_{t-k}=\chi(-1)c_k$ from circulancy and the
alternating rule. \emph{Pair identity.} The top-left block of $MM^*=qI$
(Lemma~\ref{lem:orth}) is $RR^*+SS^*$.
\end{proof}

\begin{corollary}[Turyn 1972, the case $m=2$]\label{cor:turyn2}
For $m=2$: $\chi(-1)=1$, $\chi(-g)=-1$, so $M=\bigl[\begin{smallmatrix}R&S\\
S&-R\end{smallmatrix}\bigr]$ is the symmetric conference matrix with $R,S$
symmetric circulants and $R^2+S^2=qI$; $M+iI$ is a Butson--Hadamard
$\mathrm{BH}(q+1,4)$, and $A=I+R$, $B=I-R$, $C=D=S$ is a Williamson quadruple of
order $4t=2(q+1)$.
\end{corollary}

\subsection{The obstruction}

\begin{theorem}[Fillability dichotomy]\label{thm:dichotomy}
A unimodular $\delta$ with $(M+\delta I)(M+\delta I)^*=(q+1)I$ exists iff $m\le2$.
\end{theorem}

\begin{proof}
By Lemma~\ref{lem:orth} the condition is $\bar\delta M+\delta M^*=0$; since
$M^*=\chi(-1)\overline M$, entrywise off the diagonal
$M_{jk}^2=-\delta\chi(-1)/\bar\delta$ is constant, i.e.\ all off-diagonal
entries lie in $\{\pm c\}$. But Theorem~\ref{thm:normalform} exhibits $c_k$ (in
$M_{11}$) and $\chi(-g)c_k=\pm\xi c_k$ (in $M_{22}$); both in $\{\pm c\}$ forces
$\xi\in\{\pm1\}$, i.e.\ $m\le2$. Conversely $m=2$ admits $\delta=\pm i$
(Corollary~\ref{cor:turyn2}).
\end{proof}

Theorem~\ref{thm:dichotomy} explains why Turyn's Hadamard--Williamson endpoint
was never extended past the quadratic character: for $m\ge3$ the invariant
content is the circulant pair $(R,S)$ itself, an \emph{almost-perfect
complementary $m$-ary pair} of odd length $t$ with a single zero entry, existing
for every prime power $2t-1\equiv1\pmod4$ and every $m\mid2t-2$. The closest
relatives are the quaternary Legendre pairs of~\cite{JP24,KKW24}, which require
fully unimodular entries and $\paf$-sum $\equiv-2$ and live at length $(q-1)/2$
(the affine side) rather than $(q+1)/2$ (the projective side).

\begin{remark}[Quaternionic completion, and the exact boundary]\label{rem:quat}
The pair can be re-assembled: $N=\bigl[\begin{smallmatrix}R&S\\
-S^*&R^*\end{smallmatrix}\bigr]$ satisfies $NN^*=qI_{q+1}$ for every $m$ (the
circulant blocks commute), has zero diagonal, and is unimodular off it. A
unimodular diagonal $D=\operatorname{diag}(d_j)$ completes it,
$(N+D)(N+D)^*=(q+1)I$, iff $N_{jk}N_{kj}=-d_jd_k$ off the diagonal. In the cross
blocks the left side is $-|S_{jk}|^2=-1$, deleting every constraint on $S$; on a
diagonal block it is $\chi(-1)c_{k-j}^2$, forcing $d$ constant there and the
condition $c_k^2$ constant. If $\chi^2$ is nontrivial, the Hadamard square
$M\circ M$ is the character matrix of $\chi^2$, so
$(R\circ R)(R\circ R)^*+(S\circ S)(S\circ S)^*=qI_t$; with $R\circ R=c^2(J-I)$
this forces $(S\circ S)(S\circ S)^*=(q-1)I-(t-2)J$, and evaluating on the
all-ones vector gives $|\sum_k s_k^2|^2=-t^2+4t-2<0$ for $t\ge5$. Hence the
quaternionic completion exists iff $m\le2$ or $t=3$; the only new case is
$(q,m)=(5,4)$, a $\mathrm{BH}(6,4)$ already reached at $q=5$. The obstruction is
intrinsic to the phase spread of $c_k^2$, not to the $2\times2$ arrangement.
\end{remark}

\subsection{Arithmetic specializations at the trivial character}

Evaluating $RR^*+SS^*=qI$ on the all-ones vector gives, with $\rho=\sum_k c_k$
and $\sigma=\sum_k s_k$ in $\Z[\zeta_m]$,
\begin{equation}\label{eq:twonorm}
  |\rho|^2+|\sigma|^2=q,
\end{equation}
a decomposition of $q$ into two $\Z[\zeta_m]$-norms produced by character sums.

\begin{proposition}\label{prop:twosquaresm4}
For $m=4$ and $q\equiv5\pmod8$: $\rho=0$ and $\sigma\in\Z[i]$ has $|\sigma|^2=q$;
the construction witnesses $q=a^2+b^2$.
\end{proposition}

\begin{proof}
$\chi(-1)=\xi^{(q-1)/2}=\xi^2=-1$ (as $(q-1)/2\equiv2\bmod4$, $\xi=\pm i$). By
Theorem~\ref{thm:normalform}(1), $c_{t-k}=-c_k$, and the pairs $(k,t-k)$
partition the nonzero indices ($t$ odd), so $\rho=0$; then~\eqref{eq:twonorm}
gives $|\sigma|^2=q$ with $\sigma\in\Z[i]$.
\end{proof}

For $q\equiv1\pmod8$ one has $\chi(-1)=1$, so $\rho\in2\Z[i]$; numerically more
holds, and the basic $2$-adic divisibility is elementary.

\begin{proposition}[The quartic $2$-adic refinement]\label{prop:Lstar}
For $m=4$ and $q\equiv1\pmod8$, $\rho\equiv0\pmod{(1+i)^3}$ in $\Z[i]$. Hence
$8$ divides $|\rho|^2$, and~\eqref{eq:twonorm} decomposes $q$ as a sum of two
Gaussian norms with the first norm divisible by $8$.
\end{proposition}

\begin{proof}
Here $\chi(-1)=1$, so Theorem~\ref{thm:normalform}(1) gives
$c_{t-k}=c_k$. Since $c_0=0$ and $t$ is odd,
\[
  \rho=2\sum_{k=1}^{(t-1)/2}c_k .
\]
Each $c_k$ in the sum is a Gaussian unit. In the quotient $\Z[i]/(1+i)$ all four
Gaussian units are congruent to $1$, while
$(t-1)/2=(q-1)/4$ is even because $q\equiv1\pmod8$. Therefore the half-row sum is
divisible by $(1+i)$. Since $2=-i(1+i)^2$, the displayed formula for $\rho$
gives $\rho\in(1+i)^3\Z[i]$.
\end{proof}

\begin{remark}
In the computed prime cases $q\le457$, $q\equiv1\pmod8$, the stronger diagonal
pattern $\rho=a(1\pm i)$ with $a$ even also holds (and $\sigma$ lies on a
coordinate axis), recovering the classical-looking form $q=8x^2+b^2$ inside
this normalization. That extra phase statement is not used here.
\end{remark}

Only $|\rho|^2$, $|\sigma|^2$, and $\rho=0$ are invariant under the global
scalar $\chi(c)$; statements about the phase of $\sigma$ alone are not, and are
avoided.

\subsection{Dictionary and related objects}\label{sec:marydict}

\begin{proposition}[Twisted-circulant form]\label{prop:dictionary}
Set $w_e=\chi(B(1,\omega^e))$, so $w_{e+2t}=\xi w_e$. Let $C$ be the
$2t\times2t$ $\xi^{-1}$-circulant $C_{jk}=w_{(k-j)\bmod2t}\,\xi^{-[k<j]}$. Then
$M$ is monomially equivalent to $C$; in particular $C$ has zero diagonal,
$m$-th-root entries off it, and $CC^*=qI$. For $m=2$, $C$ is the negacyclic
conference matrix of~\cite{GS67,DGS71}.
\end{proposition}

Before applying $\chi$, the matrix $[B(x_i,x_j)]$ organized by the Singer cycle
encodes Bose's affine difference set~\cite{Bose42}: a $(q+1,q-1,q,1)$ relative
difference set in $E^\times\cong\Z_{q^2-1}$ relative to $\F^\times$. By
Jungnickel's theorem~\cite{Jungnickel82} (quoted as~\cite[Thm 3.6]{JK}), cyclic
relative difference sets are equivalent to $\omega$-circulant balanced
generalized weighing matrices; the family $\mathrm{BGW}(q+1,q,q-1)$ relevant here
is the $d=1$ case, appearing in~\cite{JK} as the generalized-conference example
$\mathrm{BGW}(5,4,3)$. Our $C$ (hence $M$) is the character image of such a
representative; Proposition~\ref{prop:dictionary} makes the untwisting into the
two-coset pair form explicit, extending Turyn's remark that his construction
``is equivalent to the one in [4]''~\cite{Turyn72} from $m=2$ to all $m$.

\begin{remark}[Single-sequence unrolling]\label{rem:unroll}
Extending $u_e=\chi(B(1,\omega^e))$ over $e\in\Z$ gives $u_{e+2t}=\xi u_e$, a
period-$m(q+1)$ sequence with $m$ zeros per period; its full-period
autocorrelation is $mq\,\xi^{-r}$ at shifts $2tr$ and $0$ otherwise---an
almost-perfect $m$-phase sequence. Almost-perfect polyphase sequences with small
phase alphabet go back to L\"uke~\cite{Luke97,Luke01}, and Zeng--Hu--Liu
construct almost-perfect quadriphase families of exactly these lengths
$m(p^J+1)$~\cite{ZHL05}; we verify only that the parameters coincide. For $m=2$
the pair $(r,s)$ is the input to the two-circulant-core construction~\cite{BAOD22};
the underlying almost-perfect autocorrelation sequences were introduced by
Wolfmann~\cite{Wolfmann92} and classified via cyclic divisible difference sets
by Bradley--Pott~\cite{PB95}. Turyn himself derived, on p.\ 321
of~\cite{Turyn72}, a perfect-autocorrelation quaternary circulant from the pair;
our $(R,S)$ for $m\ge3$ is the character-theoretic extension of that page.
\end{remark}

Adjacent modern objects: the odd-length single quaternary sequences of
Armario--Flannery~\cite{AF19} are fully unimodular with $\paf=\pm1$ (no quartic
characters, no pair); negaperiodic Golay pairs~\cite{BD15} are binary of even
length; the quaternary Legendre pairs~\cite{JP24,KKW24} live on the affine side.
Every clause of Theorems~\ref{thm:normalform}--\ref{thm:dichotomy},
Corollary~\ref{cor:turyn2}, and Propositions~\ref{prop:twosquaresm4}
and~\ref{prop:Lstar} have been
verified entry by entry, in exact arithmetic, for $59$ pairs $(q,m)$ with $q$
prime up to $113$ and $m\mid q-1$, $m\le12$ (zero failures), and
Remark~\ref{rem:quat} likewise.

\begin{remark}[Further directions on the construction side]\label{rem:maryfurther}
Two questions are left open. \emph{(1)} For $q\equiv3\pmod4$ ($t$ even) the
two-coset split in~\eqref{eq:ordering} breaks; an offset-$1$ ordering with $m$
odd appears to be the correct skew analogue. \emph{(2)} The fillability
dichotomy suggests a general question: among character images of
$\omega$-circulant $\mathrm{BGW}$ matrices, which admit \emph{monomial}
completions to Butson--Hadamard matrices? Theorem~\ref{thm:dichotomy} settles
diagonal fillings; row/column rescalings before filling are not addressed here.
\end{remark}


\section{Verification status and provenance}

\textbf{Tiers.} Tier~0: unconditional, complete proof in this document.
Tier~1: conditional on GRH. Tier~2: conditional on a named explicit
hypothesis (EQ$_3$, hypothesis (C), Dickson). Tier~3: numerical evidence
only, no theorem claimed.

\textbf{Principal claims.} Selection theorem, T0. Two-orbit criterion,
T0. $h=3$ hierarchy, T0. Row-sum bound and exact form, T0. Rigidity
criterion, T0. $q\le2000$ closure and first determinations $1469, 1937,
1325$, T0 with a closed-table novelty caveat. $q=5185$, $q=62305$, T1.
Semiprime lane theorem, T0. Lane collapse at general $\omega$
(Proposition~\ref{prop:lane-kill}), T0. Dickson family and the
hypothesis-(C) dichotomy, T2. Small-factor sieve count
(Proposition~\ref{prop:small-factor-sieve}), T2 (EQ$_3$); GRH implies
EQ$_3(\theta)$ for $\theta<1/8$ (T1) but does not close the lane's
remaining structure count, so the lane density program is not
GRH-complete. Census calibrations at $q\le10^4$--$10^7$, T3.

\textbf{Citations (source-verified 2026-07-07).} The bibliography entries
for Halberstam--Richert, Friedlander--Iwaniec, Lagarias--Odlyzko,
Stephens, and Heath-Brown--Patterson were checked against public records
and are bibliographically accurate (e.g.\ Lagarias--Odlyzko pp.~409--464;
Stephens \emph{Mathematika} \textbf{16} (1969) 178--188;
Heath-Brown--Patterson \emph{J. Reine Angew. Math.} \textbf{310} (1979)
111--130). One in-text precision point remains: the chapter pointer for
Halberstam--Richert could not be independently confirmed, so those
citations are now given without a chapter number.

\textbf{Provenance.} This document was written by large language models
(primarily Claude, with some work by OpenAI's Codex GPT-5.5) under human
direction, with machine-verified audits for the
finite claims and independent-session cross-verification for the
theoretical ones; one conditional theorem announced during development
(a conditional infinite family on the lane) was retracted after
verification and does not appear in this document. The working record,
including the retraction, is preserved in the companion notes.


\end{document}
